% BIBKEYS: Allard1972, Federer1969, Reifenberg1960a, Almgren1968
%
% Sections 6--8 of the blueprint of Allard, On the first variation of a varifold (1972).
% New macro (checked unused in preambule.tex, macros.tex and the master):
\providecommand{\affgrass}[2]{\mathbf{A}(#1,#2)}   % Allard's A(n,k), see 8.7

\section{The compactness theorem for integral varifolds}
\label{sec:compactness}

\begin{miniremark}
    \label{mr:sec6}
    This section treats~\cite[section~6]{Allard1972}. Lemma~\ref{al:6.1} splits a~varifold with
    small tilt and small first variation near a~finite set~$Y$ into two pieces, each of which
    concentrates near fewer points of~$Y$; by induction, \ref{al:6.2} bounds the sum of the
    densities at the points of~$Y$ by the weight of a~neighbourhood of~$Y$. From this,
    \ref{al:6.3} shows that integral varifolds converging to a~multiple $d \var{}(\im T)$ of a~plane
    with vanishing first variation have integral~$d$, and \ref{al:6.4} proves the compactness
    theorem for integral varifolds. Tangent cones are discussed in~\ref{al:6.5}. The monotonicity
    inequality~\ref{new:mono} restates~\ref{al:5.1(3)} for points which need not lie in the
    support. Throughout the section $\vdim \ge 1$. Contrary to Allard, we state~\ref{al:6.1}
    and~\ref{al:6.2} for varifolds in an~open set~$W$ containing the closed $R$~neighbourhood
    of~$Y$; Allard's statements are the case $W = \R^{\adim}$, and the general case is needed
    in~\ref{al:6.3} and~\ref{al:6.4}, where Allard reduces to $U = \R^{\adim}$ without proof.
\end{miniremark}

\begin{lemma}[new]
    \label{new:mono}
    Suppose $W$ is an open subset of~$\R^{\adim}$, $V \in \Var{\vdim}(W)$, $\|\delta V\|$ is
    a~Radon measure on~$W$, $y \in W$, $0 < R < \infty$, $\cball{y}{R} \subseteq W$, $0 \le
    \Lambda < \infty$, and $\|\delta V\| \cball{y}{r} \le \Lambda \|V\| \cball{y}{r}$ for $0 < r
    < R$. Then $[r \mapsto r^{-\vdim} \|V\| \cball{y}{r} \exp(\Lambda r)]$ is nondecreasing on
    $\{ r : 0 < r < R \}$, $\density^{\vdim}(\|V\|,y) \in \R$, and
    \begin{displaymath}
        \unitmeasure{\vdim} \density^{\vdim}(\|V\|,y) \le r^{-\vdim} \|V\| \cball{y}{r}
        \exp(\Lambda r) \le R^{-\vdim} \|V\| \cball{y}{R} \exp(\Lambda R) \quad \text{for $0 < r
        \le R$} \,.
    \end{displaymath}
\end{lemma}

\begin{proof}
    Let $\phi(r) = r^{-\vdim} \|V\| \cball{y}{r} \exp(\Lambda r)$ for $0 < r \le R$. Since
    $\cball{y}{R}$ is a~compact subset of~$W$, $R < \dist(y, \R^{\adim} \without W)$. If $y \in
    \spt \|V\|$, then $\phi$ is nondecreasing on $\{ r : 0 < r < R \}$ by~\ref{al:5.1(3)}. Suppose
    $y \notin \spt \|V\|$ and let $d = \dist(y, \spt \|V\|) > 0$. Then $\phi(r) = 0$ for $0 < r <
    d$. If $d < r < s < R$, then $|\delta V(\kappa(V,y,t))| \le \|\delta V\| \cball{y}{t} /
    \|V\| \cball{y}{t} \le \Lambda$ for $r \le t \le s$ by~\ref{al:4.3}\eqref{eq:4.3:2}, because
    $|\kappa(V,y,t)| \le \|V\| \cball{y}{t}^{-1}$; the right side of~\ref{al:5.1(1)} is
    nonnegative, hence $s^{-\vdim} \|V\| \cball{y}{s} \exp(\Lambda (s - r)) \ge r^{-\vdim} \|V\|
    \cball{y}{r}$, that is $\phi(r) \le \phi(s)$. Since $\cball{y}{d} = \tbcap \{ \cball{y}{r} :
    d < r \}$, $\phi(d) = \lim_{r \to d+} \phi(r)$, so $\phi$ is nondecreasing on $\{ r : 0 < r <
    R \}$ in this case as well. A~nondecreasing nonnegative function has a~finite limit at~$0+$,
    and $\exp(\Lambda r) \to 1$ as $r \to 0+$; hence $\density^{\vdim}(\|V\|,y) \in \R$ and
    $\unitmeasure{\vdim} \density^{\vdim}(\|V\|,y) = \lim_{r \to 0+} \phi(r) \le \phi(r)$ for $0
    < r < R$. Finally $\phi(r) \le \phi(s) \le s^{-\vdim} \|V\| \cball{y}{R} \exp(\Lambda s)$ for
    $r < s < R$, and letting $s \to R$ gives $\phi(r) \le \phi(R)$.
\end{proof}

\begin{lemma}[{\cite[6.1]{Allard1972}}]
    \label{al:6.1}
    Suppose:
    \begin{enumerate}[label=(\arabic*)]
        \item \label{al:6.1(1)} $\nu$ is a~positive integer, $1 \le \mu < \infty$, $0 < \xi <
            1$, $1 < M < \infty$, $0 < R < \infty$, $T \in \pgrass{\adim}{\vdim}$, $Y \subseteq
            \R^{\adim}$, $W$ is an open subset of~$\R^{\adim}$ with $\R^{\adim} \cap \{ x :
            \dist(x,Y) \le R \} \subseteq W$, $V \in \Var{\vdim}(W)$, and $\|\delta V\|$ is
            a~Radon measure on~$W$;
        \item \label{al:6.1(2)} $Y$ has at least two and at most $\nu + 1$ elements;
        \item \label{al:6.1(3)} $|y - z| \le \mu |T^\perp(y - z)|$ whenever $y, z \in Y$;
        \item \label{al:6.1(4)} $(M + 1) \diam Y \le R$;
        \item \label{al:6.1(5)} $R \|\delta V\| \cball{y}{r} \le \xi \|V\| \cball{y}{r}$
            whenever $y \in Y$ and $0 < r < R$;
        \item \label{al:6.1(6)} $\textint{\cball{y}{r} \times \pgrass{\adim}{\vdim}}{} \|S -
            T\| \ud V(x,S) \le \xi \|V\| \cball{y}{r}$ whenever $y \in Y$ and $0 < r < R$.
    \end{enumerate}
    Then $\density^{\vdim}(\|V\|,y) \in \R$ for $y \in Y$, and there are $V_1, V_2 \in
    \Var{\vdim}(W)$ such that $\|\delta V_1\|$ and $\|\delta V_2\|$ are Radon measures on~$W$,
    and a~partition of~$Y$ into subsets $Y_0$, $Y_1$, $Y_2$, such that
    \begin{enumerate}[label=(\arabic*),start=7]
        \item \label{al:6.1(7)} $V \ge V_1 + V_2$;
        \item \label{al:6.1(8)} neither $Y_1$ nor~$Y_2$ has more than $\nu$ elements;
        \item \label{al:6.1(9)} $(M \diam Y) \|\delta V_j\| \cball{y}{r} \le 2 M (\nu + 1)
            (\exp \xi) (3 \nu M \mu)^{\vdim+1} \xi \|V_j\| \cball{y}{r}$ whenever $j \in \{1,2\}$,
            $y \in Y_j$ and $0 < r < M \diam Y$;
        \item \label{al:6.1(10)} $\textint{\cball{y}{r} \times \pgrass{\adim}{\vdim}}{} \|S -
            T\| \ud V_j(x,S) \le M (\exp \xi) (3 \nu M \mu)^{\vdim} \xi \|V_j\| \cball{y}{r}$
            whenever $j \in \{1,2\}$, $y \in Y_j$ and $0 < r < M \diam Y$;
        \item \label{al:6.1(11)} $\density^{\vdim}(\|V\|,y) = \density^{\vdim}(\|V_j\|,y)$
            whenever $j \in \{1,2\}$ and $y \in Y_j$;
        \item \label{al:6.1(12)} with $A = \R^{\adim} \cap \{ x : \dist(x,Y) \le R \}$ and $A_j =
            \R^{\adim} \cap \{ x : \dist(x,Y_j) \le M \diam Y \}$ for $j \in \{1,2\}$,
            \begin{multline*}
                [(1 + 1/M)^{\vdim} + (\nu + 1)/M] (\exp \xi) \frac{\|V\|(A)}{\unitmeasure{\vdim}
                R^{\vdim}} \ge \textsum{y \in Y_0}{} \density^{\vdim}(\|V\|,y) \\ +
                \frac{\|V_1\|(A_1)}{\unitmeasure{\vdim} (M \diam Y)^{\vdim}} + \frac{\|V_2\|(A_2)}
                {\unitmeasure{\vdim} (M \diam Y)^{\vdim}} \,.
            \end{multline*}
    \end{enumerate}
\end{lemma}

\begin{proof}
    For $y \in Y$ we have $\cball{y}{R} \subseteq W$ by~\ref{al:6.1(1)}, and~\ref{al:6.1(5)}
    allows us to apply~\ref{new:mono} with $\Lambda = \xi / R$. Hence
    $\density^{\vdim}(\|V\|,y) \in \R$ and, since $\Lambda s \le \xi$ for $s \le R$,
    \begin{equation}
        \tag{13} \label{al:6.1(13)}
        r^{-\vdim} \|V\| \cball{y}{r} \le s^{-\vdim} \|V\| \cball{y}{s} \exp \xi \,, \qquad
        \unitmeasure{\vdim} \density^{\vdim}(\|V\|,y) \le s^{-\vdim} \|V\| \cball{y}{s} \exp \xi
    \end{equation}
    whenever $y \in Y$ and $0 < r < s \le R$.

    Since $Y$ is finite, there are $\tilde y, \tilde z \in Y$ with $\diam Y = |\tilde y - \tilde
    z|$; by~\ref{al:6.1(2)}, $\diam Y > 0$. Let $f(x) = |T^\perp(x - \tilde y)|$ for $x \in
    \R^{\adim}$, $\eta = M \diam Y$ and $\sigma = |T^\perp(\tilde z - \tilde y)| / (3\nu) = f(\tilde
    z)/(3\nu)$. By~\ref{al:6.1(3)}, $f(\tilde z) \ge \mu^{-1} \diam Y > 0$, so $\sigma > 0$ and
    \begin{equation}
        \tag{14} \label{al:6.1(14)}
        \frac{\eta}{\sigma} = 3 \nu \frac{M |\tilde y - \tilde z|}{|T^\perp(\tilde y - \tilde
        z)|} \le 3 \nu M \mu \,.
    \end{equation}
    Moreover $\eta < (M+1) \diam Y \le R$ and $\sigma \le f(\tilde z) \le \diam Y < R$
    by~\ref{al:6.1(4)}. The $\nu$ open intervals $\{ s : 3 (i-1) \sigma < s < 3 i \sigma \}$, $i =
    1, \ldots, \nu$, are disjoint and contain neither $f(\tilde y) = 0$ nor $f(\tilde z) = 3 \nu
    \sigma$; since $Y \without \{ \tilde y, \tilde z \}$ has at most $\nu - 1$ elements, one of
    these intervals contains no $f(y)$ with $y \in Y$. Let $\rho = 3(i-1) \sigma$ for such~$i$.
    Then $0 \le \rho$, $\rho + 3 \sigma \le f(\tilde z)$, and each $y \in Y$ satisfies either $f(y)
    \le \rho$ or $f(y) \ge \rho + 3 \sigma$. Let
    \begin{gather*}
        Y_0 = Y \cap \{ y : \|V\| \cball{y}{\eta} > M \density^{\vdim}(\|V\|,y)
        \unitmeasure{\vdim} \eta^{\vdim} \} \,, \\
        Y_1 = (Y \cap \{ y : f(y) \le \rho \}) \without Y_0 \,, \qquad Y_2 = (Y \cap \{ y : f(y)
        \ge \rho + 3 \sigma \}) \without Y_0 \,.
    \end{gather*}
    These sets form a~partition of~$Y$. Since $\tilde y \notin Y_2$ and $\tilde z \notin Y_1$,
    \ref{al:6.1(2)} implies~\ref{al:6.1(8)}.

    For $\tau \in \R$ let $V_{1,\tau} = V \restrict \{ (x,S) : f(x) < \tau \}$ and $V_{2,\tau}
    = V \restrict \{ (x,S) : f(x) > \tau \}$. The function~$f$ is continuous, and it is of
    class~$\infty$ on the open set $\{ x : f(x) > 0 \}$, with $\grad f(x) = f(x)^{-1}
    T^\perp(x - \tilde y)$ there. For such~$x$ and $S \in \pgrass{\adim}{\vdim}$, since $|\grad
    f(x)| = 1$ and $T \circ T^\perp = 0$,
    \begin{displaymath}
        |S(\grad f(x))| = |(S - T)(T^\perp(\grad f(x)))| \le \|S - T\| \,.
    \end{displaymath}
    We apply~\ref{al:4.10(2)} to $f$ with $a = \rho + \sigma$ and $b = \rho + 2\sigma$, and to
    $-f$ with $a = - \rho - 2 \sigma$ and $b = - \rho - \sigma$; the hypotheses hold because
    $\rho + \sigma > 0$. Taking $B = \cball{y}{\eta}$ and using~\ref{al:6.1(5)}
    and~\ref{al:6.1(6)} with $r = \eta$, we infer that $\|\delta V_{1,\tau}\|$ and $\|\delta
    V_{2,\tau}\|$ are Radon measures on~$W$ for $\LM^1$ almost all~$\tau$, and that
    \begin{multline*}
        \textint{\rho+\sigma}{\rho+2\sigma} \|\delta V_{j,\tau}\| \cball{y}{\eta} \ud \LM^1 \tau
        \le \textint{\cball{y}{\eta} \times \pgrass{\adim}{\vdim}}{} \|S - T\| \ud V(x,S) +
        \sigma \|\delta V\| \cball{y}{\eta} \\ \le \xi \|V\| \cball{y}{\eta} + \sigma R^{-1} \xi
        \|V\| \cball{y}{\eta} \le 2 \xi \|V\| \cball{y}{\eta}
    \end{multline*}
    for $j \in \{1,2\}$ and $y \in Y$. Let $Q$ be the set of pairs $(j,y)$ with $j \in \{1,2\}$
    and $y \in Y_j$; it has at most $\nu + 1$ elements. For $(j,y) \in Q$ with $\|V\|
    \cball{y}{\eta} = 0$, the last estimate gives $\|\delta V_{j,\tau}\| \cball{y}{\eta} = 0$
    for $\LM^1$ almost all~$\tau$. Let $Q' = Q \cap \{ (j,y) : \|V\| \cball{y}{\eta} > 0 \}$ and
    $\Gamma(\tau) = \textsum{(j,y) \in Q'}{} \|\delta V_{j,\tau}\| \cball{y}{\eta} / (2 \xi
    \|V\| \cball{y}{\eta})$. Then $\textint{\rho+\sigma}{\rho+2\sigma} \Gamma \ud \LM^1 \le
    \nu + 1$, so the set of $\tau$ with $\rho + \sigma < \tau < \rho + 2 \sigma$ and
    $\Gamma(\tau) \le (\nu+1)/\sigma$ has positive $\LM^1$ measure. We may therefore choose~$t$
    with $\rho + \sigma < t < \rho + 2 \sigma$ such that $\|\delta V_{1,t}\|$ and $\|\delta
    V_{2,t}\|$ are Radon measures on~$W$ and
    \begin{equation}
        \tag{15} \label{al:6.1(15)}
        \|\delta V_{j,t}\| \cball{y}{\eta} \le \frac{\nu + 1}{\sigma} 2 \xi \|V\|
        \cball{y}{\eta} \quad \text{whenever $(j,y) \in Q$} \,.
    \end{equation}
    Let $V_1 = V_{1,t}$ and $V_2 = V_{2,t}$. Since $\{ x : f(x) < t \}$ and $\{ x : f(x) > t \}$
    are disjoint, \ref{al:6.1(7)} holds.

    Suppose $y \in Y_1$. Since $|f(x) - f(y)| \le |T^\perp(x - y)| \le |x - y|$ and $t - f(y) >
    \sigma$, the open ball $O = \oball{y}{t - f(y)}$ contains $\cball{y}{\sigma}$ and $f(x) < t$
    for $x \in O$. Hence $V_1 \restrict (O \times \pgrass{\adim}{\vdim}) = V \restrict (O \times
    \pgrass{\adim}{\vdim})$; consequently $\|V_1\| \restrict O = \|V\| \restrict O$ and, since
    $\delta V_1(g) = \delta V(g)$ for $g \in \VF(W)$ with $\spt g \subseteq O$, also $\|\delta
    V_1\| \restrict O = \|\delta V\| \restrict O$. For $0 < r < \sigma$ we infer
    from~\ref{al:6.1(5)}, \ref{al:6.1(6)} and $\eta \le R$ that
    \begin{gather*}
        \eta \|\delta V_1\| \cball{y}{r} \le R \|\delta V\| \cball{y}{r} \le \xi \|V_1\|
        \cball{y}{r} \,, \\ \textint{\cball{y}{r} \times \pgrass{\adim}{\vdim}}{} \|S - T\| \ud
        V_1(x,S) \le \xi \|V_1\| \cball{y}{r} \,,
    \end{gather*}
    which implies~\ref{al:6.1(9)} and~\ref{al:6.1(10)} for these~$r$, because $M (\exp \xi) (3 \nu
    M \mu)^{\vdim} \ge 1$. Since $y \notin Y_0$, \eqref{al:6.1(13)} with $s = \sigma$
    and~\eqref{al:6.1(14)} give
    \begin{multline*}
        \|V\| \cball{y}{\eta} \le M \density^{\vdim}(\|V\|,y) \unitmeasure{\vdim} \eta^{\vdim}
        \le M (\sigma^{-\vdim} \|V\| \cball{y}{\sigma} \exp \xi) \eta^{\vdim} \\ \le M (\exp \xi) (3
        \nu M \mu)^{\vdim} \|V_1\| \cball{y}{\sigma} \,.
    \end{multline*}
    Combining this with~\eqref{al:6.1(15)} and~\eqref{al:6.1(14)} we obtain, for $\sigma \le r <
    \eta$,
    \begin{gather*}
        \eta \|\delta V_1\| \cball{y}{r} \le \eta \frac{\nu+1}{\sigma} 2 \xi \|V\|
        \cball{y}{\eta} \le 2 M (\nu + 1) (\exp \xi) (3 \nu M \mu)^{\vdim + 1} \xi \|V_1\|
        \cball{y}{r} \,, \\
        \textint{\cball{y}{r} \times \pgrass{\adim}{\vdim}}{} \|S - T\| \ud V_1(x,S) \le \xi
        \|V\| \cball{y}{\eta} \le M (\exp \xi) (3 \nu M \mu)^{\vdim} \xi \|V_1\| \cball{y}{r}
        \,;
    \end{gather*}
    here we used~\ref{al:6.1(6)} with $r = \eta$. If $y \in Y_2$, then $f(y) \ge \rho + 3 \sigma >
    t + \sigma$, the open ball $\oball{y}{f(y) - t}$ contains $\cball{y}{\sigma}$, $f(x) > t$ on
    this ball, and the same reasoning with $V_2$ in place of~$V_1$ applies. Thus~\ref{al:6.1(9)}
    and~\ref{al:6.1(10)} hold. Since $V_j$ and~$V$ agree on $\oball{y}{\sigma} \times
    \pgrass{\adim}{\vdim}$ for $y \in Y_j$, \ref{al:6.1(11)} holds.

    Pick $y \in Y$. Since $A_1 \cup A_2 \subseteq \cball{y}{\eta + \diam Y} = \cball{y}{(M+1)
    \diam Y}$, \ref{al:6.1(7)} and~\eqref{al:6.1(13)} with $r = (M+1) \diam Y$ and $s = R$, or
    trivially if $r = R$, give
    \begin{displaymath}
        \frac{\|V_1\|(A_1) + \|V_2\|(A_2)}{\unitmeasure{\vdim} \eta^{\vdim}} \le \frac{\|V\|
        \cball{y}{(M+1) \diam Y}}{\unitmeasure{\vdim} \eta^{\vdim}} \le (1 + 1/M)^{\vdim} (\exp
        \xi) \frac{\|V\|(A)}{\unitmeasure{\vdim} R^{\vdim}} \,,
    \end{displaymath}
    because $\cball{y}{R} \subseteq A$. Finally, the definition of~$Y_0$, \eqref{al:6.1(13)}
    with $r = \eta$ and $s = R$, and~\ref{al:6.1(2)} give
    \begin{displaymath}
        \textsum{y \in Y_0}{} \density^{\vdim}(\|V\|,y) \le \frac{1}{M} \textsum{y \in Y_0}{}
        \frac{\|V\| \cball{y}{\eta}}{\unitmeasure{\vdim} \eta^{\vdim}} \le \frac{\nu + 1}{M} (\exp
        \xi) \frac{\|V\|(A)}{\unitmeasure{\vdim} R^{\vdim}} \,.
    \end{displaymath}
    Adding the last two estimates yields~\ref{al:6.1(12)}.
\end{proof}

\begin{remark}
    \label{rem:6.1}
    Allard assumes only that $Y$ has at most $\nu + 1$ elements. If $Y$ has fewer than two
    elements, then $\diam Y = 0$, the right side of~\ref{al:6.1}\ref{al:6.1(12)} is undefined, and
    the choice of $\tilde y$, $\tilde z$, $\sigma$ in the proof fails; this case is not needed,
    since~\ref{al:6.2} treats it directly. Allard also leaves implicit that $\|\delta V\|$ is
    a~Radon measure, which is needed for~\ref{al:4.10(2)}, and states~\ref{al:6.1(9)} with $\|V\|$
    instead of $\|V_j\|$ on the right; the proof gives $\|V_j\|$, which is what~\ref{al:6.2} needs.
\end{remark}

\begin{theorem}[{\cite[6.2]{Allard1972}}]
    \label{al:6.2}
    Corresponding to each $\lambda$, $\mu$ and~$\nu$ such that $1 < \lambda < \infty$, $1 \le \mu
    < \infty$ and $\nu$ is a~positive integer, there is $\gamma > 0$ with the following property.
    If $W$ is an open subset of~$\R^{\adim}$, $V \in \Var{\vdim}(W)$, $\|\delta V\|$ is a~Radon
    measure on~$W$, $T \in \pgrass{\adim}{\vdim}$, $Y \subseteq \R^{\adim}$, $0 < R < \infty$,
    $\R^{\adim} \cap \{ x : \dist(x,Y) \le R \} \subseteq W$, $Y$ has at most $\nu$ elements,
    $\diam Y \le \gamma R$,
    \begin{displaymath}
        |y - z| \le \mu |T^\perp(y - z)| \quad \text{whenever $y, z \in Y$} \,,
    \end{displaymath}
    and
    \begin{displaymath}
        R \|\delta V\| \cball{y}{r} \le \gamma \|V\| \cball{y}{r} \,, \qquad
        \textint{\cball{y}{r} \times \pgrass{\adim}{\vdim}}{} \|S - T\| \ud V(x,S) \le \gamma
        \|V\| \cball{y}{r}
    \end{displaymath}
    whenever $y \in Y$ and $0 < r < R$, then $\density^{\vdim}(\|V\|,y) \in \R$ for $y \in Y$ and
    \begin{displaymath}
        \lambda R^{-\vdim} \|V\|(\R^{\adim} \cap \{ x : \dist(x,Y) \le R \}) \ge
        \unitmeasure{\vdim} \textsum{y \in Y}{} \density^{\vdim}(\|V\|,y) \,.
    \end{displaymath}
\end{theorem}

\begin{proof}
    We say that $\gamma$ is admissible for $(\lambda,\mu,\nu)$ if it has the stated property. If
    $\gamma$ is admissible, so is every $\gamma'$ with $0 < \gamma' < \gamma$. Fix~$\mu$. We prove
    by induction on~$\nu$ that for every $\lambda > 1$ some $\gamma$ is admissible for
    $(\lambda,\mu,\nu)$.

    Suppose $\nu$ is a~positive integer, $1 < \lambda < \infty$, $0 < \gamma \le \log \lambda$,
    and $V$, $W$, $T$, $Y$, $R$ satisfy the hypotheses with $\gamma$, where $Y$ has at most one
    element. If $Y = \{ y \}$, then~\ref{new:mono} with $\Lambda = \gamma / R$ gives
    $\density^{\vdim}(\|V\|,y) \in \R$ and $\unitmeasure{\vdim} \density^{\vdim}(\|V\|,y) \le
    R^{-\vdim} \|V\| \cball{y}{R} \exp \gamma \le \lambda R^{-\vdim} \|V\| \cball{y}{R}$; if $Y =
    \varnothing$ there is nothing to prove. In particular every $\gamma$ with $0 < \gamma \le \log
    \lambda$ is admissible for $(\lambda,\mu,1)$.

    Suppose now that $\nu$ is a~positive integer and that for every $\lambda' > 1$ some
    $\gamma'$ is admissible for $(\lambda',\mu,\nu)$. Let $1 < \lambda < \infty$ and
    $\lambda' = \lambda^{1/2}$, and choose $\gamma'$ admissible for $(\lambda',\mu,\nu)$. Choose
    $M$ with $1 < M < \infty$, $M \gamma' \ge 1$ and $[(1 + 1/M)^{\vdim} + (\nu+1)/M] <
    \lambda^{1/4}$; then choose $\xi$ with $0 < \xi < 1$, $\exp \xi \le \lambda^{1/4}$,
    \begin{displaymath}
        2 M (\nu+1) (\exp \xi) (3 \nu M \mu)^{\vdim + 1} \xi \le \gamma' \quad \text{and} \quad M
        (\exp \xi) (3 \nu M \mu)^{\vdim} \xi \le \gamma' \,.
    \end{displaymath}
    Let $\gamma = \inf \{ \xi, (M+1)^{-1}, \log \lambda \}$. Suppose $V$, $W$, $T$, $Y$, $R$
    satisfy the hypotheses with $\nu + 1$ and~$\gamma$. If $Y$ has at most one element, we
    conclude by the preceding paragraph. Otherwise all hypotheses of~\ref{al:6.1} hold:
    \ref{al:6.1(4)} because $(M+1) \diam Y \le (M+1) \gamma R \le R$, and~\ref{al:6.1(5)},
    \ref{al:6.1(6)} because $\gamma \le \xi$. Let $V_1$, $V_2$, $Y_0$, $Y_1$, $Y_2$ be as
    in~\ref{al:6.1}, and $\eta = M \diam Y$. For $j \in \{1,2\}$ we verify the hypotheses of the
    present theorem for $V_j$, $W$, $T$, $Y_j$ and~$\eta$ in place of $V$, $W$, $T$, $Y$, $R$,
    with $(\lambda', \mu, \nu)$ and~$\gamma'$: $\|\delta V_j\|$ is a~Radon measure on~$W$; $\{ x :
    \dist(x,Y_j) \le \eta \} \subseteq \{ x : \dist(x,Y) \le R \} \subseteq W$ because $\eta <
    R$; $Y_j$ has at most $\nu$ elements by~\ref{al:6.1(8)}; $\diam Y_j \le \diam Y = \eta/M \le
    \gamma' \eta$; the condition on $|y - z|$ is inherited from~$Y$; and the two remaining
    inequalities with $\gamma'$ follow from~\ref{al:6.1(9)}, \ref{al:6.1(10)} and the choice
    of~$\xi$. Since $\gamma'$ is admissible, \ref{al:6.1(11)} gives
    \begin{displaymath}
        \unitmeasure{\vdim} \textsum{y \in Y_j}{} \density^{\vdim}(\|V\|,y) =
        \unitmeasure{\vdim} \textsum{y \in Y_j}{} \density^{\vdim}(\|V_j\|,y) \le \lambda'
        \eta^{-\vdim} \|V_j\|(\R^{\adim} \cap \{ x : \dist(x,Y_j) \le \eta \})
    \end{displaymath}
    for $j \in \{1,2\}$. Adding $\unitmeasure{\vdim} \textsum{y \in Y_0}{}
    \density^{\vdim}(\|V\|,y)$ and using~\ref{al:6.1(12)} and $\lambda' \ge 1$, we obtain
    \begin{displaymath}
        \unitmeasure{\vdim} \textsum{y \in Y}{} \density^{\vdim}(\|V\|,y) \le \lambda' [(1 +
        1/M)^{\vdim} + (\nu+1)/M] (\exp \xi) R^{-\vdim} \|V\|(\R^{\adim} \cap \{ x : \dist(x,Y)
        \le R \}) \,,
    \end{displaymath}
    and $\lambda' [(1 + 1/M)^{\vdim} + (\nu+1)/M] \exp \xi \le \lambda$. Hence $\gamma$ is
    admissible for $(\lambda,\mu,\nu+1)$.
\end{proof}

\begin{remark}
    \label{rem:6.2}
    Allard's proof reads ``to carry out the inductive step, use~6.2''; the intended reference
    is~\cite[6.1]{Allard1972}. His statement omits the quantifier on~$T$ and the hypothesis that
    $\|\delta V\|$ is a~Radon measure.
\end{remark}

\begin{lemma}[{\cite[6.3]{Allard1972}}]
    \label{al:6.3}
    Suppose $V_1, V_2, \ldots \in \IVar{\vdim}(\R^{\adim})$, $0 < d < \infty$, $T \in
    \pgrass{\adim}{\vdim}$, $\lim_{i \to \infty} V_i = d \var{}(\im T)$ in
    $\Var{\vdim}(\R^{\adim})$, $W$ is an open neighbourhood of~$0$ in~$\R^{\adim}$, and $\lim_{i
    \to \infty} \|\delta V_i\|(W) = 0$. Then $d$ is a~positive integer.
\end{lemma}

\begin{proof}
    Let $\nu$ be the smallest positive integer greater than~$d$, so that $\nu - 1 \le d < \nu$.
    Choose $R$ and~$\lambda$ with $0 < R < \infty$, $1 < \lambda < \infty$, $\cball{0}{\lambda
    R} \subseteq W$ and $\lambda^{\vdim+2} d < \nu$. Let $\gamma$ be as in~\ref{al:6.2} with
    $\lambda$, $\mu = 1$ and~$\nu$. Discarding finitely many~$i$ we assume $\|\delta V_i\|(W) <
    \infty$, so that $\|\delta V_i\|$ is a~Radon measure on~$W$; we write $V_i$ also for $V_i
    \restrict (W \times \pgrass{\adim}{\vdim}) \in \Var{\vdim}(W)$. Let $B = \cball{0}{(\lambda -
    1) R}$, and let $A_i$ be the set of those $x \in B$ such that $2 |T^\perp(x)| < \gamma R$,
    $\density^{\vdim}(\|V_i\|,x)$ is a~positive integer, and
    \begin{gather*}
        R \|\delta V_i\| \cball{x}{r} \le \gamma \|V_i\| \cball{x}{r} \,, \\
        \textint{\cball{x}{r} \times \pgrass{\adim}{\vdim}}{} \|S - T\| \ud V_i(x',S) \le \gamma
        \|V_i\| \cball{x}{r} \quad \text{whenever $0 < r < R$} \,.
    \end{gather*}
    Let $\psi_i$ be the Radon measure on~$W$ with $\psi_i(E) = R \|\delta V_i\|(E) +
    \textint{E \times \pgrass{\adim}{\vdim}}{} \|S - T\| \ud V_i(x,S)$ for $E \subseteq W$.
    By~\ref{al:3.5(1)}(a)(c), the set of $x \in B$ at which $\density^{\vdim}(\|V_i\|,x)$ is not
    a~positive integer is $\|V_i\|$ null. Each other point~$x$ of $(B \cap \spt \|V_i\| \cap \{ x
    : 2 |T^\perp(x)| < \gamma R \}) \without A_i$ is the centre of a~ball $\cball{x}{r_x}$ with
    $0 < r_x < R$ and $\gamma \|V_i\| \cball{x}{r_x} < \psi_i(\cball{x}{r_x})$, and
    $\cball{x}{r_x} \subseteq \cball{0}{\lambda R}$. By~\ref{al:2.7(1)} these balls admit
    disjointed subfamilies $\mathcal{C}_1, \ldots, \mathcal{C}_{\besicovitch{\adim}}$ covering
    the set of these points; each $\mathcal{C}_l$ is countable, because its members have positive
    $\|V_i\|$ measure and lie in the compact set $\cball{0}{\lambda R}$. Hence
    \begin{equation}
        \label{eq:6.3:1}
        \begin{multlined}
        \|V_i\|((B \cap \{ x : 2 |T^\perp(x)| < \gamma R \}) \without A_i) \\ \le \gamma^{-1}
        \besicovitch{\adim} \Big[ R \|\delta V_i\| \cball{0}{\lambda R} + \textint{\cball{0}{\lambda
        R} \times \pgrass{\adim}{\vdim}}{} \|S - T\| \ud V_i(x,S) \Big] \,.
        \end{multlined}
    \end{equation}
    The function $[(x,S) \mapsto \|S - T\|]$ is continuous and vanishes on $\spt (d \var{}(\im
    T)) \subseteq \im T \times \{ T \}$; by~\ref{al:2.6(2)(a)}(c) applied to the measures $\|S -
    T\| V_i$, the integral on the right of~\eqref{eq:6.3:1} tends to~$0$. Also $\|\delta V_i\|
    \cball{0}{\lambda R} \le \|\delta V_i\|(W) \to 0$. Moreover $K = B \cap \{ x : 2 |T^\perp(x)|
    \ge \gamma R \}$ is compact and disjoint from $\im T$, and $\|V_i\| \to d \HM^{\vdim} \restrict
    \im T$, so $\limsup_{i \to \infty} \|V_i\|(K) \le d \HM^{\vdim}(K \cap \im T) = 0$
    by~\ref{al:2.6(2)(a)}(c). Consequently $\lim_{i \to \infty} \|V_i\|(B \without A_i) = 0$.

    The sets~$A_i$ are $\|V_i\|$ measurable: the functions $[x \mapsto \|V_i\| \cball{x}{r}]$
    and $[x \mapsto \psi_i(\cball{x}{r})]$ are upper semicontinuous, both sides are continuous
    from the right in~$r$, so the conditions for $0 < r < R$ may be required for rational~$r$
    only, and $\density^{\vdim}(\|V_i\|,\cdot)$ is $\|V_i\|$ measurable by~\ref{al:5.1(4)}. Let
    $W_i = V_i \restrict (A_i \times \pgrass{\adim}{\vdim})$ and $W_\infty = d \var{}(\im T \cap
    B)$. Since $d \HM^{\vdim}(\im T \cap \Bdry B) = 0$, \ref{al:2.6(2)(a)}(d) applied to the
    measures $f^+ V_i$ and $f^- V_i$ gives $\lim_{i \to \infty} (V_i \restrict (B \times
    \pgrass{\adim}{\vdim}))(f) = W_\infty(f)$ for $f \in \ccspace{\R^{\adim} \times
    \pgrass{\adim}{\vdim}}$; since $|(V_i \restrict (B \times \pgrass{\adim}{\vdim}))(f) -
    W_i(f)| \le \sup |f| \, \|V_i\|(B \without A_i)$, we conclude $\lim_{i \to \infty} W_i =
    W_\infty$.

    We claim $\lim_{i \to \infty} T_\# W_i = T_\# W_\infty = W_\infty$. Let $\zeta \in
    \ccspace{\R^{\adim}}$ with $0 \le \zeta \le 1$ and $\zeta = 1$ on~$B$. For $\alpha \in
    \ccspace{\R^{\adim} \times \pgrass{\adim}{\vdim}}$ let $\beta(x,S) = \zeta(x) \alpha(T(x),
    P(S)) |\extpower{\vdim}(T \circ S)|$ if $\extpower{\vdim}(T \circ S) \ne 0$, where $P(S)$ is
    the projection onto $T[\im S]$, and $\beta(x,S) = 0$ otherwise. Since $P$ is continuous on the
    open set of those~$S$ with $\extpower{\vdim}(T \circ S) \ne 0$ and $|\extpower{\vdim}(T
    \circ S)|$ is continuous and vanishes elsewhere, $\beta \in \ccspace{\R^{\adim} \times
    \pgrass{\adim}{\vdim}}$. By~\ref{al:3.2(1)}, $T_\# W_i(\alpha) = W_i(\beta)$ and $T_\#
    W_\infty(\alpha) = W_\infty(\beta)$, because $\spt \|W_i\| \cup \spt \|W_\infty\| \subseteq
    B$; this proves the claim.

    For $z \in \im T$ let $Y_i(z) = A_i \cap T^{-1}\{ z \}$. If $y, y' \in Y_i(z)$, then $|y - y'|
    = |T^\perp(y - y')| < \gamma R$. Since $\R^{\adim} \cap \{ x : \dist(x, Y_i(z)) \le R \}
    \subseteq \cball{0}{\lambda R}$ and $\lim_{i \to \infty} \|V_i\| \cball{0}{\lambda R} = d
    \unitmeasure{\vdim} (\lambda R)^{\vdim}$ by~\ref{al:2.6(2)(a)}(d), there is~$i_0$ such that
    \begin{displaymath}
        \frac{\|V_i\|(\R^{\adim} \cap \{ x : \dist(x,Y_i(z)) \le R \})}{\unitmeasure{\vdim}
        R^{\vdim}} < \lambda^{\vdim+1} d \quad \text{for $z \in \im T$ and $i \ge i_0$} \,.
    \end{displaymath}
    Suppose $i \ge i_0$, $z \in \im T$ and $Y \subseteq Y_i(z)$ has at most $\nu$ elements. All
    hypotheses of~\ref{al:6.2} hold for $V_i$, $W$, $T$, $Y$, $R$, with $\mu = 1$: in particular
    $\{ x : \dist(x,Y) \le R \} \subseteq \cball{0}{\lambda R} \subseteq W$, $\diam Y \le \gamma
    R$, and $|y - y'| = |T^\perp(y - y')|$ for $y, y' \in Y$. Hence $\textsum{y \in Y}{}
    \density^{\vdim}(\|V_i\|,y) \le \lambda^{\vdim+2} d < \nu$. Since the densities at points
    of~$A_i$ are positive integers, $Y_i(z)$ has fewer than $\nu$ elements, and applying the
    estimate with $Y = Y_i(z)$ gives $\textsum{y \in Y_i(z)}{} \density^{\vdim}(\|V_i\|,y) \le
    \nu - 1$.

    Let $i \ge i_0$. For $y \in A_i$ we have $\density^{\vdim}(\|W_i\|,y) \le
    \density^{\vdim}(\|V_i\|,y)$ whenever the former exists, because $\|W_i\| \le \|V_i\|$. The set
    $N_i = \R^{\adim} \cap \{ y : y \notin A_i,\ \udensity^{\vdim}(\|W_i\|,y) > 0 \}$ satisfies
    $\HM^{\vdim}(N_i) = 0$ by~\ref{al:2.8(2)}(c), since $\|W_i\| = \|V_i\| \restrict A_i$ and
    $\|V_i\|(A_i) < \infty$; hence $\HM^{\vdim}(T[N_i]) = 0$ by~\cite[2.10.11]{Federer1969}.
    Since $W_i \in \IVar{\vdim}(\R^{\adim})$ by~\ref{al:3.5(1)}(c), \ref{al:3.5(3)}
    and~\ref{al:3.5(1)}(b) give
    \begin{multline*}
        \|T_\# W_i\|(\R^{\adim}) = \textint{}{} \density^{\vdim}(\|T_\# W_i\|,z) \ud \HM^{\vdim}
        z = \textint{\im T \without T[N_i]}{} \textsum{y \in T^{-1}\{z\}}{}
        \density^{\vdim}(\|W_i\|,y) \ud \HM^{\vdim} z \\ \le \textint{T[A_i]}{} \textsum{y \in
        Y_i(z)}{} \density^{\vdim}(\|V_i\|,y) \ud \HM^{\vdim} z \le (\nu - 1)
        \HM^{\vdim}(T[A_i]) \le (\nu - 1) \unitmeasure{\vdim} ((\lambda - 1) R)^{\vdim} \,,
    \end{multline*}
    where we used $T[A_i] \subseteq \im T \cap B$ and $\HM^{\vdim}(\im T \cap B) =
    \unitmeasure{\vdim} ((\lambda-1)R)^{\vdim}$ by~\ref{al:2.8(1)}. All $T_\# W_i$ and~$W_\infty$
    are supported in the compact set $B \times \pgrass{\adim}{\vdim}$, so the claim gives $\lim_{i
    \to \infty} \|T_\# W_i\|(\R^{\adim}) = \|W_\infty\|(\R^{\adim}) = d \unitmeasure{\vdim}
    ((\lambda - 1) R)^{\vdim}$. Therefore $d \le \nu - 1$, and $d = \nu - 1$ is a~positive
    integer.
\end{proof}

\begin{remark}
    \label{rem:6.3}
    Allard concludes that $d$ is a~nonnegative integer; since $d > 0$ it is positive. In his
    proof the multiplicity formula is attributed to ``3.5(3) and 2.8(5)''; the formula used
    is~\ref{al:3.5(3)}, which rests on the area formula~\ref{al:2.8(6)}, while 2.8(5) is the
    rectifiability theorem. The passage from $V_i \restrict (A_i \times \pgrass{\adim}{\vdim})$ to
    its image under~$T$ and the treatment of points outside~$A_i$ are added here.
\end{remark}

\begin{theorem}[{\cite[6.4]{Allard1972}}]
    \label{al:6.4}
    Suppose $U$ is an open subset of~$\R^{\adim}$, $G_1, G_2, \ldots$ are open subsets of~$U$, $U
    = \tbcup_{i=1}^\infty G_i$, and $M_1, M_2, \ldots$ are nonnegative real numbers. Then
    \begin{displaymath}
        \mathcal{I} = \IVar{\vdim}(U) \cap \{ V : (\|V\| + \|\delta V\|)(G_i) \le M_i \text{ for
        } i = 1, 2, \ldots \}
    \end{displaymath}
    is compact.
\end{theorem}

\begin{proof}
    Let $\mathcal{M}$ be the set of Radon measures~$\nu$ on $U \times \pgrass{\adim}{\vdim}$ with
    $\nu(G_i \times \pgrass{\adim}{\vdim}) \le M_i$ for every~$i$. By~\ref{al:2.6(2)(a)}(a)
    and~\ref{new:weak_metric}, $\mathcal{M}$ is compact and its topology is induced by a~metric.
    Since $\mathcal{I} \subseteq \mathcal{M}$, it suffices to show that $V \in \mathcal{I}$
    whenever $V_1, V_2, \ldots \in \mathcal{I}$ and $\lim_{j \to \infty} V_j = V$. Each~$V_j$
    belongs to the set~$\mathcal{R}$ of~\ref{al:5.6} with $\vartheta = 1$,
    by~\ref{al:3.5(1)}(a)(c); since $\mathcal{R}$ is compact and the weak topology is Hausdorff,
    $V \in \mathcal{R}$. Thus $V \in \RVar{\vdim}(U)$ and $(\|V\| + \|\delta V\|)(G_i) \le M_i$
    for every~$i$; in particular $\|\delta V\|$ and every $\|\delta V_j\|$ are Radon measures
    on~$U$, as every compact subset of~$U$ is covered by finitely many~$G_i$. By~\ref{al:3.5(1)}(c)
    it remains to show that $\density^{\vdim}(\|V\|,a)$ is a~positive integer for $\|V\|$ almost
    all $a \in U$.

    Fix~$i$. The weights $\|V_j\|$ converge to~$\|V\|$, because $[(x,S) \mapsto f(x)] \in
    \ccspace{U \times \pgrass{\adim}{\vdim}}$ for $f \in \ccspace{U}$. Applying~\ref{al:2.7(2)}
    with $U$, $\mu_j$, $\nu_j$ replaced by $G_i$, $\|V_j\| \restrict G_i$, $\|\delta V_j\|
    \restrict G_i$, and noting $\|\delta V_j\|(G_i) \le M_i$, we see that for $\|V\|$ almost all
    $a \in G_i$ there is a~positive integer~$p$ with $\liminf_{j \to \infty} \dist(a, \spt(\|V_j\|
    \restrict A_{j,p})) = 0$, where
    \begin{displaymath}
        A_{j,p} = G_i \cap \{ x : \|\delta V_j\| \cball{x}{r} \le p \|V_j\| \cball{x}{r} \text{
        whenever } 0 < r < \dist(x, \R^{\adim} \without G_i) \} \,.
    \end{displaymath}
    Since $V \in \RVar{\vdim}(U)$ and $\|\delta V\|$ is a~Radon measure, \ref{al:3.5(1)}(a),
    \ref{al:5.5(1)} and its proof show that for $\|V\|$ almost all $a \in U$ we have $0 <
    \density^{\vdim}(\|V\|,a) < \infty$ and every member of $\VarTan(V,a)$ equals
    $\density^{\vdim}(\|V\|,a) \var{}(\im T)$ for some $T \in \pgrass{\adim}{\vdim}$; moreover
    $\VarTan(V,a) \ne \varnothing$ by~\ref{al:3.4(1)}(a). As $U = \tbcup_{i=1}^\infty G_i$, it
    suffices to show that $\theta = \density^{\vdim}(\|V\|,a)$ is a~positive integer whenever $a
    \in G_i$ has all these properties.

    Fix such $a$ and~$p$, choose $\rho > 0$ with $\cball{a}{3\rho} \subseteq G_i$, choose $C \in
    \VarTan(V,a)$, $T \in \pgrass{\adim}{\vdim}$ with $C = \theta \var{}(\im T)$, and positive
    numbers $r_1 < r_2 < \ldots$ with $\lim_{m \to \infty} r_m = \infty$ as in~\ref{al:3.4},
    so that $\lim_{m \to \infty} r_m^{\vdim} \textint{}{} \alpha(r_m(x - a), S) \ud V(x,S) =
    C(\alpha)$ for $\alpha \in \ccspace{\R^{\adim} \times \pgrass{\adim}{\vdim}}$. For $m = 1, 2,
    \ldots$ and $V' \in \Var{\vdim}(U)$ let $Z_m(V') = (\scale{r_m} \circ \trans{-a})_\# (V'
    \restrict (\oball{a}{\rho} \times \pgrass{\adim}{\vdim})) \in \Var{\vdim}(\R^{\adim})$, see
    \ref{al:3.2(1)}, \ref{al:3.2(2)}. Then $Z_m(V')(\alpha) = r_m^{\vdim} \textint{}{}
    \alpha(r_m(x-a), S) \ud V'(x,S)$ whenever $\spt \alpha \subseteq \oball{0}{r_m \rho} \times
    \pgrass{\adim}{\vdim}$, and for such~$\alpha$ the function $[(x,S) \mapsto \alpha(r_m(x-a),
    S)]$ belongs to $\ccspace{U \times \pgrass{\adim}{\vdim}}$, so that $\lim_{j \to \infty}
    Z_m(V_j)(\alpha) = Z_m(V)(\alpha)$. Let $D = \{ d_1, d_2, \ldots \}$ be as
    in~\ref{new:countable_dense} with $X = \R^{\adim} \times \pgrass{\adim}{\vdim}$. Using these
    facts and the choice of~$p$ we choose inductively integers $\lambda(1) < \lambda(2) < \ldots$
    and points $a_{\lambda(m)}$ such that
    \begin{displaymath}
        |Z_m(V_{\lambda(m)})(d_l) - Z_m(V)(d_l)| < 1/m \quad \text{whenever $l \le m$ and $\spt d_l
        \subseteq \oball{0}{r_m \rho} \times \pgrass{\adim}{\vdim}$} \,,
    \end{displaymath}
    and $a_{\lambda(m)} \in A_{\lambda(m),p} \cap \spt \|V_{\lambda(m)}\| \cap \oball{a}{1/r_m}$;
    the second condition can be met for infinitely many indices because every point of $\spt
    (\|V_j\| \restrict A_{j,p})$ is a~limit of points of $A_{j,p} \cap \spt \|V_j\|$. Let $W_m =
    Z_m(V_{\lambda(m)})$ and $b_m = a_{\lambda(m)}$. Each $W_m$ belongs to
    $\IVar{\vdim}(\R^{\adim})$ by~\ref{al:3.5}, \ref{al:3.5(3)}.

    For large~$m$ we have $1/r_m < \rho$ and $\dist(b_m, \R^{\adim} \without G_i) > 2 \rho$, so
    the definition of $A_{\lambda(m),p}$ and~\ref{new:mono} with $W = G_i$, $R = \rho$ and $\Lambda
    = p$ give
    \begin{displaymath}
        r^{-\vdim} \|V_{\lambda(m)}\| \cball{b_m}{r} \le \rho^{-\vdim} \|V_{\lambda(m)}\|
        \cball{b_m}{\rho} \exp(p \rho) \le \rho^{-\vdim} M_i \exp(p\rho) \quad \text{for $0 < r \le
        \rho$} \,,
    \end{displaymath}
    because $\cball{b_m}{\rho} \subseteq \cball{a}{2\rho} \subseteq G_i$. For $0 < s < \infty$
    and large~$m$ this yields $\|W_m\| \cball{0}{s} = r_m^{\vdim} \|V_{\lambda(m)}\|
    \cball{a}{s/r_m} \le r_m^{\vdim} \|V_{\lambda(m)}\| \cball{b_m}{(s+1)/r_m} \le (s+1)^{\vdim}
    \rho^{-\vdim} M_i \exp(p\rho)$. Since $\lim_{m \to \infty} Z_m(V)(d_l) = C(d_l)$ for
    every~$l$, the choice of $\lambda(m)$ gives $\lim_{m \to \infty} W_m(d_l) = C(d_l)$ for
    every~$l$; together with the bound on $\|W_m\| \cball{0}{s}$ and the approximation property
    of~$D$ in~\ref{new:countable_dense}, this gives $\lim_{m \to \infty} W_m = C = \theta
    \var{}(\im T)$.

    For $g \in \VF(\R^{\adim})$ with $\spt g \subseteq \oball{0}{1}$ and $1/r_m < \rho$, the
    computation in the proof of~\ref{al:4.12} gives $\delta W_m(g) = r_m^{\vdim-1} \delta
    V_{\lambda(m)}(g \circ \scale{r_m} \circ \trans{-a})$, and $\spt(g \circ \scale{r_m} \circ
    \trans{-a}) \subseteq \oball{a}{1/r_m} \subseteq \oball{b_m}{2/r_m}$. Hence, for large~$m$,
    \begin{multline*}
        \|\delta W_m\| \oball{0}{1} \le r_m^{\vdim-1} \|\delta V_{\lambda(m)}\|
        \cball{b_m}{2/r_m} \le r_m^{\vdim-1} p \|V_{\lambda(m)}\| \cball{b_m}{2/r_m} \\ \le
        r_m^{\vdim-1} p (2/r_m)^{\vdim} \rho^{-\vdim} M_i \exp(p\rho) = 2^{\vdim} p r_m^{-1}
        \rho^{-\vdim} M_i \exp(p\rho) \,,
    \end{multline*}
    which tends to~$0$. By~\ref{al:6.3} with $W = \oball{0}{1}$, $\theta$ is a~positive integer.
\end{proof}

\begin{remark}
    \label{rem:6.4}
    Allard assumes $U = \R^{\adim}$ and remarks that the general case follows by a~trivial
    modification; the localisation to~$G_i$ and the cut-off $V' \restrict (\oball{a}{\rho} \times
    \pgrass{\adim}{\vdim})$ above carry it out. The diagonal choice of $\lambda(m)$ is stated
    without details in~\cite[6.4]{Allard1972}.
\end{remark}

\begin{corollary}[{\cite[6.4, Corollary]{Allard1972}}]
    \label{al:6.4:cor}
    Corresponding to every $\gamma$ and~$\varepsilon$ with $0 < \gamma < \infty$ and $0 <
    \varepsilon < 1$ there is $\delta > 0$ with the following property. If $1 \le m \le \adim$, $T
    \in \pgrass{\adim}{m}$, $M$ is a~smooth $m$~dimensional submanifold of~$\R^{\adim}$ with
    $\oball{0}{1} \cap (\Clos M \without M) = \varnothing$,
    \begin{gather*}
        \HM^m(M \cap \oball{0}{1}) < \gamma \,, \qquad \textint{M \cap \oball{0}{1}}{} |h(M;x)|^{1
        + \varepsilon} \ud \HM^m x < \gamma \,, \\ \sup \{ \dist(x, \im T) : x \in M \cap
        \oball{0}{1} \} < \delta \,,
    \end{gather*}
    then $\HM^m(M \cap \oball{0}{1-\varepsilon}) / \HM^m(\im T \cap \oball{0}{1-\varepsilon})$ is
    within~$\varepsilon$ of some integer.
\end{corollary}

\begin{proof}
    Suppose the assertion fails for some $\gamma$ and~$\varepsilon$. Then for $i = 1, 2, \ldots$
    there are $T_i$ and~$M_i$ satisfying the hypotheses with $\delta = 1/i$ for which the quotient
    has distance at least~$\varepsilon$ from every integer. Passing to a~subsequence we may assume
    that $m$ is the same for all~$i$ and, by compactness of $\pgrass{\adim}{m}$, that $\lim_{i \to
    \infty} T_i = T$. Let $V_i = \var{}(M_i \cap \oball{0}{1}) \in \IVar{m}(\oball{0}{1})$; this
    is defined because $M_i \cap \oball{0}{1}$ is a~relatively closed submanifold
    of~$\oball{0}{1}$ of finite $\HM^m$ measure. Since $M_i \cap \oball{0}{1}$ has no boundary
    in~$\oball{0}{1}$ and $h(M_i;x) \in \Nor(M_i,x)$, \ref{al:4.4} and~\ref{new:divergence},
    the latter applied after a~partition of unity~\cite[3.1.13]{Federer1969} subordinate to
    charts, give
    \begin{displaymath}
        \delta V_i(g) = - \textint{M_i \cap \oball{0}{1}}{} g(x) \bullet h(M_i;x) \ud \HM^m x
        \quad \text{for $g \in \VF(\oball{0}{1})$} \,.
    \end{displaymath}
    Hence $\|\delta V_i\| = \HM^m \restrict (|h(M_i;\cdot)| \CF{M_i \cap \oball{0}{1}})$
    and, by H\"older's inequality~\cite[2.4.14]{Federer1969}, $\|\delta V_i\|(\oball{0}{1}) \le
    \gamma^{1/(1+\varepsilon)} \gamma^{\varepsilon/(1+\varepsilon)} = \gamma$. By~\ref{al:6.4}
    with $U = G_1 = \oball{0}{1}$ and $M_1 = 2\gamma$, a~subsequence of $V_1, V_2, \ldots$, still
    denoted $V_i$, converges to some $V \in \IVar{m}(\oball{0}{1})$, and $\|\delta
    V\|(\oball{0}{1}) \le \gamma$ by~\ref{al:4.11}.

    If $K$ is a~compact subset of $\oball{0}{1} \without \im T$, then $\dist(K, \im T) > 0$ and
    $|\dist(x, \im T) - \dist(x, \im T_i)| \le \|T - T_i\|$ for $|x| \le 1$, so $K \cap \spt
    \|V_i\| = \varnothing$ for large~$i$; by~\ref{al:2.6(2)(a)}(c), $\|V\|(\Int K) = 0$. Thus $\spt
    \|V\| \subseteq \im T$. Let $u \in \im T^\perp$ and $\zeta \in \dspace{\oball{0}{1}}{\R}$, and
    let $f(x) = \zeta(x) (u \bullet x)$. Then $f \in \dspace{\oball{0}{1}}{\R}$, $\spt \|V\|
    \subseteq \{ x : f(x) = 0 \}$, and $\uD f(x)(v) = \zeta(x) (u \bullet v)$ for $x \in \im T$;
    by~\ref{al:4.6(1)}, $u \bullet v = 0$ for $v \in \im S$ whenever $\zeta(x) \ne 0$, for $V$
    almost all $(x,S)$. Letting $u$ run through a~basis of~$\im T^\perp$ and $\zeta$ through
    a~sequence with $\tbcup_i \{ x : \zeta_i(x) \ne 0 \} = \oball{0}{1}$, we obtain $\im S
    \subseteq \im T$, hence $S = T$, for $V$ almost all $(x,S)$.

    Suppose $g \in \VF(\oball{0}{1})$ and $g(x) \in \im T$ for $x \in \oball{0}{1}$. Let $q = (1 +
    \varepsilon)/\varepsilon$ and choose $\phi \in \ccspace{\oball{0}{1}}$ with $0 \le \phi \le 1$
    and $\phi = 1$ on $\spt g$. Since $h(M_i;x) \in \im \Tan(M_i,x)^\perp$ and $g(x) = T(g(x))$,
    we have $g(x) \bullet h(M_i;x) = (T - S_x)(g(x)) \bullet h(M_i;x)$ with $S_x$ the projection
    onto $\Tan(M_i,x)$, hence by H\"older's inequality
    \begin{displaymath}
        |\delta V_i(g)| \le \sup |g| \, \Big( \textint{}{} \phi(x) \|T - S\|^q \ud V_i(x,S)
        \Big)^{1/q} \gamma^{1/(1+\varepsilon)} \,.
    \end{displaymath}
    Since $[(x,S) \mapsto \phi(x) \|T - S\|^q] \in \ccspace{\oball{0}{1} \times
    \pgrass{\adim}{m}}$ and $S = T$ for $V$ almost all $(x,S)$, the right side tends to~$0$; thus
    $\delta V(g) = \lim_{i \to \infty} \delta V_i(g) = 0$. By~\ref{al:4.6(3)} with the connected
    submanifold $\im T \cap \oball{0}{1}$ there is $0 \le c < \infty$ with $V = c \var{}(\im T
    \cap \oball{0}{1})$, and $c$ is an integer by~\ref{al:3.5(1)}(c), because
    $\density^m(\|V\|,x) = c$ for $x \in \im T \cap \oball{0}{1}$. Finally $\|V\|(\Bdry
    \oball{0}{1-\varepsilon}) = 0$, so~\ref{al:2.6(2)(a)}(d) gives
    \begin{displaymath}
        \lim_{i \to \infty} \HM^m(M_i \cap \oball{0}{1-\varepsilon}) = \|V\|
        \oball{0}{1-\varepsilon} = c \HM^m(\im T \cap \oball{0}{1-\varepsilon}) \,,
    \end{displaymath}
    and $\HM^m(\im T_i \cap \oball{0}{1-\varepsilon}) = \HM^m(\im T \cap
    \oball{0}{1-\varepsilon})$ by~\ref{al:2.8(1)}. This contradicts the choice of $M_i$.
\end{proof}

\begin{remark}
    \label{rem:6.4:cor}
    Allard's proof reads ``Apply 4.7, 4.4, the preceding theorem, and 4.5''. Since $M \cap
    \oball{0}{1}$ has no boundary in $\oball{0}{1}$, the exterior normal of~\ref{al:4.7}, and
    with it~\cite[4.5.6, 4.5.12]{Federer1969}, is not needed. The hypothesis on $\textint{}{}
    |h(M;x)|^{1+\varepsilon} \ud \HM^m x$, rather than $\textint{}{} |h(M;x)| \ud \HM^m x$, is
    used to show that the limit is stationary with respect to fields tangent to~$\im T$.
\end{remark}

\begin{example}[{\cite[6.4, Remark]{Allard1972}}]
    \label{al:6.4:rem}
    Let $\adim = 4$, $m = 2$, $T \in \pgrass{4}{2}$ with $\im T = \R^4 \cap \{ x : x_3 = x_4 = 0
    \}$, and identify $\R^4$ with $\C^2$ through $Z_1(x) = x_1 + \sqrt{-1} x_2$ and $Z_2(x) = x_3
    + \sqrt{-1} x_4$. For $0 < t < \infty$ let $M_t = \R^4 \cap \{ x : Z_1(x) = t Z_2(x)^2 \}$.
    Then $\var{}(M_t) \in \IVar{2}(\R^4)$ is stationary and $\lim_{t \to \infty} \var{}(M_t) = 2
    \var{}(\im T)$. Thus in~\ref{al:6.4:cor} the integer may exceed~$1$ although the~$M$ are
    connected minimal surfaces lying close to~$\im T$ on compact sets.
\end{example}

\begin{proof}
    Let $P_t(\zeta) = (\zeta^2, t^{-1/2} \zeta) \in \C^2$ for $\zeta \in \C$. Then $P_t$ is an
    injective holomorphic map with $P_t[\C] = M_t$, $P_t$ is proper, and $P_t'(\zeta) = (2\zeta,
    t^{-1/2}) \ne 0$. By the Cauchy--Riemann equations, with $\zeta = u + \sqrt{-1} v$ the
    partial derivatives satisfy $|\partial_u P_t| = |\partial_v P_t| = |P_t'(\zeta)|$ and
    $\partial_u P_t \bullet \partial_v P_t = 0$, so $|\extpower{2} \uD P_t(\zeta)| = |P_t'(\zeta)|^2 =
    4|\zeta|^2 + t^{-1}$, and the components of~$P_t$ are harmonic. Suppose $g \in \VF(\R^4)$ and
    $F_s = P_t + s\, g \circ P_t$ for $s \in \R$. For small~$|s|$ the map $\id{\R^4} + s g$ is
    a~diffeomorphism, so by~\ref{al:4.1} and the area formula~\cite[3.2.3, 3.2.5]{Federer1969}
    \begin{displaymath}
        \delta \var{}(M_t)(g) = \frac{\ud}{\ud s} A(s) \Big|_{s=0} \,, \qquad A(s) = \textint{}{}
        |\partial_u F_s \wedge \partial_v F_s| \ud \LM^2 \,,
    \end{displaymath}
    where all integrands vanish outside the compact set $P_t^{-1}[\spt g]$. Let $E(s) =
    \frac{1}{2} \textint{}{} (|\partial_u F_s|^2 + |\partial_v F_s|^2) \ud \LM^2$. Since $|a \wedge
    b| \le |a| |b| \le \frac12 (|a|^2 + |b|^2)$, we have $A \le E$, with equality at $s = 0$ by
    conformality; both functions are differentiable, so $A'(0) = E'(0)$. Integration by parts on
    lines~\cite[2.6.2, 2.9.20]{Federer1969} gives $E'(0) = - \textint{}{} (\partial_u^2 P_t +
    \partial_v^2 P_t) \bullet g(P_t) \ud \LM^2 = 0$. Hence $\var{}(M_t)$ is stationary.

    Let $\alpha \in \ccspace{\R^4 \times \pgrass{4}{2}}$, $\spt \alpha \subseteq \cball{0}{\rho}
    \times \pgrass{4}{2}$. The area formula gives $\var{}(M_t)(\alpha) = \textint{}{}
    \alpha(P_t(\zeta), Q_t(\zeta)) (4 |\zeta|^2 + t^{-1}) \ud \LM^2 \zeta$, where $Q_t(\zeta)$ is
    the projection onto the complex line $\C P_t'(\zeta)$. The integrand vanishes for $|\zeta|^2 >
    \rho$ and is bounded independently of $t \ge 1$. As $t \to \infty$, $P_t(\zeta) \to
    (\zeta^2,0)$ and, for $\zeta \ne 0$, $Q_t(\zeta) \to T$. By the bounded convergence
    theorem~\cite[2.4.9]{Federer1969}, $\var{}(M_t)(\alpha) \to \textint{}{} \alpha((\zeta^2,0),T)
    4 |\zeta|^2 \ud \LM^2 \zeta$, which equals $2 \textint{\im T}{} \alpha(z,T) \ud \HM^2 z =
    2 \var{}(\im T)(\alpha)$ by the area formula applied to the two-to-one map $[\zeta \mapsto
    (\zeta^2,0)]$, whose Jacobian is $4 |\zeta|^2$.
\end{proof}

\begin{theorem}[{\cite[6.5]{Allard1972}}]
    \label{al:6.5}
    Suppose $U$ is an open subset of~$\R^{\adim}$, $V \in \Var{\vdim}(U)$, $\|\delta V\|$ is
    a~Radon measure on~$U$, $a \in U$, $\density^{\vdim}(\|V\|,a)$ exists and is finite, and
    $\density^{\vdim-1}(\|\delta V\|,a) = 0$. Suppose $C \in \VarTan(V,a)$. Then:
    \begin{enumerate}
        \item if $\density^{\vdim}(\|V\|,a) > 0$, then for every $0 < R < \infty$ the varifold $C
            \restrict (\oball{0}{R} \times \pgrass{\adim}{\vdim})$ satisfies the hypotheses
            of~\ref{al:5.2(2)};
        \item if $0 < d < \infty$ and $\density^{\vdim}(\|V\| \restrict \{ x :
            \density^{\vdim}(\|V\|,x) < d \}, a) = 0$, then $C \in \RVar{\vdim}(\R^{\adim})$ and
            $\scale{r \#} C = C$ for $0 < r < \infty$;
        \item if $V \in \IVar{\vdim}(U)$, then $C \in \IVar{\vdim}(\R^{\adim})$.
    \end{enumerate}
\end{theorem}

\begin{proof}
    Let $r_1, r_2, \ldots$ and $V_i = ((\scale{r_i} \circ \trans{-a})|U)_\# V$ be as
    in~\ref{al:3.4} with $\lim_{i \to \infty} V_i = C$; here $\Tan(U,a) = \R^{\adim}$. Fix $0 < R
    < \infty$; for large~$i$, $\oball{0}{R} \subseteq r_i(U - a)$, and we write $V_i$ for the
    restriction of~$V_i$ to $\oball{0}{R} \times \pgrass{\adim}{\vdim}$. As in the proof
    of~\ref{al:4.12}, $\|\delta V_i\| \oball{0}{R} = r_i^{\vdim-1} \|\delta V\| \oball{a}{R/r_i}$,
    which tends to~$0$ because $\density^{\vdim-1}(\|\delta V\|,a) = 0$; and $\|V_i\| \oball{0}{R}
    = r_i^{\vdim} \|V\| \oball{a}{R/r_i}$ is bounded.

    By~\ref{al:4.11}, $C$ is stationary, and by~\ref{al:3.4(1)}(b), $\|C\| \cball{0}{r} =
    \density^{\vdim}(\|V\|,a) \unitmeasure{\vdim} r^{\vdim}$ for $0 < r < \infty$. Hence
    $\density^{\vdim}(\|C\|,0) = \density^{\vdim}(\|V\|,a)$ and $\|C\| \oball{0}{R} \le
    \density^{\vdim}(\|C\|,0) \unitmeasure{\vdim} R^{\vdim}$. If $\density^{\vdim}(\|V\|,a) > 0$,
    then $0 \in \spt \|C\|$, which proves~(a).

    Under the hypothesis of~(b), $\density^{\vdim}(\|V_i\|, r_i(x-a)) =
    \density^{\vdim}(\|V\|,x)$ for $x \in U$ by~\ref{al:3.2(2)}, hence
    \begin{displaymath}
        \|V_i\|(\oball{0}{R} \cap \{ y : \density^{\vdim}(\|V_i\|,y) < d \}) = r_i^{\vdim}
        \|V\|(\oball{a}{R/r_i} \cap \{ x : \density^{\vdim}(\|V\|,x) < d \}) \to 0 \,.
    \end{displaymath}
    By~\ref{al:5.4} with $W = \oball{0}{R}$, $\ldensity^{\vdim}(\|C\|,y) \ge d$ for $\|C\|$
    almost all $y \in \oball{0}{R}$, and~\ref{al:5.5(1)}, applied to the stationary varifold $C
    \restrict (\oball{0}{R} \times \pgrass{\adim}{\vdim})$, shows that this varifold is
    rectifiable. As $R$ is arbitrary, \ref{al:3.5(1)}(a) gives $C \in \RVar{\vdim}(\R^{\adim})$.
    If $\density^{\vdim}(\|V\|,a) = 0$ then $C = 0$. Otherwise~(a) and~\ref{al:5.2(2)}(b) give
    $\|C\|(A) = \lambda^{-\vdim} \|C\|(\{ \lambda x : x \in A \})$ for $0 < \lambda < 1$ and
    bounded~$A$; hence $\|\scale{r \#} C\| = r^{\vdim} \scale{r \#} \|C\| = \|C\|$
    by~\ref{al:3.2(2)}. By~\ref{al:3.5(3)}, $\scale{r \#} C$ is rectifiable and
    $\Tan^{\vdim}(\|\scale{r \#} C\|, r y) = \Tan^{\vdim}(\|C\|, y)$ for $\|C\|$ almost all~$y$,
    since $\uD \scale{r}(y) = r \id{\R^{\adim}}$. By~\ref{al:3.5(1)}(b) a~rectifiable varifold is
    determined by its weight and its approximate tangent planes; therefore $\scale{r \#} C = C$.

    If $V \in \IVar{\vdim}(U)$, then $V_i \in \IVar{\vdim}(\oball{0}{R})$ by~\ref{al:3.5(3)},
    $(\|V_i\| + \|\delta V_i\|)(\oball{0}{R})$ is bounded, and~\ref{al:6.4} with $U = G_1 =
    \oball{0}{R}$ shows that $C \restrict (\oball{0}{R} \times \pgrass{\adim}{\vdim})$ is
    integral. As $R$ is arbitrary, \ref{al:3.5(1)}(c) gives~(c).
\end{proof}

\begin{remark}
    \label{rem:6.5}
    Allard writes the hypothesis as $0 \le \density^{\vdim}(\|V\|,a) < \infty$, which we read as
    the existence and finiteness of the density; the hypothesis that $\|\delta V\|$ is a~Radon
    measure is implicit. He states that he knows of no varifold~$V$ and point~$a$ with
    $\density^{\vdim}(\|V\|,a) < \infty$ and $\density^{\vdim-1}(\|\delta V\|,a) = 0$ for which
    $\VarTan(V,a)$ has more than one element.
\end{remark}

\section{The isoperimetric inequality}
\label{sec:isoperimetric}

\begin{miniremark}
    \label{mr:sec7}
    This section treats~\cite[section~7]{Allard1972}. Theorem~\ref{al:7.1} estimates the
    $\vdim/(\vdim-1)$~power of the weight of a~varifold by the total variation of its first
    variation; \ref{al:7.3} deduces a~Sobolev inequality for functions on a~varifold whose first
    variation lies in a~Lebesgue space with respect to its weight. Corollaries~\ref{al:7.2}
    and~\ref{al:7.4} restate both for submanifolds. In~\ref{al:7.5(1)}--\ref{al:7.5(6)} we
    introduce $V$~subharmonic functions and bound the supremum of the square of a~positive
    $V$~subharmonic function by its mean square. Throughout the section $\vdim \ge 1$ and
    $\mathbf{c}_1 = \besicovitch{\adim} e \vdim^{-1} \unitmeasure{\vdim}^{-1/\vdim}$. The cut-off
    functions are provided by~\ref{new:cutoff}.
\end{miniremark}

\begin{lemma}[new]
    \label{new:cutoff}
    Suppose $K$ is a~bounded subset of~$\R^{\adim}$, $0 < \rho < \infty$ and $1 < \lambda <
    \infty$. Then there is $\phi \in \dspace{\R^{\adim}}{\R}$ with $0 \le \phi \le 1$, $\phi = 1$
    on a~neighbourhood of~$K$, $\spt \phi \subseteq \R^{\adim} \cap \{ x : \dist(x,K) < \rho \}$,
    and $|\grad \phi(x)| \le \lambda / \rho$ for $x \in \R^{\adim}$.
\end{lemma}

\begin{proof}
    If $K = \varnothing$, let $\phi = 0$. Otherwise choose $\tau$ with $0 < 4 \tau < \rho$ and
    $\rho \le \lambda (\rho - 4 \tau)$, let $d(x) = \dist(x,K)$ and
    \begin{displaymath}
        \psi(x) = \inf \{ 1, \sup \{ 0, (\rho - 2\tau - d(x)) / (\rho - 4 \tau) \} \} \quad
        \text{for $x \in \R^{\adim}$} \,.
    \end{displaymath}
    Then $\Lip \psi \le (\rho - 4\tau)^{-1}$, $\psi(x) = 1$ if $d(x) \le 2 \tau$, and $\psi(x) = 0$
    if $d(x) \ge \rho - 2 \tau$. Choose $\zeta \in \dspace{\R^{\adim}}{\R}$ with $\zeta \ge 0$,
    $\spt \zeta \subseteq \cball{0}{\tau}$ and $\textint{}{} \zeta \ud \LM^{\adim} = 1$, and let
    $\phi(x) = \textint{}{} \psi(x - z) \zeta(z) \ud \LM^{\adim} z$ for $x \in \R^{\adim}$. By~\cite[4.1.2]{Federer1969},
    $\phi$ is of class~$\infty$. Since $0 \le \psi \le 1$, $\zeta \ge 0$ and $\textint{}{} \zeta \ud \LM^{\adim} = 1$, we have $0 \le \phi \le 1$. Since $d(x - z) \le d(x) + |z|$
    and $d(x) \le d(x-z) + |z|$, we have $\phi(x) = 1$ if $d(x) < \tau$ and $\phi(x) = 0$ if $d(x)
    > \rho - \tau$; hence $\phi = 1$ on the open set $\{ x : d(x) < \tau \} \supseteq K$, and $\spt
    \phi$ is contained in the compact set $\{ x : d(x) \le \rho - \tau \}$. Finally $|\phi(x) -
    \phi(y)| \le \textint{}{} |\psi(x - z) - \psi(y - z)| \zeta(z) \ud \LM^{\adim} z \le (\rho - 4
    \tau)^{-1} |x - y|$, so $|\grad \phi(x)| = \sup \{ \uD \phi(x)(v) : |v| \le 1 \} \le (\rho - 4
    \tau)^{-1} \le \lambda / \rho$.
\end{proof}

\begin{theorem}[{\cite[7.1]{Allard1972}}]
    \label{al:7.1}
    Suppose $U$ is an open subset of~$\R^{\adim}$, $V \in \Var{\vdim}(U)$, $\|\delta V\|$ is
    a~Radon measure on~$U$, $\phi \in \dspace{U}{\R}$, $\phi \ge 0$, and $A = U \cap \{ x :
    \phi(x) \udensity^{\vdim}(\|V\|,x) \ge 1 \}$. Then $A$ is a~Borel set and
    \begin{displaymath}
        \textint{A}{} \phi \ud \|V\| \le \mathbf{c}_1 \Big( \textint{}{} \phi \ud \|V\| \Big)^{1/\vdim}
        \Big( \textint{}{} \phi \ud \|\delta V\| + \textint{}{} |S(\grad \phi(x))| \ud V(x,S) \Big)
        \,.
    \end{displaymath}
\end{theorem}

\begin{proof}
    For $0 < r < \infty$ the function $[x \mapsto \|V\| \cball{x}{r}]$ is upper semicontinuous on~$U$
    wherever $\cball{x}{r} \subseteq U$, and $[r \mapsto \|V\| \cball{x}{r}]$ is continuous from the
    right. Hence $\sup \{ r^{-\vdim} \|V\| \cball{x}{r} : 0 < r < 1/j \}$ equals the supremum over
    rational~$r$, and $\udensity^{\vdim}(\|V\|,\cdot)$ is the infimum over~$j$ of these Borel
    functions; thus $A$ is a~Borel set.

    Define $V_\phi \in \Var{\vdim}(\R^{\adim})$ by $V_\phi(B) = \textint{B \cap \Gk{U}}{} \phi(x)
    \ud V(x,S)$ for $B \subseteq \Gk{\R^{\adim}}$; this is a~Radon measure because $\spt \phi$ is
    a~compact subset of~$U$, and $\|V_\phi\| = \|V\| \restrict \phi$. Suppose $g \in
    \VF(\R^{\adim})$. Then $\phi g | U \in \VF(U)$ and $\uD(\phi g)(x) = \phi(x) \uD g(x) + \uD
    \phi(x) \cdot g(x)$ for $x \in U$, so~\ref{al:4.2:lem}, \ref{al:2.3(1)}\eqref{eq:2.3:4}
    and~\eqref{eq:2.3:1} give
    \begin{displaymath}
        \delta V_\phi(g) = \textint{}{} \phi(x) \uD g(x) \bullet S \ud V(x,S) = \delta V(\phi g|U) -
        \textint{}{} g(x) \bullet S(\grad \phi(x)) \ud V(x,S) \,.
    \end{displaymath}
    By~\ref{al:4.3}\eqref{eq:4.3:2}, $|\delta V(\phi g | U)| \le \sup |g| \textint{}{} \phi \ud
    \|\delta V\|$. Therefore
    \begin{displaymath}
        \|\delta V_\phi\|(\R^{\adim}) \le \textint{}{} \phi \ud \|\delta V\| + \textint{}{}
        |S(\grad \phi(x))| \ud V(x,S) \,,
    \end{displaymath}
    and $\|\delta V_\phi\|$ is a~Radon measure on~$\R^{\adim}$ by~\ref{al:4.2:lem}. If $\|V_\phi\|(\R^{\adim})
    = 0$, the assertion is trivial; assume $\|V_\phi\|(\R^{\adim}) > 0$.

    For $a \in U$ with $\phi(a) > 0$ we have $|\|V_\phi\| \cball{a}{r} - \phi(a) \|V\|
    \cball{a}{r}| \le \sup \{ |\phi(x) - \phi(a)| : x \in \cball{a}{r} \} \|V\| \cball{a}{r}$ for
    small~$r$, hence $\udensity^{\vdim}(\|V_\phi\|,a) = \phi(a) \udensity^{\vdim}(\|V\|,a)$. Thus
    $A = \R^{\adim} \cap \{ a : \udensity^{\vdim}(\|V_\phi\|,a) \ge 1 \}$.

    Let $1 < \lambda' < \lambda < \infty$ and $s = (\lambda \|V_\phi\|(\R^{\adim}) /
    \unitmeasure{\vdim})^{1/\vdim}$. Suppose $a \in A$. There is $r$ with $0 < r < s$ and
    $r^{-\vdim} \|V_\phi\| \cball{a}{r} > \unitmeasure{\vdim} \lambda' / \lambda$; in particular $a
    \in \spt \|V_\phi\|$. Since $s^{-\vdim} \|V_\phi\| \cball{a}{s} \le \unitmeasure{\vdim} /
    \lambda$, \ref{al:5.1(1)} applied to~$V_\phi$ yields
    \begin{displaymath}
        \exp \textint{r}{s} \delta V_\phi(\kappa(V_\phi,a,t)) \ud \LM^1 t \ge \frac{r^{-\vdim}
        \|V_\phi\| \cball{a}{r}}{s^{-\vdim} \|V_\phi\| \cball{a}{s}} > \lambda' \,.
    \end{displaymath}
    By~\ref{al:4.3}\eqref{eq:4.3:2} and $|\kappa(V_\phi,a,t)| \le (\|V_\phi\|
    \cball{a}{t})^{-1}$, we have $\delta V_\phi(\kappa(V_\phi,a,t)) \le \|\delta V_\phi\|
    \cball{a}{t} / \|V_\phi\| \cball{a}{t}$. If this quotient were at most $(\log \lambda') / s$ for
    every $t$ with $r < t < s$, the integral would be at most $(s - r) s^{-1} \log \lambda' <
    \log \lambda'$. Hence there is $t(a)$ with $0 < t(a) < s$ and
    \begin{displaymath}
        \|V_\phi\| \cball{a}{t(a)} < (s / \log \lambda') \|\delta V_\phi\| \cball{a}{t(a)} \,.
    \end{displaymath}
    By~\ref{al:2.7(1)} there are disjointed subfamilies $\mathcal{C}_1, \ldots,
    \mathcal{C}_{\besicovitch{\adim}}$ of $\{ \cball{a}{t(a)} : a \in A \}$ covering~$A$. Each
    $\mathcal{C}_l$ is countable, since its members are disjoint and have positive
    $\|V_\phi\|$~measure, and $\|V_\phi\|(\R^{\adim}) < \infty$. Hence
    \begin{displaymath}
        \|V_\phi\|(A) \le \textsum{l=1}{\besicovitch{\adim}} \textsum{C \in \mathcal{C}_l}{}
        \|V_\phi\|(C) \le \frac{s}{\log \lambda'} \textsum{l=1}{\besicovitch{\adim}} \textsum{C \in
        \mathcal{C}_l}{} \|\delta V_\phi\|(C) \le \besicovitch{\adim} \frac{s}{\log \lambda'}
        \|\delta V_\phi\|(\R^{\adim}) \,.
    \end{displaymath}
    Letting $\lambda' \to \lambda$ and taking $\lambda = e^{\vdim}$, so that $s / \log \lambda = e
    \vdim^{-1} \unitmeasure{\vdim}^{-1/\vdim} \|V_\phi\|(\R^{\adim})^{1/\vdim}$, we obtain the
    assertion, because $\|V_\phi\|(A) = \textint{A}{} \phi \ud \|V\|$ and $\|V_\phi\|(\R^{\adim}) =
    \textint{}{} \phi \ud \|V\|$.
\end{proof}

\begin{remark}
    \label{rem:7.1}
    Allard states~\ref{al:7.1} with the density $\density^{\vdim}(\|V_\phi\|,x) = \phi(x)
    \density^{\vdim}(\|V\|,x)$ in place of the upper density, and without the hypothesis that
    $\|\delta V\|$ is a~Radon measure; the proof above applies to the larger set~$A$. The
    choice $\lambda = e^{\vdim}$ minimises $\lambda^{1/\vdim} / \log \lambda$.
\end{remark}

\begin{corollary}[{\cite[7.2]{Allard1972}}]
    \label{al:7.2}
    Suppose $m \ge 1$, $M$ is a~nonempty smooth $m$~dimensional submanifold of~$\R^{\adim}$,
    $\HM^m(M) < \infty$, $F = \Clos M \without M$, with $\dist(x,\varnothing) = \infty$, and
    \begin{displaymath}
        P = \liminf_{r \to 0+} \HM^{m-1}(M \cap \{ x : \dist(x,F) = r \}) \,.
    \end{displaymath}
    Then
    \begin{displaymath}
        \HM^m(M)^{(m-1)/m} \le \mathbf{c}_1 \Big( P + \textint{M}{} |h(M;x)| \ud \HM^m x \Big) \,,
    \end{displaymath}
    where $\mathbf{c}_1$ is computed with $\vdim = m$.
\end{corollary}

\begin{proof}
    Each $a \in M$ has a~chart $(W,\psi,\phi)$ as in~\ref{al:2.5}; the set $\im \phi = \{ x :
    \phi(\psi(x)) = x \}$ is closed, and $W \cap \Clos M \subseteq \Clos(W \cap M) \cap W \subseteq
    W \cap \im \phi = W \cap M$. Hence $F$ is closed, $U = \R^{\adim} \without F$ is open, $M
    \subseteq U$, and $M = U \cap \Clos M$, so that the inclusion map $i : M \to U$ is proper. Let
    $d(x) = \dist(x,F)$; then $d$ is continuous and $d(x) > 0$ for $x \in M$. For $0 < r < \infty$
    let $M_r = M \cap \{ x : d(x) > r \}$ and $B_r = M \cap \{ x : d(x) = r \}$.

    Suppose $0 < r < \infty$ and $\HM^{m-1}(B_r) < \infty$. If $g \in \VF(M)$ and $\spt g \subseteq M
    \cap \{ x : d(x) \ne r \}$, then $g = g_1 + g_2$ with $g_1, g_2 \in \VF(M)$, $\spt g_1
    \subseteq M_r$ and $\spt g_2 \subseteq \{ x : d(x) < r \}$; hence $\delta \var{}(M_r)(g_2) =
    0$, and $\delta \var{}(M_r)(g_1) = \delta \var{}(M)(g_1) = 0$ by~\ref{new:divergence},
    applied after a~partition of unity~\cite[3.1.13]{Federer1969} subordinate to finitely many
    charts covering $\spt g_1$. Thus $\|\delta \var{}(M_r)\|(M \cap \{ x : d(x) \ne r \}) = 0$. If
    $b \in M$ and $d(b) \ne r$, then either $M_r$ or $M \without M_r$ does not meet a~neighbourhood
    of~$b$, so $b$ belongs to the set $Q \cup R$ of~\ref{al:4.7} with $E = M_r$; hence $K
    \without (Q \cup R) \subseteq B_r$ for compact $K \subseteq M$, and~\ref{al:4.7} gives
    $\|\delta \var{}(M_r)\| = \HM^{m-1} \restrict \{ b : n(M;M_r,b) \ne 0 \}$. Combining, $\|\delta
    \var{}(M_r)\| \le \HM^{m-1} \restrict B_r$. Now~\ref{al:4.4}, with $|\Tan(M,g)| \le |g|$ and
    $h(M;x,\project{\Tan(M,x)}) = h(M;x)$, yields
    \begin{displaymath}
        \|\delta(i_\# \var{}(M_r))\|(U) \le \HM^{m-1}(B_r) + \textint{M}{} |h(M;x)| \ud \HM^m x \,.
    \end{displaymath}
    Let $W_r \in \Var{m}(\R^{\adim})$ be $i_\# \var{}(M_r)$ extended by zero. Since $\spt
    \|W_r\| \subseteq \Clos M_r \subseteq \Clos M \cap \{ x : d(x) \ge r \} \subseteq M$, for $g \in
    \VF(\R^{\adim})$ we may choose $\zeta \in \dspace{U}{\R}$ with $0 \le \zeta \le 1$ and $\zeta =
    1$ near the compact set $\spt \|W_r\| \cap \spt g$, by~\ref{new:cutoff}; then $\delta W_r(g) =
    \delta(i_\# \var{}(M_r))(\zeta g|U)$, hence $\|\delta W_r\|(\R^{\adim}) \le \HM^{m-1}(B_r) +
    \textint{M}{} |h(M;x)| \ud \HM^m x$.

    Fix $\varphi \in \dspace{\R^{\adim}}{\R}$ with $0 \le \varphi \le 1$, $\varphi = 1$ on
    $\cball{0}{1}$ and $|\grad \varphi| \le 2$, by~\ref{new:cutoff}, and let $\varphi_j(x) =
    \varphi(x/j)$ for $j = 1, 2, \ldots$. For $x \in M_r$, $\|W_r\|$ agrees with $\HM^m \restrict M$
    near~$x$, and $\density^m(\HM^m \restrict M, x) = 1$ for $\HM^m$ almost all $x \in M$
    by~\cite[3.2.19]{Federer1969}, because $M$ is a~countable union of $(\HM^m,m)$ rectifiable
    compact sets $\phi[C]$ with $C$ compact. Hence \ref{al:7.1} applied to $W_r$ and~$\varphi_j$
    gives
    \begin{displaymath}
        \HM^m(M_r \cap \cball{0}{j}) \le \mathbf{c}_1 \HM^m(M)^{1/m} \Big( \HM^{m-1}(B_r) +
        \textint{M}{} |h(M;x)| \ud \HM^m x + 2 j^{-1} \HM^m(M) \Big) \,.
    \end{displaymath}
    This also holds if $\HM^{m-1}(B_r) = \infty$. Letting $j \to \infty$ and then $r \to 0+$ along
    a~sequence with $\HM^{m-1}(B_r) \to P$, and noting $M = \tbcup_{r > 0} M_r$, we obtain
    $\HM^m(M) \le \mathbf{c}_1 \HM^m(M)^{1/m} (P + \textint{M}{} |h(M;x)| \ud \HM^m x)$. Since $M$ is
    nonempty, $\HM^m(M) > 0$ by the area formula~\cite[3.2.3, 3.2.5]{Federer1969} applied to
    a~chart, and the assertion follows.
\end{proof}

\begin{remark}
    \label{rem:7.2}
    The hypothesis $M \ne \varnothing$ is needed for $m = 1$. The proof uses~\ref{al:4.7}, whose
    proof rests on~\cite[4.5.6, 4.5.12]{Federer1969}, beyond the formalised range.
\end{remark}

\begin{theorem}[{\cite[7.3]{Allard1972}}]
    \label{al:7.3}
    Suppose $U$ is an open subset of~$\R^{\adim}$, $V \in \Var{\vdim}(U)$, $0 \le \alpha <
    \infty$, $f \in \dspace{U}{\R}$, $\beta = \vdim/(\vdim-1)$ if $\vdim > 1$,
    \begin{gather*}
        \delta V(g) \le \alpha \sup \{ |g(x)| : x \in U \} \quad \text{for $g \in \VF(U)$, in case
        $\vdim = 1$} \,, \\
        \delta V(g) \le \alpha \Big( \textint{}{} |g|^\beta \ud \|V\| \Big)^{1/\beta} \quad
        \text{for $g \in \VF(U)$, in case $\vdim > 1$} \,,
    \end{gather*}
    and $\udensity^{\vdim}(\|V\|,x) \ge 1$ for $\|V\|$ almost all $x \in U$. Then
    \begin{gather*}
        (1 - \mathbf{c}_1 \alpha) \sup \{ |f(x)| : x \in \spt \|V\| \} \le \mathbf{c}_1 \textint{}{}
        |S(\grad f(x))| \ud V(x,S) \quad \text{in case $\vdim = 1$} \,, \\
        (1 - \mathbf{c}_1 \alpha) \Big( \textint{}{} |f|^\beta \ud \|V\| \Big)^{1/\beta} \le
        \mathbf{c}_1 \textint{}{} |S(\grad f(x))| \ud V(x,S) \quad \text{in case $\vdim > 1$} \,.
    \end{gather*}
\end{theorem}

\begin{proof}
    We may assume $1 - \mathbf{c}_1 \alpha > 0$. The hypothesis implies $\|\delta V\|(G) \le \alpha
    \|V\|(G)^{1/\beta}$ for open $G \subseteq U$ if $\vdim > 1$, and $\|\delta V\|(U) \le \alpha$ if
    $\vdim = 1$; in particular $\|\delta V\|$ is a~Radon measure on~$U$.

    We first assume $f \ge 0$. Let $s = \sup \{ f(x) : x \in U \}$, and for $0 \le r < \infty$ let $V_r
    = V \restrict \{ (x,S) : f(x) > r \}$. Since $f$ is smooth on~$U$, \ref{al:4.10(2)} with $a = 0$,
    $b = s$ and $B = U$ gives that $\|\delta V_r\|$ is a~Radon measure for $\LM^1$ almost all $r$
    with $0 < r < s$, and
    \begin{equation}
        \label{eq:7.3:1}
        \textint{0}{s} \|\delta V_r\|(U) \ud \LM^1 r \le \textint{}{} |S(\grad f(x))| \ud V(x,S) +
        \textint{0}{s} \|\delta V\|(U \cap \{ x : f(x) > r \}) \ud \LM^1 r \,.
    \end{equation}
    By~\ref{new:cutoff} there is $\phi \in \dspace{U}{\R}$ with $0 \le \phi \le 1$ and $\phi = 1$
    on a~neighbourhood of $\spt f$. Suppose $0 < r < s$ and $\|\delta V_r\|$ is a~Radon measure.
    Since $\{ x : f(x) > r \}$ is open, $\udensity^{\vdim}(\|V_r\|,x) = \udensity^{\vdim}(\|V\|,x)
    \ge 1$ for $\|V_r\|$ almost all~$x$, and $\phi = 1$ on $\spt \|V_r\| \subseteq \spt f$; also
    $\grad \phi = 0$ near $\spt \|V_r\|$. Hence~\ref{al:7.1} applied to $V_r$ and~$\phi$ gives
    \begin{equation}
        \label{eq:7.3:2}
        \|V_r\|(U) \le \mathbf{c}_1 \|V_r\|(U)^{1/\vdim} \|\delta V_r\|(U) \,.
    \end{equation}

    Suppose $\vdim > 1$. For $0 \le r < \infty$ let $f_r = \inf \{ f, r \}$ and $\theta(r) =
    (\textint{}{} f_r^\beta \ud \|V\|)^{1/\beta}$. If $0 \le r$ and $0 < h$, then $0 \le f_{r+h} -
    f_r \le h \CF{\{ x : f(x) > r \}}$, so Minkowski's inequality~\cite[2.4.15]{Federer1969} gives
    $0 \le \theta(r+h) - \theta(r) \le h \|V_r\|(U)^{1/\beta}$. Thus $\theta$ is Lipschitzian, hence
    absolutely continuous, $\theta'(r) \le \|V_r\|(U)^{1/\beta}$ for $\LM^1$ almost all~$r$, and
    $\theta(s) = (\textint{}{} f^\beta \ud \|V\|)^{1/\beta}$. By~\cite[2.9.20]{Federer1969},
    \begin{displaymath}
        \Big( \textint{}{} f^\beta \ud \|V\| \Big)^{1/\beta} = \textint{0}{s} \theta' \ud \LM^1 \le
        \textint{0}{s} \|V_r\|(U)^{1/\beta} \ud \LM^1 r \,.
    \end{displaymath}
    Since $1 - 1/\vdim = 1/\beta$, \eqref{eq:7.3:2} gives $\|V_r\|(U)^{1/\beta} \le \mathbf{c}_1
    \|\delta V_r\|(U)$ for $\LM^1$ almost all $0 < r < s$, and $\|\delta V\|(U \cap \{ x : f(x) > r
    \}) \le \alpha \|V_r\|(U)^{1/\beta}$. Combining with~\eqref{eq:7.3:1},
    \begin{displaymath}
        \textint{0}{s} \|V_r\|(U)^{1/\beta} \ud \LM^1 r \le \mathbf{c}_1 \textint{}{} |S(\grad
        f(x))| \ud V(x,S) + \mathbf{c}_1 \alpha \textint{0}{s} \|V_r\|(U)^{1/\beta} \ud \LM^1 r \,,
    \end{displaymath}
    where the integrals are finite, because $\|V_r\|(U) \le \|V\|(\spt f) < \infty$. This proves the
    assertion for $f \ge 0$.

    Suppose $\vdim = 1$, and let $s' = \sup \{ f(x) : x \in \spt \|V\| \}$. We may assume $s' > 0$.
    For $0 < r < s'$ the set $\spt \|V\| \cap \{ x : f(x) > r \}$ is a~nonempty relatively open
    subset of $\spt \|V\|$, so $0 < \|V_r\|(U) < \infty$, and~\eqref{eq:7.3:2} gives $1 \le
    \mathbf{c}_1 \|\delta V_r\|(U)$ for $\LM^1$ almost all such~$r$. Since $\|\delta V\|(U) \le
    \alpha$ and $s' \le s$, \eqref{eq:7.3:1} yields $s' \le \mathbf{c}_1 \textint{0}{s} \|\delta
    V_r\|(U) \ud \LM^1 r \le \mathbf{c}_1 \textint{}{} |S(\grad f(x))| \ud V(x,S) + \mathbf{c}_1
    \alpha s$. Since $\|\delta V_r\| = 0$ for $r \ge s'$, \eqref{eq:7.3:1} also holds with $s'$ in
    place of~$s$; hence $(1 - \mathbf{c}_1 \alpha) s' \le \mathbf{c}_1 \textint{}{} |S(\grad f(x))|
    \ud V(x,S)$.

    For general~$f$ and $0 < \epsilon < \infty$ let $f_\epsilon = (f^2 + \epsilon^2)^{1/2} -
    \epsilon$. Then $f_\epsilon \in \dspace{U}{\R}$, $f_\epsilon \ge 0$, $\grad f_\epsilon = f (f^2 +
    \epsilon^2)^{-1/2} \grad f$, so $|S(\grad f_\epsilon(x))| \le |S(\grad f(x))|$, and $f_\epsilon$
    increases to~$|f|$ as $\epsilon \to 0+$, with $f_\epsilon \ge |f| - \epsilon$. The assertion
    for~$f$ follows from that for the $f_\epsilon$ by Lebesgue's increasing convergence
    theorem~\cite[2.4.7]{Federer1969} in case $\vdim > 1$, and directly in case $\vdim = 1$.
\end{proof}

\begin{remark}
    \label{rem:7.3}
    Allard assumes $\density^{\vdim}(\|V\|,x) \ge 1$ and applies~\ref{al:7.1} with approximations
    to the characteristic function of $\{ x : f(x) > r \}$; a~single cut-off equal to~$1$ near $\spt
    f$ suffices. He leaves the case $\vdim = 1$ and the reduction to $f \ge 0$ to the reader.
    He remarks that the introduction of the function~$\theta$ is an idea of Federer;
    see~\cite[4.5.9(18)]{Federer1969}, beyond the formalised range, which is cited for credit only.
\end{remark}

\begin{corollary}[{\cite[7.4]{Allard1972}}]
    \label{al:7.4}
    Suppose $m \ge 1$, $\beta = m/(m-1)$ if $m > 1$, $M$ is a~smooth $m$~dimensional submanifold
    of~$\R^{\adim}$, and $f : M \to \R$ is smooth with compact support. Then, with $\mathbf{c}_1$
    computed with $\vdim = m$,
    \begin{gather*}
        \Big( 1 - \mathbf{c}_1 \textint{\spt f}{} |h(M;x)| \ud \HM^1 x \Big) \sup \{ |f(x)| : x \in
        M \} \le \mathbf{c}_1 \textint{M}{} |\grad f(x)| \ud \HM^1 x \quad \text{in case $m = 1$} \,,
        \\
        \Big( 1 - \mathbf{c}_1 \Big( \textint{\spt f}{} |h(M;x)|^m \ud \HM^m x \Big)^{1/m} \Big)
        \Big( \textint{M}{} |f|^\beta \ud \HM^m \Big)^{1/\beta} \le \mathbf{c}_1 \textint{M}{} |\grad
        f(x)| \ud \HM^m x \quad \text{in case $m > 1$} \,.
    \end{gather*}
\end{corollary}

\begin{proof}
    Let $F$, $U$ and $i : M \to U$ be as in the proof of~\ref{al:7.2}. For compact $K \subseteq U$
    the set $K \cap M$ is compact and covered by finitely many chart images, so $\HM^m(K \cap M) <
    \infty$; hence $\var{}(M) \in \Var{m}(M)$, and $V = i_\# \var{}(M) \in \Var{m}(U)$. By~\ref{al:4.4},
    \ref{new:divergence} with a~partition of unity~\cite[3.1.13]{Federer1969}, and $h(M;x) \in
    \Nor(M,x)$,
    \begin{displaymath}
        \delta V(g) = - \textint{M}{} g(x) \bullet h(M;x) \ud \HM^m x \quad \text{for $g \in \VF(U)$}
        \,.
    \end{displaymath}
    By~\ref{al:2.5}, each $a \in \spt f$ has an open neighbourhood $W_a \subseteq U$ and a~map
    $G_a$ of class~$\infty$ on~$W_a$ with $G_a | W_a \cap M = f | W_a \cap M$. Choose finitely many
    $W_{a_1}, \ldots, W_{a_l}$ covering $\spt f$ and $\zeta_j \in \dspace{W_{a_j}}{\R}$ with
    $\textsum{j=1}{l} \zeta_j = 1$ on a~neighbourhood of $\spt f$~\cite[3.1.13]{Federer1969}, and
    let $\tilde f = \textsum{j=1}{l} \zeta_j G_{a_j} \in \dspace{U}{\R}$. Then $\tilde f | M = f$,
    and $\project{\Tan(M,x)}(\grad \tilde f(x)) = \grad f(x)$ for $x \in M$, by~\ref{al:2.5}(b).
    Since only $S = \project{\Tan(M,x)}$ occurs, $\textint{}{} |S(\grad \tilde f(x))| \ud V(x,S) =
    \textint{M}{} |\grad f| \ud \HM^m$. Also $\density^m(\|V\|,x) = 1$ for $\|V\|$ almost all~$x$,
    as in the proof of~\ref{al:7.2}, and $\spt \|V\| = M$, because nonempty relatively open subsets
    of~$M$ have positive $\HM^m$~measure.

    Let $0 < \rho < \infty$ be so small that $\{ x : \dist(x, \spt f) \le \rho \}$ is a~compact
    subset of~$U$, and let $G = \R^{\adim} \cap \{ x : \dist(x, \spt f) < \rho \}$. In the
    construction of~$\tilde f$ we take $W_{a_j} \subseteq G$, so that $\tilde f | G \in
    \dspace{G}{\R}$. Since $M \cap \Clos G$ is compact and $h(M;\cdot)$ is continuous, $\alpha_\rho =
    (\textint{G \cap M}{} |h(M;x)|^m \ud \HM^m x)^{1/m} < \infty$. For $g \in \VF(G)$, H\"older's
    inequality~\cite[2.4.14]{Federer1969} gives
    \begin{displaymath}
        \delta V(g) \le \alpha_\rho \Big( \textint{}{} |g|^\beta \ud \|V\| \Big)^{1/\beta}
    \end{displaymath}
    if $m > 1$, and $\delta V(g) \le \alpha_\rho \sup |g|$ if $m = 1$. Moreover $\spt \|V \restrict (G
    \times \pgrass{\adim}{m})\| = G \cap M$, and $\tilde f = f$ on~$M$. Applying~\ref{al:7.3} to $V
    \restrict (G \times \pgrass{\adim}{m}) \in \Var{m}(G)$ and $\tilde f | G$, we obtain the
    assertion with $\alpha_\rho$ in place of $(\textint{\spt f}{} |h(M;x)|^m \ud \HM^m x)^{1/m}$.
    As $\rho \to 0+$, the sets $G \cap M$ decrease to $\spt f$, and Lebesgue's dominated
    convergence theorem~\cite[2.4.9]{Federer1969} gives the assertion.
\end{proof}

\begin{remark}
    \label{rem:7.4}
    Allard's proof reads ``Apply 4.4 and 4.7 with 7.3''. Since $f$ has compact support in~$M$,
    \ref{al:4.7} is not needed. The extension~$\tilde f$ supplies the passage from functions
    on~$M$ to functions on an~open set required by~\ref{al:7.3}.
\end{remark}

\begin{miniremark}
    \label{mr:7.5}
    In~\ref{al:7.5(1)}--\ref{al:7.5(6)}, which treat~\cite[7.5]{Allard1972}, we suppose $\vdim \ge
    1$, $V \in \Var{\vdim}(\R^{\adim})$, and $U$ is an~open subset of~$\R^{\adim}$. For $0 < \delta <
    \infty$ let $U(\delta) = \tbcup \{ \oball{x}{\delta} : x \in U \}$; then $U(\delta)(\delta') =
    U(\delta + \delta')$. A~member of $\VF(U)$ or $\dspace{U}{\R}$ is regarded as a~member of
    $\VF(\R^{\adim})$ or $\dspace{\R^{\adim}}{\R}$ by extension by zero. Allard remarks that the
    method bears formal resemblance to, and was suggested by, Moser's proof of the regularity
    theorem of De Giorgi and Nash.
\end{miniremark}

\begin{definition}[{\cite[7.5(1)]{Allard1972}}]
    \label{al:7.5(1)}
    A~smooth function $f : U \to \R$ is \emph{$V$~subharmonic in~$U$} if and only if
    \begin{displaymath}
        \textint{}{} S(\grad f(x)) \bullet \grad \phi(x) \ud V(x,S) \le 0 \quad \text{whenever
        $\phi \in \dspace{U}{\R}$ and $\phi \ge 0$} \,;
    \end{displaymath}
    it is \emph{$V$~harmonic in~$U$} if and only if the integral vanishes for every $\phi \in
    \dspace{U}{\R}$. Here the integrand is taken to be~$0$ for $x \notin U$.
\end{definition}

\begin{lemma}[{\cite[7.5(2)]{Allard1972}}]
    \label{al:7.5(2)}
    Suppose $f : U \to \R$ is smooth. Then
    \begin{equation}
        \label{eq:7.5(2):1}
        \begin{multlined}
        \textint{}{} S(\grad f(x)) \bullet \grad \phi(x) \ud V(x,S) = \delta V(\phi \grad f) \\ -
        \textint{}{} \phi(x) \uD(\grad f)(x) \bullet S \ud V(x,S) \quad \text{for $\phi \in
        \dspace{U}{\R}$} \,.
        \end{multlined}
    \end{equation}
    Moreover:
    \begin{enumerate}
        \item If $\|\delta V\| \restrict U$ is a~Radon measure and $\|\delta V\|_{\mathrm{sing}}(U) =
            0$, then $f$ is $V$~subharmonic in~$U$ if and only if
            \begin{gather*}
                \textint{}{} \phi(x) \big( \grad f(x) \bullet h(V;x) + \uD(\grad f)(x) \bullet S
                \big) \ud V(x,S) \ge 0 \\ \text{whenever $\phi \in \dspace{U}{\R}$ and $\phi \ge 0$}
                \,,
            \end{gather*}
            and $f$ is $V$~harmonic in~$U$ if and only if this integral vanishes for every $\phi
            \in \dspace{U}{\R}$; in particular, $f$ is $V$~subharmonic in~$U$ if $- \grad f(x)
            \bullet h(V;x) \le \uD(\grad f)(x) \bullet S$ for $V$ almost all $(x,S) \in U \times
            \pgrass{\adim}{\vdim}$.
        \item If $V$ is stationary in~$U$, that is $\delta V(g) = 0$ for $g \in \VF(\R^{\adim})$ with
            $\spt g \subseteq U$, then every smooth convex function on~$U$ is $V$~subharmonic in~$U$,
            and every function $[x \mapsto u \bullet x + c]$ with $u \in \R^{\adim}$, $c \in \R$ is
            $V$~harmonic in~$U$. Conversely, if each such function is $V$~harmonic in~$U$, then $V$
            is stationary in~$U$.
        \item If $0 < R < \infty$, $U = \oball{0}{R}$, $\|\delta V\| \restrict U$ is a~Radon measure,
            $\|\delta V\|_{\mathrm{sing}}(U) = 0$, $K : \R^{\adim} \to \R$ is smooth and convex, $0 <
            \mu < \infty$, $0 < t < \infty$, $|h(V;x)| \le \mu$ for $\|V\|$ almost all $x \in U$,
            $\mu (\sup \{ |\grad K(x)| : |x| < R \} + 2 t R) \le 2 \vdim t$, and $f(x) = K(x) + t
            |x|^2$ for $|x| < R$, then $f$ is $V$~subharmonic in~$U$.
    \end{enumerate}
\end{lemma}

\begin{proof}
    For $\phi \in \dspace{U}{\R}$ we have $\phi \grad f \in \VF(U)$ and $\uD(\phi \grad f)(x) =
    \phi(x) \uD(\grad f)(x) + \uD \phi(x) \cdot \grad f(x)$; by~\ref{al:2.3(1)}\eqref{eq:2.3:4}
    and~\eqref{eq:2.3:1}, $(\uD \phi(x) \cdot \grad f(x)) \bullet S = \uD \phi(x)(S(\grad f(x))) =
    S(\grad f(x)) \bullet \grad \phi(x)$. Integrating with respect to~$V$ and
    using~\ref{al:4.2:lem} gives~\eqref{eq:7.5(2):1}.

    Under the hypotheses of~(a), \ref{al:4.3} gives $\delta V(g) = - \textint{}{} g(x) \bullet
    h(V;x) \ud \|V\| x$ for $g \in \VF(\R^{\adim})$ with $\spt g \subseteq U$; with $g = \phi
    \grad f$, \eqref{eq:7.5(2):1} becomes
    \begin{displaymath}
        \textint{}{} S(\grad f(x)) \bullet \grad \phi(x) \ud V(x,S) = - \textint{}{} \phi(x) \big(
        \grad f(x) \bullet h(V;x) + \uD(\grad f)(x) \bullet S \big) \ud V(x,S) \,,
    \end{displaymath}
    which implies~(a).

    If $V$ is stationary in~$U$, \eqref{eq:7.5(2):1} reduces to $\textint{}{} S(\grad f(x)) \bullet
    \grad \phi(x) \ud V(x,S) = - \textint{}{} \phi(x) \uD(\grad f)(x) \bullet S \ud V(x,S)$. If $f$
    is convex, then $\uD(\grad f)(x)$ is self adjoint with $v \bullet \uD(\grad f)(x)(v) =
    \frac{\ud^2}{\ud t^2} f(x + t v)|_{t=0} \ge 0$; choosing an orthonormal basis $e_1, \ldots,
    e_{\vdim}$ of $\im S$, \ref{al:2.3(1)}\eqref{eq:2.3:5} gives $\uD(\grad f)(x) \bullet S =
    \trace(S \circ \uD(\grad f)(x) \circ S) = \textsum{i=1}{\vdim} e_i \bullet \uD(\grad f)(x)(e_i)
    \ge 0$. If $f(x) = u \bullet x + c$, then $\uD(\grad f) = 0$. Conversely, if $[x \mapsto u
    \bullet x]$ is $V$~harmonic in~$U$ for each $u \in \R^{\adim}$, \eqref{eq:7.5(2):1} gives
    $\delta V(\phi u) = 0$ for $\phi \in \dspace{U}{\R}$; every $g \in \VF(\R^{\adim})$ with $\spt g
    \subseteq U$ is a~sum $\textsum{i=1}{\adim} \phi_i e_i$ with $\phi_i \in \dspace{U}{\R}$ and
    $e_1, \ldots, e_{\adim}$ the standard basis, so $V$ is stationary in~$U$. This proves~(b).

    Under the hypotheses of~(c), $\grad f(x) = \grad K(x) + 2 t x$ and $\uD(\grad f)(x) = \uD(\grad
    K)(x) + 2 t \id{\R^{\adim}}$, so $\uD(\grad f)(x) \bullet S \ge 2 t \vdim$
    by~\ref{al:2.3(1)}\eqref{eq:2.3:3} and the computation in~(b), while $- \grad f(x) \bullet
    h(V;x) \le \mu (|\grad K(x)| + 2 t R) \le 2 \vdim t$ for $\|V\|$ almost all $x \in U$. Now
    apply~(a).
\end{proof}

\begin{remark}
    \label{rem:7.5(2)}
    Allard asserts that $f$ is $V$~subharmonic if and only if $\grad f(x) \bullet h(V;x) \le
    \uD(\grad f)(x) \bullet S$ for $V$ almost all $(x,S)$, and similarly for $V$~harmonic functions.
    With $h(V;\cdot)$ as in~\ref{al:4.3}, the sign of the left side is misprinted: for the unit
    sphere $M$ in~$\R^3$, $V = \var{}(M)$ and $f(x) = |x|^2/2$, which is $V$~harmonic, one has
    $\grad f(x) \bullet h(V;x) = -2$ and $\uD(\grad f)(x) \bullet S = 2$. Moreover the pointwise
    condition is only sufficient, since the test functions depend on~$x$ alone; the integrated
    condition in~(a) is necessary and sufficient. The hypothesis $\|\delta V\|_{\mathrm{sing}}(U) =
    0$ is needed in~(c), and Allard writes $\oball{0}{r}$ for $\oball{0}{R}$ there.
\end{remark}

\begin{lemma}[{\cite[7.5(3)(4)]{Allard1972}}]
    \label{al:7.5(3)}
    Suppose $f : U \to \R$ is smooth, $I$ is an~open interval with $f[U] \subseteq I$, and $C : I \to
    \R$ is smooth and convex. If $f$ is $V$~subharmonic in~$U$ and $C$ is nondecreasing, or if $f$
    is $V$~harmonic in~$U$, then $C \circ f$ is $V$~subharmonic in~$U$.
\end{lemma}

\begin{proof}
    Let $\phi \in \dspace{U}{\R}$, $\phi \ge 0$. Then $(C' \circ f) \phi \in \dspace{U}{\R}$, and
    $\grad (C \circ f) = (C' \circ f) \grad f$, $\grad((C' \circ f) \phi) = (C' \circ f) \grad \phi
    + \phi (C'' \circ f) \grad f$. Hence, by~\ref{al:2.3(1)}\eqref{eq:2.3:1},\eqref{eq:2.3:2},
    \begin{displaymath}
        S(\grad (C \circ f)(x)) \bullet \grad \phi(x) = S(\grad f(x)) \bullet \grad((C' \circ f)
        \phi)(x) - (C'' \circ f)(x) \phi(x) |S(\grad f(x))|^2
    \end{displaymath}
    for $(x,S) \in U \times \pgrass{\adim}{\vdim}$. The integral of the first term with respect
    to~$V$ is nonpositive if $f$ is $V$~subharmonic and $C' \ge 0$, and vanishes if $f$ is
    $V$~harmonic; the second term is nonpositive because $C'' \ge 0$.
\end{proof}

\begin{remark}
    \label{rem:7.5(3)}
    Allard requires $C$ to be defined on~$\R$; in~\ref{al:7.5(6)} the lemma is applied to $C(t) =
    t^{\beta^\nu}$ on $\{ t : t > 0 \}$, which is not smooth on~$\R$ when $\beta^\nu$ is not an
    integer. Hence the interval~$I$.
\end{remark}

\begin{lemma}[{\cite[7.5(5)]{Allard1972}}]
    \label{al:7.5(5)}
    If $0 < \delta < \infty$, $f$ is $V$~subharmonic in $U(\delta)$, and $f \ge 0$, then
    \begin{displaymath}
        \delta^2 \textint{U \times \pgrass{\adim}{\vdim}}{} |S(\grad f(x))|^2 \ud V(x,S) \le 4
        \textint{U(\delta)}{} f^2 \ud \|V\| \,.
    \end{displaymath}
\end{lemma}

\begin{proof}
    Suppose first that $U$ is bounded. Let $1 < \lambda < \infty$, and let $\phi$ be as
    in~\ref{new:cutoff} with $K = U$ and $\rho = \delta$; then $\phi \in \dspace{U(\delta)}{\R}$,
    $\phi = 1$ on~$U$ and $|\grad \phi| \le \lambda / \delta$. Since $\grad(\phi^2 f) = \phi^2 \grad
    f + 2 \phi f \grad \phi$, \ref{al:2.3(1)}\eqref{eq:2.3:1},\eqref{eq:2.3:2} give
    \begin{displaymath}
        |\phi(x) S(\grad f(x))|^2 = S(\grad f(x)) \bullet \grad(\phi^2 f)(x) - 2 \phi(x) f(x)
        S(\grad f(x)) \bullet \grad \phi(x)
    \end{displaymath}
    for $(x,S) \in U(\delta) \times \pgrass{\adim}{\vdim}$. As $\phi^2 f \in \dspace{U(\delta)}{\R}$
    and $\phi^2 f \ge 0$, the integral of the first term is nonpositive; by Schwarz's inequality,
    with $Q = \textint{}{} |\phi(x) S(\grad f(x))|^2 \ud V(x,S) < \infty$,
    \begin{displaymath}
        Q \le 2 \textint{}{} |\phi(x) S(\grad f(x))| \, f(x) |\grad \phi(x)| \ud V(x,S) \le 2
        Q^{1/2} \Big( \textint{}{} f^2 |\grad \phi|^2 \ud \|V\| \Big)^{1/2} \,.
    \end{displaymath}
    Hence $Q \le 4 \lambda^2 \delta^{-2} \textint{U(\delta)}{} f^2 \ud \|V\|$, and the left side of
    the assertion is at most $\delta^2 Q$. Let $\lambda \to 1$. For general~$U$, apply this to the
    bounded sets $U \cap \oball{0}{j}$, $j = 1, 2, \ldots$, noting $(U \cap \oball{0}{j})(\delta)
    \subseteq U(\delta)$, and use the increasing convergence theorem~\cite[2.4.7]{Federer1969}.
\end{proof}

\begin{theorem}[{\cite[7.5(6)]{Allard1972}}]
    \label{al:7.5(6)}
    Suppose $\vdim \ge 2$, $\beta = \vdim/(\vdim-1)$, $0 < \delta < \infty$, $0 \le \alpha < \infty$,
    \begin{displaymath}
        \delta V(g) \le \alpha \Big( \textint{}{} |g|^\beta \ud \|V\| \Big)^{1/\beta} \quad
        \text{whenever $g \in \VF(\R^{\adim})$ and $\spt g \subseteq U(\delta)$} \,,
    \end{displaymath}
    $\mathbf{c}_1 \alpha < 1$, $\udensity^{\vdim}(\|V\|,x) \ge 1$ for $\|V\|$ almost all $x \in
    U(\delta)$, and $f$ is a~positive function which is $V$~subharmonic in~$U(\delta)$. Then, with
    $\mathbf{c}(\alpha) = [20 (\mathbf{c}_1 / (1 - \mathbf{c}_1 \alpha)) 2^{\vdim-1}]^{\vdim}$,
    \begin{displaymath}
        \sup \{ f(x)^2 : x \in U \cap \spt \|V\| \} \le \mathbf{c}(\alpha) \delta^{-\vdim}
        \textint{U(\delta)}{} f^2 \ud \|V\| \,.
    \end{displaymath}
\end{theorem}

\begin{proof}
    If $U'$ is an~open subset of~$U$, the hypotheses hold with $U'$ in place of~$U$, and $U'(\delta)
    \subseteq U(\delta)$. Since $U = \tbcup_{j=1}^\infty U \cap \oball{0}{j}$, we may assume that
    $U$ is bounded. Let $U_\nu = U(2^{-\nu} \delta)$ for $\nu = 0, 1, 2, \ldots$, and $W_\nu =
    U(3 \cdot 2^{-\nu-1} \delta)$ for $\nu = 1, 2, \ldots$. Then $U \subseteq U_{\nu+1} \subseteq
    W_{\nu+1} \subseteq U_\nu$, $U_{\nu+1}(2^{-\nu-2} \delta) = W_{\nu+1}$ and $W_{\nu+1}(2^{-\nu-2}
    \delta) = U_\nu$. Let
    \begin{displaymath}
        I(\nu) = \Big( \textint{U_\nu}{} f^{2 \beta^\nu} \ud \|V\| \Big)^{1/\beta^\nu} \quad
        \text{for $\nu = 0, 1, 2, \ldots$} \,, \qquad \gamma = 20 \delta^{-1} \mathbf{c}_1 / (1 -
        \mathbf{c}_1 \alpha) \,.
    \end{displaymath}
    We may assume $I(0) < \infty$. We show
    \begin{equation}
        \label{eq:7.5(6):a}
        I(\nu + 1) \le \gamma^{1/\beta^\nu} 2^{\nu / \beta^\nu} I(\nu) \quad \text{for $\nu = 0,
        1, 2, \ldots$} \,.
    \end{equation}
    Fix~$\nu$ and let $g = f^{\beta^\nu}$. By~\ref{al:7.5(3)} with $C(t) = t^{\beta^\nu}$ on $\{ t
    : t > 0 \}$, $g$ is a~positive $V$~subharmonic function in~$U(\delta)$. By~\ref{al:7.5(5)} with
    $W_{\nu+1}$ and $2^{-\nu-2} \delta$ in place of $U$ and~$\delta$, and by Schwarz's inequality,
    \begin{multline*}
        \textint{W_{\nu+1} \times \pgrass{\adim}{\vdim}}{} |S(\grad g^2(x))| \ud V(x,S) \le 2 \Big(
        \textint{W_{\nu+1} \times \pgrass{\adim}{\vdim}}{} |S(\grad g(x))|^2 \ud V(x,S) \Big)^{1/2}
        \Big( \textint{W_{\nu+1}}{} g^2 \ud \|V\| \Big)^{1/2} \\ \le 2 \cdot 2 \cdot 2^{\nu+2}
        \delta^{-1} \textint{U_\nu}{} g^2 \ud \|V\| = 16 \cdot 2^\nu \delta^{-1} \textint{U_\nu}{}
        f^{2\beta^\nu} \ud \|V\| \,.
    \end{multline*}
    Let $1 < \lambda < \infty$ and let $\phi$ be as in~\ref{new:cutoff} with $K = U_{\nu+1}$ and
    $\rho = 2^{-\nu-2} \delta$; then $\phi g^2 \in \dspace{U(\delta)}{\R}$, $\phi = 1$
    on~$U_{\nu+1}$, $\spt \phi \subseteq W_{\nu+1}$ and $|\grad \phi| \le \lambda 2^{\nu+2}
    \delta^{-1}$. Since $S(\grad(\phi g^2)) = \phi S(\grad g^2) + g^2 S(\grad \phi)$,
    \ref{al:7.3} applied to $V \restrict (U(\delta) \times \pgrass{\adim}{\vdim}) \in
    \Var{\vdim}(U(\delta))$ and $\phi g^2$ yields
    \begin{displaymath}
        \Big( \textint{U_{\nu+1}}{} f^{2 \beta^{\nu+1}} \ud \|V\| \Big)^{1/\beta} \le
        \frac{\mathbf{c}_1}{1 - \mathbf{c}_1 \alpha} (16 + 4 \lambda) 2^\nu \delta^{-1}
        \textint{U_\nu}{} f^{2 \beta^\nu} \ud \|V\| \,.
    \end{displaymath}
    Letting $\lambda \to 1$ and raising to the power $1/\beta^\nu$ gives~\eqref{eq:7.5(6):a}.

    By induction, $I(\nu) \le \gamma^{s_\nu} 2^{t_\nu} I(0)$ with $s_\nu =
    \textsum{i=0}{\nu-1} \beta^{-i}$ and $t_\nu = \textsum{i=0}{\nu-1} i \beta^{-i}$. As $\nu \to
    \infty$, $s_\nu \to (1 - 1/\beta)^{-1} = \vdim$ and $t_\nu \uparrow \beta^{-1} (1 -
    \beta^{-1})^{-2} = \vdim(\vdim-1)$; hence
    \begin{displaymath}
        \limsup_{\nu \to \infty} I(\nu) \le \gamma^{\vdim} 2^{\vdim(\vdim-1)} I(0) =
        \mathbf{c}(\alpha) \delta^{-\vdim} \textint{U(\delta)}{} f^2 \ud \|V\| \,.
    \end{displaymath}
    Finally, if $x \in U \cap \spt \|V\|$ and $0 < c < f(x)^2$, there is $\rho > 0$ with
    $\cball{x}{\rho} \subseteq U$ and $f^2 > c$ on $\cball{x}{\rho}$; then $I(\nu) \ge c (\|V\|
    \cball{x}{\rho})^{1/\beta^\nu}$ for every~$\nu$, and $0 < \|V\| \cball{x}{\rho} < \infty$, so
    $\limsup_{\nu \to \infty} I(\nu) \ge c$.
\end{proof}

\begin{remark}
    \label{rem:7.5(6)}
    Allard writes $U_0 = U$; the definition of $I(0)$ and the conclusion require $U_0 =
    U(\delta)$. He remarks that~\ref{al:4.4}, \ref{al:4.7} and~\ref{al:7.5(2)} allow one to apply the
    estimate to submanifolds of~$\R^{\adim}$.
\end{remark}

\section{The regularity theorem}
\label{sec:regularity}

\begin{miniremark}
    \label{mr:sec8}
    This section treats~\cite[section~8]{Allard1972}. Throughout it, $1 \le \vdim < p < \infty$,
    $q = p/(p-1)$, $p \ge 2$ in case $\vdim = 1$, and $\mathbf{c}_1$ is as in~\ref{mr:sec7}. The
    purpose of the section is the regularity theorem~\ref{al:8.RT}: a~varifold whose density is
    at least~$d$ almost everywhere in a~ball, whose weight in the ball is close to that of
    a~$\vdim$~disc of multiplicity~$d$, and whose first variation is small in $\Lp{p}$, coincides
    in a~smaller ball with the graph of a~continuously differentiable map with H\"older
    continuous derivative. Theorem~\ref{al:8.8} provides density bounds and approximating planes
    at all points and scales; \ref{al:8.12} and~\ref{al:8.16} contain the approximation by
    Lipschitzian graphs and the decay of the height excess; \ref{al:8.17} iterates the decay, and
    \ref{al:8.19} yields the graph representation. Allard states the regularity theorem at the
    beginning of the section and comments on it in~\cite[8.1]{Allard1972}; here the theorem is
    proved in~\ref{al:8.RT} after~\ref{al:8.19}, and the comments follow it in~\ref{al:8.1(1)}--\ref{al:8.1(5)}.
    The harmonic function theory used in~\ref{al:8.16}, which Allard quotes from Federer's
    chapter~5, is supplied by the new lemmas~\ref{new:ibp}--\ref{new:rellich}.
\end{miniremark}

\begin{lemma}[{\cite[8.2]{Allard1972}}]
    \label{al:8.2}
    Suppose $V \in \Var{\vdim}(\R^{\adim})$, $0 \le \alpha < \infty$, $0 < r < \infty$, $b \in
    \R^{\adim}$, $\iota \in \Hom(\R^{\adim},\R^{\adim})$ is a~linear isometry, and $G$ is an open
    subset of~$\R^{\adim}$. Then
    \begin{displaymath}
        \delta V(g) \le \alpha \Big( \textint{}{} |g|^q \ud \|V\| \Big)^{1/q} \quad \text{for $g \in
        \VF(\R^{\adim})$ with $\spt g \subseteq G$}
    \end{displaymath}
    if and only if
    \begin{displaymath}
        \delta(\scale{r \#} V)(g) \le \alpha r^{\vdim/p - 1} \Big( \textint{}{} |g|^q \ud
        \|\scale{r \#} V\| \Big)^{1/q} \quad \text{for $g \in \VF(\R^{\adim})$ with $\spt g
        \subseteq \scale{r}[G]$} \,;
    \end{displaymath}
    the same holds with $\scale{r}$ and $\alpha r^{\vdim/p-1}$ replaced by $\trans{b}$ and~$\alpha$,
    or by $\iota$ and~$\alpha$.
\end{lemma}

\begin{proof}
    Let $g \in \VF(\R^{\adim})$ and $\tilde g = r^{-1} g \circ \scale{r}$. Since $\uD \tilde g(x) =
    \uD g(r x)$, \ref{al:3.2(2)} and~\ref{al:4.2:lem} give $\delta(\scale{r \#} V)(g) = r^{\vdim}
    \textint{}{} \uD g(r x) \bullet S \ud V(x,S) = r^{\vdim} \delta V(\tilde g)$, and $\textint{}{}
    |\tilde g|^q \ud \|V\| = r^{-q-\vdim} \textint{}{} |g|^q \ud \|\scale{r \#} V\|$. Since $\spt
    \tilde g = \scale{1/r}[\spt g]$ and $r^{\vdim} r^{-1-\vdim/q} = r^{\vdim/p-1}$, the first
    equivalence follows, the converse by applying the argument to $\scale{1/r}$. For $\trans{b}$
    and~$\iota$ one uses $\tilde g = g \circ \trans{b}$ and $\tilde g = \iota^{-1} \circ g \circ
    \iota$, together with $\iota_\# V(\alpha') = \textint{}{} \alpha'(\iota(x), \iota \circ S \circ
    \iota^{-1}) \ud V(x,S)$ by~\ref{al:3.2(1)}, and $(\iota^{-1} \circ \uD g(\iota(x)) \circ \iota)
    \bullet S = \uD g(\iota(x)) \bullet (\iota \circ S \circ \iota^{-1})$.
\end{proof}

\begin{lemma}[{\cite[8.3]{Allard1972}}]
    \label{al:8.3}
    Suppose $0 < d < \infty$, $0 \le \alpha < \infty$, $\Delta = d^{1/\vdim} / \mathbf{c}_1 - \alpha
    > 0$, $V \in \Var{\vdim}(\R^{\adim})$, $a \in \spt \|V\|$, $0 < s < \infty$,
    \begin{gather*}
        \delta V(g) \le \alpha \sup \{ |g(x)| : x \in \R^{\adim} \} \quad \text{for $g \in
        \VF(\R^{\adim})$ with $\spt g \subseteq \oball{a}{s}$, in case $\vdim = 1$} \,, \\
        \delta V(g) \le \alpha \Big( \textint{}{} |g|^{\vdim/(\vdim-1)} \ud \|V\|
        \Big)^{(\vdim-1)/\vdim} \quad \text{for $g \in \VF(\R^{\adim})$ with $\spt g \subseteq
        \oball{a}{s}$, in case $\vdim > 1$} \,,
    \end{gather*}
    and $\udensity^{\vdim}(\|V\|,x) \ge d$ for $\|V\|$ almost all $x \in \oball{a}{s}$. Then
    \begin{displaymath}
        r^{-\vdim} \|V\| \cball{a}{r} \ge (\Delta/\vdim)^{\vdim} \quad \text{for $0 < r < s$} \,.
    \end{displaymath}
\end{lemma}

\begin{proof}
    Let $\psi(r) = \|V\| \oball{a}{r}$ for $0 < r \le s$; $\psi$ is nondecreasing and positive,
    because $a \in \spt \|V\|$. The hypothesis gives $\|\delta V\| \oball{a}{r} \le \alpha
    \psi(r)^{1 - 1/\vdim}$ for $0 < r \le s$, with $\psi(r)^0 = 1$; in particular $\|\delta V\|$ is
    a~Radon measure on $\oball{a}{s}$. Let $V_r = V \restrict (\oball{a}{r} \times
    \pgrass{\adim}{\vdim})$. By~\cite[2.9.19]{Federer1969}, $\psi'(r)$ exists for $\LM^1$ almost all
    $0 < r < s$. Fix such~$r$. For $g \in \VF(\R^{\adim})$ with $|g| \le 1$, \ref{al:4.10(1)} applied
    to the restriction of~$V$ to $\oball{a}{s}$, with $f(x) = -|x - a|$ and level~$-r$, gives
    \begin{displaymath}
        |\delta V_r(g)| \le \|\delta V\| \oball{a}{r} + \limsup_{h \to 0+} h^{-1} \|V\| \{ x : r -
        h < |x - a| < r \} \le \alpha \psi(r)^{1-1/\vdim} + \psi'(r) \,;
    \end{displaymath}
    here $\spt g$ need not lie in $\oball{a}{s}$, since $V_r$ only involves $g$ near $\cball{a}{r}$,
    and we may multiply~$g$ by a~cut-off function equal to~$1$ near $\cball{a}{r}$ with support in
    $\oball{a}{s}$, by~\ref{new:cutoff}. Hence $\|\delta V_r\|(\R^{\adim}) \le \alpha
    \psi(r)^{1-1/\vdim} + \psi'(r)$.

    Let $\phi$ be as in~\ref{new:cutoff} with $K = \cball{a}{r}$. Then $\grad \phi = 0$ near $\spt
    \|V_r\|$, and $\udensity^{\vdim}(\|d^{-1} V_r\|,x) \ge 1$ for $\|V_r\|$ almost all~$x$, because
    $\oball{a}{r}$ is open. Applying~\ref{al:7.1} to $d^{-1} V_r$ and~$\phi$ we obtain $d^{-1}
    \psi(r) \le \mathbf{c}_1 (d^{-1} \psi(r))^{1/\vdim} d^{-1} \|\delta V_r\|(\R^{\adim})$, that is
    \begin{displaymath}
        d^{1/\vdim} \psi(r)^{1 - 1/\vdim} \le \mathbf{c}_1 \big( \alpha \psi(r)^{1-1/\vdim} +
        \psi'(r) \big) \,, \qquad \Delta \psi(r)^{1 - 1/\vdim} \le \psi'(r) \,.
    \end{displaymath}
    Consequently $(\psi^{1/\vdim})'(r) = \vdim^{-1} \psi(r)^{1/\vdim - 1} \psi'(r) \ge \Delta /
    \vdim$ for $\LM^1$ almost all $0 < r < s$. Since $\psi^{1/\vdim}$ is nondecreasing,
    \cite[2.9.19]{Federer1969} gives $(r - u) \Delta / \vdim \le \textint{u}{r}
    (\psi^{1/\vdim})' \ud \LM^1 \le \psi(r)^{1/\vdim} - \psi(u)^{1/\vdim}$ for $0 < u < r < s$.
    Letting $u \to 0+$ yields $\|V\| \cball{a}{r} \ge \psi(r) \ge (\Delta / \vdim)^{\vdim} r^{\vdim}$.
\end{proof}

\begin{remark}
    \label{rem:8.3}
    Allard treats only the case $\vdim > 1$ and assumes $\density^{\vdim}(\|V\|,x) \ge d$. He
    remarks that a~similar use of an isoperimetric inequality occurs in the theory of integral
    currents, and cites~\cite[4.1.25]{Federer1969}. That item is a~lemma on the push-forward of
    the current $\HM^{\vdim} \wedge \xi$ under a~Lipschitzian map and contains no such argument;
    the printed locator appears to be in error.
    % UNVERIFIED: the item of Federer which Allard intended here (4.1.25 was read and does not fit).
\end{remark}

\begin{lemma}[new]
    \label{new:conv_plane}
    Suppose $0 < L < \infty$, $T \in \pgrass{\adim}{\vdim}$, $0 < c < \infty$, $V_1, V_2, \ldots \in
    \Var{\vdim}(\oball{0}{L})$, $\lim_{i \to \infty} V_i = \var{}(\im T \cap \oball{0}{L})$, and
    \begin{displaymath}
        r^{-\vdim} \|V_i\| \cball{b}{r} \ge c \quad \text{whenever $i$ is a positive integer, $b \in
        \spt \|V_i\|$ and $0 < r < L - |b|$} \,.
    \end{displaymath}
    Then for $0 < \rho < L$ and $0 < \delta < \infty$ there is~$i_0$ such that $\dist(x, \im T) <
    \delta$ whenever $i \ge i_0$ and $x \in \spt \|V_i\| \cap \cball{0}{\rho}$, and $\dist(y, \spt
    \|V_i\|) < \delta$ whenever $i \ge i_0$ and $y \in \im T \cap \cball{0}{\rho}$.
\end{lemma}

\begin{proof}
    We may assume $\delta < L - \rho$. Suppose the first assertion fails. Then there are $i(1) < i(2)
    < \ldots$ and $x_j \in \spt \|V_{i(j)}\| \cap \cball{0}{\rho}$ with $\dist(x_j, \im T) \ge
    \delta$; passing to a~subsequence, $x_j \to x$ with $|x| \le \rho$ and $\dist(x, \im T) \ge
    \delta$. Let $r = \delta / 4$ and $K = \cball{x}{2r}$; then $K$ is a~compact subset of
    $\oball{0}{L}$ with $K \cap \im T = \varnothing$, and $\cball{x_j}{r} \subseteq K$ for
    large~$j$. By~\ref{al:2.6(2)(a)}(c), $0 = \HM^{\vdim}(K \cap \im T) \ge \limsup_{j \to \infty}
    \|V_{i(j)}\|(K) \ge c r^{\vdim}$, a~contradiction. Suppose the second assertion fails. Then there
    are $i(1) < i(2) < \ldots$ and $y_j \in \im T \cap \cball{0}{\rho}$ with $\oball{y_j}{\delta}
    \cap \spt \|V_{i(j)}\| = \varnothing$; passing to a~subsequence, $y_j \to y$, and
    $\oball{y}{\delta/2} \cap \spt \|V_{i(j)}\| = \varnothing$ for large~$j$. By~\ref{al:2.6(2)(a)}(c),
    $\HM^{\vdim}(\im T \cap \oball{y}{\delta/2}) \le \liminf_{j \to \infty} \|V_{i(j)}\|
    \oball{y}{\delta/2} = 0$, a~contradiction, because $y \in \im T$ and $\oball{y}{\delta/2}
    \subseteq \oball{0}{L}$.
\end{proof}

\begin{lemma}[{\cite[8.4]{Allard1972}}]
    \label{al:8.4}
    Whenever $0 < \epsilon < 1$, there is $\eta > 0$ with the following property. If $V \in
    \Var{\vdim}(\oball{0}{1})$, $0 \in \spt \|V\|$, $\|\delta V\| \oball{0}{1} \le \eta$,
    $\udensity^{\vdim}(\|V\|,x) \ge 1$ for $\|V\|$ almost all $x \in \oball{0}{1}$, $\|V\|
    \oball{0}{1} \le (1 + \eta) \unitmeasure{\vdim}$,
    \begin{gather*}
        \delta V(g) \le (\mathbf{c}_1 + \epsilon)^{-1} \sup \{ |g(x)| : x \in \oball{0}{1} \} \quad
        \text{for $g \in \VF(\oball{0}{1})$, in case $\vdim = 1$} \,, \\
        \delta V(g) \le (\mathbf{c}_1 + \epsilon)^{-1} \Big( \textint{}{} |g|^{\vdim/(\vdim-1)} \ud
        \|V\| \Big)^{(\vdim-1)/\vdim} \quad \text{for $g \in \VF(\oball{0}{1})$, in case $\vdim >
        1$} \,,
    \end{gather*}
    then there is $T \in \pgrass{\adim}{\vdim}$ such that $\dist(x, \im T) \le \epsilon$ for $x \in
    \spt \|V\| \cap \oball{0}{1 - \epsilon}$, and $\dist(y, \spt \|V\|) \le \epsilon$ for $y \in
    \im T \cap \oball{0}{1-\epsilon}$.
\end{lemma}

\begin{proof}
    Suppose the assertion fails for some~$\epsilon$. Then for $i = 1, 2, \ldots$ there is $V_i$
    satisfying the hypotheses with $\eta = 1/i$, for which every $T \in \pgrass{\adim}{\vdim}$
    violates one of the two conclusions. By~\ref{al:2.6(2)(a)}(a) and~\ref{new:weak_metric}, since
    $\|V_i\| \oball{0}{1} \le 2 \unitmeasure{\vdim}$, we may assume that $V_i$ converges to some $C
    \in \Var{\vdim}(\oball{0}{1})$. By~\ref{al:4.11}, $C$ is stationary; by~\ref{al:2.6(2)(a)}(c),
    $\|C\| \oball{0}{1} \le \unitmeasure{\vdim}$. Since $\|V_i\|(\oball{0}{1} \cap \{ x :
    \udensity^{\vdim}(\|V_i\|,x) < 1 \}) = 0$, the proof of~\ref{al:5.4} applies with upper
    densities and shows $\density^{\vdim}(\|C\|,x) \ge 1$ for $\|C\|$ almost all~$x$; the density
    exists by~\ref{al:5.1(2)}.

    Let $\Delta = \mathbf{c}_1^{-1} - (\mathbf{c}_1 + \epsilon)^{-1} > 0$. Applying~\ref{al:8.3} with
    $d = 1$, $\alpha = (\mathbf{c}_1 + \epsilon)^{-1}$ and $s = 1 - |b|$ to $V_i$ extended by zero to
    $\R^{\adim}$, we obtain
    \begin{equation}
        \label{eq:8.4:a}
        r^{-\vdim} \|V_i\| \cball{b}{r} \ge (\Delta/\vdim)^{\vdim} \quad \text{whenever $b \in \spt
        \|V_i\|$ and $0 < r < 1 - |b|$} \,.
    \end{equation}
    By~\ref{al:2.6(2)(a)}(c), $\|C\| \cball{0}{r} \ge \limsup_{i \to \infty} \|V_i\| \cball{0}{r}
    \ge (\Delta/\vdim)^{\vdim} r^{\vdim}$ for $0 < r < 1$, so $0 \in \spt \|C\|$. Thus~\ref{al:5.3}
    with $R = 1$ applies and yields $T \in \pgrass{\adim}{\vdim}$ with $C = \density^{\vdim}(\|C\|,0)
    \var{}(\im T \cap \oball{0}{1})$. Since $\density^{\vdim}(\|C\|,\cdot) \ge 1$ almost everywhere
    and $\|C\| \oball{0}{1} \le \unitmeasure{\vdim}$, $\density^{\vdim}(\|C\|,0) = 1$.
    Now~\ref{new:conv_plane} with $L = 1$, $\rho = 1 - \epsilon$ and~\eqref{eq:8.4:a} shows that
    both conclusions hold with this~$T$ for large~$i$, a~contradiction.
\end{proof}

\begin{theorem}[{\cite[8.5]{Allard1972}}]
    \label{al:8.5}
    Suppose $0 < d < \infty$, $0 \le \alpha < \infty$, $V \in \Var{\vdim}(\R^{\adim})$, $a \in \spt
    \|V\|$, $0 < s < \infty$,
    \begin{displaymath}
        \delta V(g) \le \alpha \Big( \textint{}{} |g|^q \ud \|V\| \Big)^{1/q} \quad \text{for $g \in
        \VF(\R^{\adim})$ with $\spt g \subseteq \oball{a}{s}$} \,,
    \end{displaymath}
    $\udensity^{\vdim}(\|V\|,x) \ge d$ for $\|V\|$ almost all $x \in \oball{a}{s}$,
    \begin{displaymath}
        \Delta(a,s) = d^{1/\vdim} / \mathbf{c}_1 - \alpha \|V\| \oball{a}{s}^{1/\vdim - 1/p} > 0
        \,, \qquad \Phi(a,s) = \frac{p}{p - \vdim} \Big( \frac{\Delta(a,s)}{\vdim} \Big)^{-\vdim/p}
        \alpha \,.
    \end{displaymath}
    Then $[r \mapsto r^{-\vdim} \|V\| \cball{a}{r} \exp(\Phi(a,s) r^{1 - \vdim/p})]$ is
    nondecreasing on $\{ r : 0 < r < s \}$.
\end{theorem}

\begin{proof}
    Suppose $g \in \VF(\R^{\adim})$ and $\spt g \subseteq \oball{a}{s}$. If $\vdim > 1$, then $q <
    \vdim/(\vdim-1)$, and H\"older's inequality~\cite[2.4.14]{Federer1969} gives
    \begin{displaymath}
        \delta V(g) \le \alpha \|V\| \oball{a}{s}^{1/\vdim - 1/p} \Big( \textint{}{}
        |g|^{\vdim/(\vdim-1)} \ud \|V\| \Big)^{(\vdim-1)/\vdim} \,,
    \end{displaymath}
    because $1/q - (\vdim-1)/\vdim = 1/\vdim - 1/p$; if $\vdim = 1$, then $\delta V(g) \le \alpha
    \|V\| \oball{a}{s}^{1/q} \sup |g|$ and $1/q = 1/\vdim - 1/p$. Hence~\ref{al:8.3} gives
    $r^{-\vdim} \|V\| \cball{a}{r} \ge (\Delta(a,s)/\vdim)^{\vdim}$ for $0 < r < s$. The hypothesis
    implies $\|\delta V\|(O) \le \alpha \|V\|(O)^{1/q}$ for open $O \subseteq \oball{a}{s}$, hence
    $\|\delta V\| \cball{a}{t} \le \alpha \|V\| \cball{a}{t}^{1/q}$ for $0 < t < s$.
    By~\ref{al:4.3}\eqref{eq:4.3:2} and $|\kappa(V,a,t)| \le \|V\| \cball{a}{t}^{-1}$,
    \begin{displaymath}
        \delta V(\kappa(V,a,t)) \le \alpha \|V\| \cball{a}{t}^{1/q - 1} \le \alpha
        (\Delta(a,s)/\vdim)^{-\vdim/p} t^{-\vdim/p} \quad \text{for $0 < t < s$} \,,
    \end{displaymath}
    since $1/q - 1 = -1/p$. Thus $\textint{u}{r} \delta V(\kappa(V,a,t)) \ud \LM^1 t \le \Phi(a,s)
    (r^{1-\vdim/p} - u^{1-\vdim/p})$ for $0 < u < r < s$, and~\ref{al:5.1(1)}, applied to the
    restriction of~$V$ to $\oball{a}{s}$, gives $r^{-\vdim} \|V\| \cball{a}{r} \exp(\Phi(a,s)
    r^{1-\vdim/p}) \ge u^{-\vdim} \|V\| \cball{a}{u} \exp(\Phi(a,s) u^{1 - \vdim/p})$.
\end{proof}

\begin{theorem}[{\cite[8.6]{Allard1972}}]
    \label{al:8.6}
    Suppose $V \in \Var{\vdim}(\R^{\adim})$, $G$ is an open subset of~$\R^{\adim}$, $0 \le \alpha <
    \infty$, $0 < d < \infty$,
    \begin{displaymath}
        \delta V(g) \le \alpha \Big( \textint{}{} |g|^q \ud \|V\| \Big)^{1/q} \quad \text{for $g \in
        \VF(\R^{\adim})$ with $\spt g \subseteq G$} \,,
    \end{displaymath}
    and $\udensity^{\vdim}(\|V\|,x) \ge d$ for $\|V\|$ almost all $x \in G$. Then
    $\density^{\vdim}(\|V\|,x)$ exists in~$\R$ for every $x \in G$, $\density^{\vdim}(\|V\|,\cdot)$
    is upper semicontinuous on~$G$, and $\density^{\vdim}(\|V\|,x) \ge d$ for $x \in G \cap \spt
    \|V\|$.
\end{theorem}

\begin{proof}
    If $x \in G \without \spt \|V\|$, the density vanishes near~$x$. Suppose $a \in G \cap \spt
    \|V\|$. First, $\|V\| \{ a \} = 0$: otherwise $m = \|V\| \{ a \} > 0$, and for $0 < u < r$ with
    $\oball{a}{r} \subseteq G$ the argument of~\ref{al:8.5} gives $\delta V(\kappa(V,a,t)) \le
    \alpha \|V\| \cball{a}{t}^{-1/p} \le \alpha m^{-1/p}$, so~\ref{al:5.1(1)} yields $u^{-\vdim} m
    \le r^{-\vdim} \|V\| \cball{a}{r} \exp(\alpha m^{-1/p} r)$, which is impossible as $u \to 0+$.
    Hence we may choose $s_0 > 0$ with $\oball{a}{2 s_0} \subseteq G$ and $\alpha \|V\| \oball{a}{2
    s_0}^{1/\vdim - 1/p} \le d^{1/\vdim}/(2 \mathbf{c}_1)$. Let $\Delta_0 = d^{1/\vdim} / (2
    \mathbf{c}_1)$ and $\Phi_0 = p (p - \vdim)^{-1} (\Delta_0/\vdim)^{-\vdim/p} \alpha$. If $b \in
    \spt \|V\| \cap \oball{a}{s_0}$, then $\oball{b}{s_0} \subseteq \oball{a}{2 s_0}$,
    $\Delta(b,s_0) \ge \Delta_0$ and $\Phi(b,s_0) \le \Phi_0$, so~\ref{al:8.5} shows that
    $\density^{\vdim}(\|V\|,b)$ exists and
    \begin{displaymath}
        \unitmeasure{\vdim} \density^{\vdim}(\|V\|,b) \le r^{-\vdim} \|V\| \cball{b}{r} \exp(\Phi_0
        r^{1-\vdim/p}) \quad \text{for $0 < r < s_0$} \,;
    \end{displaymath}
    this also holds for $b \in \oball{a}{s_0} \without \spt \|V\|$. Since $[b \mapsto \|V\|
    \cball{b}{r}]$ is upper semicontinuous, $\limsup_{b \to a} \density^{\vdim}(\|V\|,b) \le
    \unitmeasure{\vdim}^{-1} r^{-\vdim} \|V\| \cball{a}{r} \exp(\Phi_0 r^{1-\vdim/p})$, and letting
    $r \to 0+$ gives upper semicontinuity at~$a$. Finally the points $x \in G \cap \spt \|V\|$ with
    $\density^{\vdim}(\|V\|,x) \ge d$ are dense in $G \cap \spt \|V\|$, because their complement in $G$
    is $\|V\|$ null; upper semicontinuity gives $\density^{\vdim}(\|V\|,a) \ge d$.
\end{proof}

\begin{definition}[{\cite[8.7]{Allard1972}}]
    \label{al:8.7}
    Let $\affgrass{\adim}{\vdim}$ be the family of sets $\trans{a}[\im S]$ corresponding to $a \in
    \R^{\adim}$ and $S \in \pgrass{\adim}{\vdim}$. Members $X$, $Y$ of $\affgrass{\adim}{\vdim}$ are
    \emph{parallel} if and only if $X = \trans{a}[Y]$ for some $a \in \R^{\adim}$; we say that $X$ is
    parallel to $S \in \pgrass{\adim}{\vdim}$ if $X$ is parallel to~$\im S$.
\end{definition}

\begin{theorem}[{\cite[8.8]{Allard1972}}]
    \label{al:8.8}
    Whenever $0 < \epsilon < 1$, there is $\eta > 0$ with the following property. Suppose $V \in
    \Var{\vdim}(\R^{\adim})$, $a \in \spt \|V\|$, $0 < R < \infty$, $0 < d < \infty$,
    \begin{gather*}
        \delta V(g) \le \eta d^{1/p} R^{\vdim/p - 1} \Big( \textint{}{} |g|^q \ud \|V\| \Big)^{1/q}
        \quad \text{for $g \in \VF(\R^{\adim})$ with $\spt g \subseteq \oball{a}{R}$} \,, \\
        \udensity^{\vdim}(\|V\|,x) \ge d \quad \text{for $\|V\|$ almost all $x \in \oball{a}{R}$} \,,
        \qquad R^{-\vdim} \|V\| \oball{a}{R} \le (1 + \eta) \unitmeasure{\vdim} d \,.
    \end{gather*}
    Then
    \begin{displaymath}
        (1 - \epsilon) \unitmeasure{\vdim} d \le r^{-\vdim} \|V\| \cball{b}{r} \le (1 + \epsilon)
        \unitmeasure{\vdim} d
    \end{displaymath}
    whenever $b \in \oball{a}{(1-\epsilon) R} \cap \spt \|V\|$ and $0 < r < (1 - \epsilon)(R - |a
    - b|)$. Moreover, to each such $b$ and~$r$ corresponds $X(b,r) \in \affgrass{\adim}{\vdim}$
    with $b \in X(b,r)$,
    \begin{gather*}
        \dist(x, X(b,r)) \le \epsilon r \quad \text{for $x \in \oball{b}{r} \cap \spt \|V\|$} \,,
        \\ \dist(y, \spt \|V\|) \le \epsilon r \quad \text{for $y \in \oball{b}{r} \cap X(b,r)$}
        \,.
    \end{gather*}
\end{theorem}

\begin{proof}
    By~\ref{al:8.2}, \ref{al:3.2(2)} and the identity $\delta(c V) = c \, \delta V$ for $0 < c <
    \infty$, the hypotheses and conclusions are invariant under replacing $V$ by $d^{-1}
    (\scale{1/R} \circ \trans{-a})_\# V$; hence we may assume $a = 0$, $R = 1$ and $d = 1$. Then
    $\density^{\vdim}(\|V\|,\cdot)$ exists on $\oball{0}{1}$ and is at least~$1$ on $\spt \|V\|
    \cap \oball{0}{1}$, by~\ref{al:8.6}. For $\eta \le 1$ and $b \in \spt \|V\| \cap \oball{0}{1}$
    let $s = 1 - |b|$; then $\Delta(b,s) \ge \Delta_1 = \mathbf{c}_1^{-1} - \eta (2
    \unitmeasure{\vdim})^{1/\vdim - 1/p}$ and $\Phi(b,s) \le \Phi_1 = p (p-\vdim)^{-1}
    (\Delta_1/\vdim)^{-\vdim/p} \eta$, provided $\Delta_1 > 0$. By~\ref{al:8.5},
    \begin{equation}
        \label{eq:8.8:a}
        \unitmeasure{\vdim} \le \unitmeasure{\vdim} \density^{\vdim}(\|V\|,b) \le r^{-\vdim} \|V\|
        \cball{b}{r} \exp(\Phi_1) \quad \text{for $0 < r < 1 - |b|$} \,.
    \end{equation}
    Choosing $\eta$ so small that $\Delta_1 > 0$ and $\exp(-\Phi_1) \ge 1 - \epsilon$, we obtain the
    lower bound for every $b \in \spt \|V\| \cap \oball{0}{1}$ and $0 < r < 1 - |b|$. Moreover, for
    such $b$ and~$r$, the monotonicity in~\ref{al:8.5} and $|b| < 1 - \epsilon$ give
    \begin{equation}
        \label{eq:8.8:b}
        r^{-\vdim} \|V\| \cball{b}{r} \le \lim_{t \to (1-|b|)-} t^{-\vdim} \|V\| \cball{b}{t}
        \exp(\Phi_1) \le \epsilon^{-\vdim} (1 + \eta) \unitmeasure{\vdim} \exp(\Phi_1) \,.
    \end{equation}

    Suppose the remaining assertions fail for some $\epsilon$ and every $\eta > 0$. Then there are
    $V_i$ satisfying the normalised hypotheses with $\eta = \eta_i$, where $\eta_i \to 0$, and $b_i
    \in \oball{0}{1-\epsilon} \cap \spt \|V_i\|$, $0 < r_i < (1 - \epsilon)(1 - |b_i|)$, such that
    either $r_i^{-\vdim} \|V_i\| \cball{b_i}{r_i} > (1 + \epsilon) \unitmeasure{\vdim}$ or no
    $X(b_i,r_i)$ exists. The corresponding numbers $\Phi_{1,i}$ tend to~$0$. As in the proof
    of~\ref{al:8.4}, using~\eqref{eq:8.8:a}, \ref{al:4.11}, \ref{al:5.4} and~\ref{al:5.3}, we may
    assume that the restrictions of $V_i$ to $\oball{0}{1}$ converge to $\var{}(\im T \cap
    \oball{0}{1})$ for some $T \in \pgrass{\adim}{\vdim}$, that $b_i \to b$ and $r_i \to r_\infty$,
    and~\eqref{eq:8.8:a} gives $r^{-\vdim} \|V_i\| \cball{x}{r} \ge \unitmeasure{\vdim}/2$ for $x
    \in \spt \|V_i\|$, $0 < r < 1 - |x|$ and large~$i$. Note $|b_i| + r_i < 1 - \epsilon^2$.

    Suppose $r_\infty > 0$. For $\delta > 0$ and large~$i$, $\cball{b_i}{r_i} \subseteq
    \cball{b}{r_\infty + \delta}$, so~\ref{al:2.6(2)(a)}(c) gives $\limsup_{i \to \infty} \|V_i\|
    \cball{b_i}{r_i} \le \HM^{\vdim}(\im T \cap \cball{b}{r_\infty + \delta}) \le
    \unitmeasure{\vdim} (r_\infty + \delta)^{\vdim}$, hence $\limsup_{i \to \infty} r_i^{-\vdim}
    \|V_i\| \cball{b_i}{r_i} \le \unitmeasure{\vdim}$. Let $X_i = \trans{T^\perp(b_i)}[\im T]$, so
    $b_i \in X_i$. By~\ref{new:conv_plane} with $L = 1$ and $\rho = 1 - \epsilon^2$,
    $\dist(b_i, \im T) \to 0$ and both suprema in the conclusion, computed with $X_i$, tend to~$0$,
    while $\epsilon r_i \to \epsilon r_\infty > 0$. Thus the conclusions hold for large~$i$,
    a~contradiction.

    Suppose $r_\infty = 0$. Let $L = (1 - \epsilon)^{-1}$, so that $\oball{b_i}{L r_i} \subseteq
    \oball{0}{1}$, and let $W_i \in \Var{\vdim}(\oball{0}{L})$ be the restriction of $(\scale{1/r_i}
    \circ \trans{-b_i})_\# V_i$ to $\oball{0}{L} \times \pgrass{\adim}{\vdim}$. Then $0 \in \spt
    \|W_i\|$, $\udensity^{\vdim}(\|W_i\|,y) \ge 1$ for $\|W_i\|$ almost all~$y$, and by~\ref{al:8.2}
    $\delta W_i(g) \le \eta_i r_i^{1 - \vdim/p} (\textint{}{} |g|^q \ud \|W_i\|)^{1/q}$ for $g \in
    \VF(\oball{0}{L})$. By~\eqref{eq:8.8:b} the masses $\|W_i\| \oball{0}{L}$ are bounded, so
    $\|\delta W_i\| \oball{0}{L} \to 0$. If $y \in \spt \|W_i\|$ and $0 < r < L - |y|$, then
    $r_i r < 1 - |b_i + r_i y|$, so $r^{-\vdim} \|W_i\| \cball{y}{r} \ge \unitmeasure{\vdim}/2$ for
    large~$i$ by~\eqref{eq:8.8:a}. Fix $t_0$ and~$\delta$ with $0 < t_0$ and $t_0 + \delta < 1 -
    |b|$. By the monotonicity in~\ref{al:8.5} at~$b_i$ and~\ref{al:2.6(2)(a)}(c),
    \begin{displaymath}
        \limsup_{i \to \infty} (L r_i)^{-\vdim} \|V_i\| \cball{b_i}{L r_i} \le \limsup_{i \to \infty}
        t_0^{-\vdim} \|V_i\| \cball{b_i}{t_0} \exp(\Phi_{1,i}) \le \unitmeasure{\vdim} (1 +
        \delta/t_0)^{\vdim} \,,
    \end{displaymath}
    hence $\limsup_{i \to \infty} \|W_i\| \cball{0}{L} \le \unitmeasure{\vdim} L^{\vdim}$. Passing to
    a~subsequence, $W_i \to C'$, where $C'$ is stationary by~\ref{al:4.11}, $\density^{\vdim}(\|C'\|,y)
    \ge 1$ for $\|C'\|$ almost all~$y$ by~\ref{al:5.4} and~\ref{al:5.1(2)}, $0 \in \spt \|C'\|$, and
    $\|C'\| \oball{0}{L} \le \unitmeasure{\vdim} L^{\vdim}$. By~\ref{al:5.3}, $C' = \var{}(\im T'
    \cap \oball{0}{L})$ for some $T' \in \pgrass{\adim}{\vdim}$. As in the case $r_\infty > 0$,
    $\limsup_{i \to \infty} r_i^{-\vdim} \|V_i\| \cball{b_i}{r_i} = \limsup_{i \to \infty} \|W_i\|
    \cball{0}{1} \le \unitmeasure{\vdim}$, and~\ref{new:conv_plane} with $\rho = 1$ shows that
    $X(b_i,r_i) = \trans{b_i}[\im T']$ satisfies the conclusions for large~$i$, a~contradiction.
\end{proof}

\begin{remark}
    \label{rem:8.8}
    Allard assumes only $a \in \R^{\adim}$, and writes $\oball{a}{r}$ for $\oball{a}{R}$ in the density
    hypothesis. Without $a \in \spt \|V\|$ the assertion fails: if $\vdim = 1$, $\adim = 2$, and $V$ is
    the sum of the varifolds of two lines meeting at a~point~$c$ with $|c| = 0.95$, both at distance
    at least $0.9$ from~$0$, then $V$ is stationary with density at least~$1$, $\|V\| \oball{0}{1} < 2
    = \unitmeasure{1}$, while $\density^1(\|V\|,c) = 2$. Allard's proof reads ``Combine 5.3, 8.2, 8.4,
    and 8.5''. If $\epsilon$ is small, Reifenberg's topological disc
    theorem~\cite[p.~64, Theorem]{Reifenberg1960a} yields $\rho > 0$ such that $\oball{b}{\rho} \cap
    \spt \|V\|$ is homeomorphic to~$\R^{\vdim}$ for $b \in \oball{a}{(1-\epsilon)R} \cap \spt \|V\|$.
\end{remark}

\begin{lemma}[{\cite[8.9(1)]{Allard1972}}]
    \label{al:8.9(1)}
    If $S, T \in \pgrass{\adim}{\vdim}$, then $|S - T|^2 = 2 S \bullet T^\perp = 2 S^\perp \bullet T
    = |S^\perp - T^\perp|^2$.
\end{lemma}

\begin{proof}
    By~\ref{al:2.3(1)}\eqref{eq:2.3:3}, $S \bullet S = \vdim = T \bullet T$ and $S \bullet
    \id{\R^{\adim}} = \vdim = T \bullet \id{\R^{\adim}}$. Hence $|S - T|^2 = 2 (\vdim - S \bullet T) =
    2 S \bullet (\id{\R^{\adim}} - T) = 2 S \bullet T^\perp$, and symmetrically $|S - T|^2 = 2 T
    \bullet S^\perp$. Finally $S - T = T^\perp - S^\perp$.
\end{proof}

\begin{lemma}[{\cite[8.9(2)]{Allard1972}}]
    \label{al:8.9(2)}
    If $S$ and~$T$ are orthogonal projections of~$\R^{\adim}$ onto linear subspaces of~$\R^{\adim}$,
    then $S \bullet T = |S \circ T|^2$.
\end{lemma}

\begin{proof}
    Since $S$ and~$T$ are self adjoint and idempotent, $|S \circ T|^2 = \trace(T \circ S \circ S
    \circ T) = \trace(T \circ S \circ T) = \trace(S \circ T \circ T) = \trace(S \circ T) = S \bullet
    T$.
\end{proof}

\begin{lemma}[{\cite[8.9(3)]{Allard1972}}]
    \label{al:8.9(3)}
    If $S, T \in \pgrass{\adim}{\vdim}$, then $\|S - T\| = \|S^\perp \circ T\| = \|S \circ
    T^\perp\| = \|T^\perp \circ S\| = \|T \circ S^\perp\| = \|S^\perp - T^\perp\|$.
\end{lemma}

\begin{proof}
    Since $S - T = T^\perp \circ S - T \circ S^\perp$ and the two terms have values in the
    orthogonal subspaces $\im T^\perp$ and $\im T$,
    \begin{displaymath}
        |(S - T)(v)|^2 = |T^\perp(S(v))|^2 + |T(S^\perp(v))|^2 \le \sup \{ \|T^\perp \circ S\|^2,
        \|T \circ S^\perp\|^2 \} (|S(v)|^2 + |S^\perp(v)|^2)
    \end{displaymath}
    for $v \in \R^{\adim}$. Conversely $T^\perp \circ S = T^\perp \circ (S - T)$ and $T \circ
    S^\perp = T \circ (T - S)$. Hence $\|S - T\| = \sup \{ \|T^\perp \circ S\|, \|T \circ
    S^\perp\| \}$. Since $\|\transpose{f}\| = \|f\|$, $\|S \circ T^\perp\| = \|T^\perp \circ S\|$
    and $\|S^\perp \circ T\| = \|T \circ S^\perp\|$. It remains to show $\|T^\perp \circ S\| =
    \|T \circ S^\perp\|$. Let $A = T | \im S$, regarded as a~map of $\im S$ into $\im T$; then
    $\transpose{A} = S | \im T$, and
    \begin{displaymath}
        \|T^\perp \circ S\|^2 = 1 - \inf \{ |A(v)|^2 : v \in \im S, |v| = 1 \} \,, \qquad \|S^\perp
        \circ T\|^2 = 1 - \inf \{ |\transpose{A}(w)|^2 : w \in \im T, |w| = 1 \} \,.
    \end{displaymath}
    The infima are the least eigenvalues of the self adjoint endomorphisms $\transpose{A} \circ A$
    of $\im S$ and $A \circ \transpose{A}$ of $\im T$. Both spaces have dimension~$\vdim$; $A$ maps
    each eigenspace of $\transpose{A} \circ A$ for a~nonzero eigenvalue injectively into the
    eigenspace of $A \circ \transpose{A}$ for the same eigenvalue, and conversely with
    $\transpose{A}$, while $\dim \ker A = \vdim - \dim \im A = \dim \ker \transpose{A}$. Hence the
    two endomorphisms have the same eigenvalues with the same multiplicities, and the infima agree.
\end{proof}

\begin{lemma}[{\cite[8.9(4)]{Allard1972}}]
    \label{al:8.9(4)}
    If $S, T \in \pgrass{\adim}{\vdim}$, $\tau \in \Hom(\R^{\adim},\R^{\adim})$ and $\tau = S^\perp
    \circ \tau \circ S$, then $2 |T \bullet \tau|^2 \le |S - T|^2 |\tau|^2$.
\end{lemma}

\begin{proof}
    Since $T = \transpose{T}$ and $\trace(f \circ g) = \trace(g \circ f)$,
    \begin{displaymath}
        T \bullet \tau = \trace(T \circ S^\perp \circ \tau \circ S) = \trace(S \circ T \circ S^\perp
        \circ \tau) = (S^\perp \circ T \circ S) \bullet \tau \,,
    \end{displaymath}
    hence $|T \bullet \tau| \le |S^\perp \circ T \circ S| \, |\tau|$ by Schwarz's inequality. If
    $e_1, \ldots, e_{\adim}$ is an orthonormal basis of~$\R^{\adim}$ with $e_1, \ldots, e_{\vdim} \in
    \im S$ and $f \in \Hom(\R^{\adim},\R^{\adim})$, then $|f \circ S|^2 = \textsum{i=1}{\vdim}
    |f(e_i)|^2 \le |f|^2$; thus $|S^\perp \circ T \circ S| \le |S^\perp \circ T|$. Finally $|S^\perp
    \circ T|^2 = S^\perp \bullet T = |S - T|^2/2$ by~\ref{al:8.9(2)} and~\ref{al:8.9(1)}.
\end{proof}

\begin{lemma}[{\cite[8.9(5)]{Allard1972}}]
    \label{al:8.9(5)}
    Suppose $S \in \pgrass{\adim}{\vdim}$, $\eta_1, \eta_2 \in \Hom(\im S, \im S^\perp)$, and $S_i
    \in \pgrass{\adim}{\vdim}$ is the projection onto $\{ y + \eta_i(y) : y \in \im S \}$ for $i \in
    \{1,2\}$. Then
    \begin{displaymath}
        \|S_1 - S_2\| \le \|\eta_1 - \eta_2\| \,, \qquad (1 - \|S_1 - S\|^2) \|\eta_1 - \eta_2\|^2
        \le (1 + \|\eta_2\|^2) \|S_1 - S_2\|^2 \,.
    \end{displaymath}
\end{lemma}

\begin{proof}
    Let $H_i(y) = y + \eta_i(y)$ for $y \in \im S$. If $v \in \im S_1$, then $v = H_1(S(v))$, so
    $|S_2^\perp(v)| = \dist(v, \im S_2) \le |H_1(S(v)) - H_2(S(v))| \le \|\eta_1 - \eta_2\| \, |v|$.
    Thus $\|S_2^\perp \circ S_1\| \le \|\eta_1 - \eta_2\|$, and $\|S_1 - S_2\| = \|S_2^\perp \circ
    S_1\|$ by~\ref{al:8.9(3)}. Next, $S_1^\perp \circ H_1 = 0$ and $S_1^\perp \circ H_2 = H_2 - S_1
    \circ H_2 = (S_2 - S_1) \circ H_2$, hence $S_1^\perp \circ (\eta_1 - \eta_2) = (S_1 - S_2) \circ
    H_2$. Since $\eta_1 - \eta_2$ has values in $\im S^\perp$, $\eta_1 - \eta_2 = S_1^\perp \circ
    (\eta_1 - \eta_2) + S_1 \circ S^\perp \circ (\eta_1 - \eta_2)$, with orthogonal terms, and
    $\|S_1 \circ S^\perp\| = \|S_1 - S\|$ by~\ref{al:8.9(3)}, $\|H_2\|^2 \le 1 + \|\eta_2\|^2$.
    Therefore $\|\eta_1 - \eta_2\|^2 \le \|S_1 - S_2\|^2 (1 + \|\eta_2\|^2) + \|S_1 - S\|^2
    \|\eta_1 - \eta_2\|^2$.
\end{proof}

\begin{definition}[{\cite[8.10]{Allard1972}}]
    \label{al:8.10}
    For $T \in \pgrass{\adim}{\vdim}$, $a \in \R^{\adim}$ and $0 < r < \infty$ let
    $\cylind{T}{a}{r} = \R^{\adim} \cap \{ x : |T(x - a)| < r \}$.
\end{definition}

\begin{lemma}[{\cite[8.11]{Allard1972}}]
    \label{al:8.11}
    Suppose $V \in \Var{\vdim}(\R^{\adim})$, $0 < r < \infty$, $0 < \zeta < 1$, $T \in
    \pgrass{\adim}{\vdim}$, $y \in \im T$, $T^\perp[\spt \|V\|]$ is bounded, and $\|V\|
    \cylind{T}{y}{\zeta r} = 0$. Then
    \begin{displaymath}
        [\zeta(\vdim+1) - \vdim] \, \|V\| \cylind{T}{y}{r} \le \textint{\cylind{T}{y}{r} \times
        \pgrass{\adim}{\vdim}}{} \|S - T\|^2 \ud V(x,S) + r \|\delta V\| \cylind{T}{y}{r} \,.
    \end{displaymath}
\end{lemma}

\begin{proof}
    We may assume $y = 0$, by~\ref{al:8.2}, and $\|V\| \cylind{T}{0}{r} < \infty$, since the set
    $\spt \|V\| \cap \cylind{T}{0}{r}$ is bounded. Suppose $\theta : \R \to \R$ is of class~$\infty$,
    $\theta \le 0$, $\theta' \ge 0$, $\spt \theta \subseteq \{ s : s < r \}$, and $\theta$ is
    constant near~$0$. By~\ref{new:cutoff} choose $\chi \in \dspace{\R^{\adim}}{\R}$ with $\chi = 1$
    near the compact set $\spt \|V\| \cap \{ x : |T(x)| \le r \}$, and let $g(x) = \chi(x)
    \theta(|T(x)|) T(x)$. Then $g \in \VF(\R^{\adim})$, $\spt g \subseteq \cylind{T}{0}{r}$, $|g| \le
    r$, and near $\spt \|V\|$, where $T(x) \ne 0$,
    \begin{displaymath}
        \uD g(x) = \theta(|T(x)|) T + \theta'(|T(x)|) \, |T(x)|^{-1} \, \omega_x \cdot T(x) \,,
        \qquad \omega_x(v) = T(x) \bullet v \,.
    \end{displaymath}
    With $u = |T(x)|^{-1} T(x) \in \im T$, \ref{al:2.3(1)}\eqref{eq:2.3:4},\eqref{eq:2.3:2}
    and~\eqref{eq:2.3:3} give
    \begin{multline*}
        \uD g(x) \bullet S = \theta'(|T(x)|) |T(x)| \, |S(u)|^2 + \theta(|T(x)|) T \bullet S \\ =
        \theta'(|T(x)|) |T(x)| (1 - |S^\perp(u)|^2) + \vdim \theta(|T(x)|) - \theta(|T(x)|) T^\perp
        \bullet S \,.
    \end{multline*}
    By~\ref{al:8.9(3)}, $|S^\perp(u)| = |S^\perp(T(u))| \le \|S - T\|$, and $T^\perp \bullet S \ge
    0$ by~\ref{al:8.9(2)}. Hence $\uD g(x) \bullet S \ge \theta'(|T(x)|) |T(x)| + \vdim
    \theta(|T(x)|) - \theta'(|T(x)|) |T(x)| \, \|S - T\|^2$, which also holds where $T(x) = 0$.
    Integrating and using $\delta V(g) \le r \|\delta V\| \cylind{T}{0}{r}$, by~\ref{al:4.2},
    \begin{equation}
        \label{eq:8.11:a}
        \textint{}{} \big( \theta'(|T(x)|) |T(x)| + \vdim \theta(|T(x)|) \big) \ud \|V\| x \le
        \textint{}{} \theta'(|T(x)|) |T(x)| \, \|S - T\|^2 \ud V(x,S) + r \|\delta V\|
        \cylind{T}{0}{r} \,.
    \end{equation}
    Let $\varphi(s) = -1$ for $s < 0$, $\varphi(s) = -1 + s/r$ for $0 \le s < r$, and $\varphi(s) = 0$
    for $s \ge r$. Let $\theta_j$ be obtained by mollifying $[s \mapsto \varphi((1 + 1/j) s)]$ with
    a~mollifier supported in $\{ s : |s| < r/(4j) \}$, see~\cite[4.1.2]{Federer1969}. Then
    $\theta_j$ satisfies the conditions above, $\theta_j \to \varphi$ uniformly, $0 \le \theta_j'
    \le (1 + 1/j) / r$, $\theta_j'(s) \to \varphi'(s)$ for $s \notin \{ 0, r \}$, and $\theta_j(s) =
    \theta_j'(s) = 0$ for $s \ge r$. Applying~\eqref{eq:8.11:a} with $\theta_j$ and Lebesgue's
    bounded convergence theorem~\cite[2.4.9]{Federer1969} on the set $\{ x : \zeta r \le |T(x)| < r
    \}$, which carries $\|V\| \restrict \cylind{T}{0}{r}$, we obtain~\eqref{eq:8.11:a} with $\varphi$
    in place of~$\theta$. For $\zeta r \le s < r$ we have $\varphi'(s) s + \vdim \varphi(s) = (\vdim +
    1) s/r - \vdim \ge \zeta(\vdim+1) - \vdim$ and $0 \le \varphi'(s) s \le 1$, which yields the
    assertion.
\end{proof}

\begin{miniremark}
    \label{mr:8.12}
    Lemma~\ref{al:8.12} approximates a~varifold which is close to a~plane, has density at least~$1$
    and small first variation, by the graph of a~Lipschitzian map; the weight off the graph and the
    measure of the set where the graph carries density less than~$1$ are bounded by the squares of
    the first variation and of the tilt. The points at which the first variation and the tilt are
    small at all scales below~$1$ form a~set~$A$; by~\ref{al:6.2}, $A$ is a~Lipschitzian graph over
    $\im T$, and Kirszbraun's theorem extends it to $\im T$. The complement of~$A$ is estimated
    by Besicovitch's covering theorem, and the part of $\im T$ not covered by the projection of the
    support by~\ref{al:8.11}.
\end{miniremark}

\begin{lemma}[{\cite[8.12]{Allard1972}}]
    \label{al:8.12}
    Corresponding to each $\epsilon$ and~$\mu$ with $0 < \epsilon < 1$ and $1 < \mu < \infty$ there
    are positive numbers $\xi$ and~$P$ with the following property. Suppose
    \begin{enumerate}[label=(\arabic*)]
        \item \label{al:8.12(1)} $V \in \Var{\vdim}(\R^{\adim})$ and $T \in \pgrass{\adim}{\vdim}$;
        \item \label{al:8.12(2)} for $\|V\|$ almost all $x \in \cylind{T}{0}{5}$, the density
            $\density^{\vdim}(\|V\|,x)$ exists, $\density^{\vdim}(\|V\|,x) \ge 1$, $\|V\| \cball{x}{1}
            \le (2 - \epsilon) \unitmeasure{\vdim}$, and $\epsilon \le r^{-\vdim} \|V\| \cball{x}{r}$
            for $0 < r < 1$;
        \item \label{al:8.12(3)} $0 \le \alpha < \infty$ and $\delta V(g) \le \alpha (\textint{}{}
            |g|^2 \ud \|V\|)^{1/2}$ for $g \in \VF(\R^{\adim})$ with $\spt g \subseteq
            \cylind{T}{0}{5}$;
        \item \label{al:8.12(4)} $\|V\| \cylind{T}{0}{1} > 0$;
        \item \label{al:8.12(5)} $\dist(x, \im T) \le \xi$ for $x \in \spt \|V\|$.
    \end{enumerate}
    Then there are $f : \im T \to \im T^\perp$ and $F : \im T \to \R^{\adim}$ such that
    \begin{enumerate}[label=(\arabic*),start=6]
        \item \label{al:8.12(6)} $T(F(y)) = y$ and $T^\perp(F(y)) = f(y)$ for $y \in \im T$;
        \item \label{al:8.12(7)} $\sup \{ |f(y)| : y \in \im T \} \le \sup \{ \dist(x, \im T) : x \in
            \spt \|V\| \}$;
        \item \label{al:8.12(8)} $|f(y) - f(z)| \le (\mu^2 - 1)^{-1/2} |y - z|$ for $y, z \in \im T$;
        \item \label{al:8.12(9)} with $Y = \im T \cap \{ y : \density^{\vdim}(\|V\|,F(y)) \text{
            exists and } \density^{\vdim}(\|V\|,F(y)) \ge 1 \}$,
            \begin{displaymath}
                \|V\|(\cylind{T}{0}{1} \without \im F) + \HM^{\vdim}(\im T \cap \oball{0}{1} \without
                Y) \le P \Big( \alpha^2 + \textint{\cylind{T}{0}{5} \times
                \pgrass{\adim}{\vdim}}{} \|S - T\|^2 \ud V(x,S) \Big) \,.
            \end{displaymath}
    \end{enumerate}
\end{lemma}

\begin{proof}
    Choose $0 < \rho < 1$ and $1 < \lambda < \infty$ with $\lambda (1 - \rho)^{-\vdim} (2 -
    \epsilon) < 2$. Let $\gamma$ be as in~\ref{al:6.2} corresponding to $\lambda$, $\mu$ and $\nu =
    2$; decreasing~$\gamma$, which is permitted by the proof of~\ref{al:6.2}, we assume $\gamma \le
    1/12$. Let $\delta_0 = \inf \{ \rho, \gamma (1 - \rho) \}$, $\xi = \delta_0 / (4 \mu)$,
    \begin{gather*}
        c = 4 \unitmeasure{\vdim} \epsilon^{-1} (2 \vdim + 2)^{\vdim} \,, \qquad P =
        \besicovitch{\adim} \gamma^{-2} (1 + \unitmeasure{\vdim} \epsilon^{-1}) +
        \besicovitch{\vdim} c (1 + 3 \besicovitch{\adim}^{1/2} \gamma^{-1}) \,,
    \end{gather*}
    where $\besicovitch{\vdim}$ is the number of~\ref{al:2.7(1)} with $\adim$ replaced by~$\vdim$.
    Suppose $V$, $T$, $\alpha$ satisfy \ref{al:8.12(1)}--\ref{al:8.12(5)}, and let
    \begin{displaymath}
        \tau = \Big( \textint{\cylind{T}{0}{5} \times \pgrass{\adim}{\vdim}}{} \|S - T\|^2 \ud
        V(x,S) \Big)^{1/2} \,, \qquad s_0 = \sup \{ \dist(x, \im T) : x \in \spt \|V\| \} \,.
    \end{displaymath}
    By~\ref{al:8.12(5)}, the set $\spt \|V\| \cap \{ x : |T(x)| \le 5 \}$ is compact, hence
    $\|V\| \cylind{T}{0}{5} < \infty$.

    \emph{Properties at all points of the support.} Suppose $a \in \spt \|V\| \cap
    \cylind{T}{0}{5}$. The points at which~\ref{al:8.12(2)} holds are dense in $\spt \|V\| \cap
    \cylind{T}{0}{5}$; choose such points $x_j \to a$. If $0 < r < 1$ and $r < r' < 1$, then
    $\cball{x_j}{r} \subseteq \cball{a}{r'}$ for large~$j$, hence $\|V\| \cball{a}{r'} \ge
    \epsilon r^{\vdim}$; letting $r' \to r+$ we obtain $\|V\| \cball{a}{r} \ge \epsilon r^{\vdim}$.
    Similarly $\oball{a}{1 - |x_j - a|} \subseteq \cball{x_j}{1}$ gives $\|V\| \oball{a}{1} \le (2
    - \epsilon) \unitmeasure{\vdim}$. Thus
    \begin{equation}
        \label{eq:8.12:a}
        \epsilon r^{\vdim} \le \|V\| \cball{a}{r} \,, \quad \|V\| \oball{a}{1} \le (2 - \epsilon)
        \unitmeasure{\vdim} \quad \text{for $a \in \spt \|V\| \cap \cylind{T}{0}{5}$ and $0 < r <
        1$} \,.
    \end{equation}
    By~\ref{al:8.12(3)} and Schwarz's inequality, $\|\delta V\|(O) \le \alpha \|V\|(O)^{1/2}$ for
    open $O \subseteq \cylind{T}{0}{5}$; by the outer regularity of the Radon measure~$\|V\|$,
    see~\cite[2.2.5]{Federer1969},
    \begin{equation}
        \label{eq:8.12:b}
        \|\delta V\|(E) \le \alpha \|V\|(E)^{1/2} \quad \text{for $E \subseteq \cylind{T}{0}{5}$} \,;
    \end{equation}
    in particular $\|\delta V\|$ is a~Radon measure on $\cylind{T}{0}{5}$ and $\|\delta V\|
    (\cylind{T}{0}{5} \without \spt \|V\|) = 0$.

    \emph{The set~$A$ and the graph.} Let $A$ be the set of those $a \in \cylind{T}{0}{4} \cap
    \spt \|V\|$ for which $\density^{\vdim}(\|V\|,a)$ exists, $\density^{\vdim}(\|V\|,a) \ge 1$, and
    for $0 < r < 1$
    \begin{gather*}
        \delta V(g) \le \gamma \|V\| \cball{a}{r}^{1/2} \Big( \textint{}{} |g|^2 \ud \|V\|
        \Big)^{1/2} \quad \text{for $g \in \VF(\R^{\adim})$ with $\spt g \subseteq \oball{a}{r}$}
        \,, \\
        \textint{\cball{a}{r} \times \pgrass{\adim}{\vdim}}{} \|S - T\|^2 \ud V(x,S) \le \gamma^2
        \|V\| \cball{a}{r} \,.
    \end{gather*}
    Let $B = \cylind{T}{0}{4} \cap \spt \|V\| \without A$, $C = T[B]$ and $D = \im T \cap
    \oball{0}{1} \without T[\spt \|V\|]$.

    Suppose $a \in A$ and $0 < r < 1$. Since $\textint{}{} |g|^2 \ud \|V\| \le \|V\| \cball{a}{r}
    \sup \im |g|^2$ for $\spt g \subseteq \oball{a}{r}$, we have $\|\delta V\| \oball{a}{r} \le
    \gamma \|V\| \cball{a}{r}$, and, letting $r' \to r+$ in this inequality for $r'$, also
    \begin{equation}
        \label{eq:8.12:c}
        \|\delta V\| \cball{a}{r} \le \gamma \|V\| \cball{a}{r} \quad \text{for $a \in A$ and $0 <
        r < 1$} \,.
    \end{equation}
    By Schwarz's inequality, $\textint{\cball{a}{r} \times \pgrass{\adim}{\vdim}}{} \|S - T\| \ud
    V(x,S) \le \gamma \|V\| \cball{a}{r}$.

    Suppose $a, \tilde a \in A$, $a \ne \tilde a$ and $|a - \tilde a| < \delta_0$. We show $|a -
    \tilde a| > \mu |T^\perp(a - \tilde a)|$. Otherwise~\ref{al:6.2} applies with $W =
    \cylind{T}{0}{5}$, the restriction of~$V$ to~$W$, $Y = \{ a, \tilde a \}$ and $R = 1 - \rho$:
    indeed $\R^{\adim} \cap \{ x : \dist(x,Y) \le 1 - \rho \} \subseteq \cylind{T}{0}{5}$, $\diam Y
    < \gamma (1 - \rho)$, and the two conditions of~\ref{al:6.2} at $y \in Y$ follow
    from~\eqref{eq:8.12:c} and the preceding inequality, because $R \le 1$. Since $|a - \tilde a| <
    \rho$, the set $\{ x : \dist(x,Y) \le 1 - \rho \}$ is contained in $\oball{a}{1}$,
    and~\eqref{eq:8.12:a} gives
    \begin{displaymath}
        2 \unitmeasure{\vdim} \le \unitmeasure{\vdim} \big( \density^{\vdim}(\|V\|,a) +
        \density^{\vdim}(\|V\|,\tilde a) \big) \le \lambda (1 - \rho)^{-\vdim} \|V\| \oball{a}{1} \le
        \lambda (1 - \rho)^{-\vdim} (2 - \epsilon) \unitmeasure{\vdim} < 2 \unitmeasure{\vdim} \,,
    \end{displaymath}
    a~contradiction. If $a, \tilde a \in A$ and $|a - \tilde a| \ge \delta_0$, then $|T^\perp(a -
    \tilde a)| \le 2 \xi < \mu^{-1} \delta_0$ by~\ref{al:8.12(5)}. Hence $|a - \tilde a|^2 > \mu^2
    |T^\perp(a - \tilde a)|^2$, that is
    \begin{displaymath}
        (\mu^2 - 1) |T^\perp(a - \tilde a)|^2 < |T(a - \tilde a)|^2 \quad \text{for $a, \tilde a \in
        A$ with $a \ne \tilde a$} \,.
    \end{displaymath}
    Consequently $T|A$ is univalent, and $f_0 = T^\perp \circ (T|A)^{-1}$ maps $T[A]$ into $\im
    T^\perp \cap \cball{0}{s_0}$ with $\Lip f_0 \le (\mu^2 - 1)^{-1/2}$. By Kirszbraun's
    theorem~\cite[2.10.43]{Federer1969}, applied after identifying $\im T$ and $\im T^\perp$ with
    Euclidean spaces through linear isometries, $f_0$ has an extension $f_1 : \im T \to \im T^\perp$
    with the same Lipschitz constant. Let $\sigma(w) = w$ for $w \in \im T^\perp$ with $|w| \le
    s_0$, and $\sigma(w) = s_0 |w|^{-1} w$ for $w \in \im T^\perp$ with $|w| > s_0$. For $w, w' \in
    \im T^\perp$ we have $(w - \sigma(w)) \bullet (\sigma(w') - \sigma(w)) \le 0$, since either $w =
    \sigma(w)$ or $w - \sigma(w)$ is a~nonnegative multiple of $\sigma(w)$ and $\sigma(w) \bullet
    \sigma(w') \le s_0^2 = |\sigma(w)|^2$; adding this inequality to the one with $w$ and~$w'$
    interchanged yields $|\sigma(w) - \sigma(w')|^2 \le (w - w') \bullet (\sigma(w) - \sigma(w'))$,
    hence $\Lip \sigma \le 1$. Let $f = \sigma \circ f_1$ and $F(y) = y + f(y)$ for $y \in \im T$.
    Then \ref{al:8.12(6)}, \ref{al:8.12(7)} and~\ref{al:8.12(8)} hold, $f | T[A] = f_0$ and
    \begin{equation}
        \label{eq:8.12:10}
        A \subseteq \im F \,.
    \end{equation}

    \emph{Estimate of $\|V\|(B)$.} Let $N$ be the set of those $x \in \cylind{T}{0}{4}$ at which
    the condition on densities in~\ref{al:8.12(2)} fails; then $\|V\|(N) = 0$. To each $b \in B
    \without N$ choose $0 < r_b < 1$ such that either some $g_b \in \VF(\R^{\adim})$ with $\spt g_b
    \subseteq \oball{b}{r_b}$ violates the first condition defining~$A$, or the second condition
    fails with $r = r_b$. By~\ref{al:2.7(1)} there are disjointed subfamilies $\mathcal{C}_1, \ldots,
    \mathcal{C}_{\besicovitch{\adim}}$ of $\{ \cball{b}{r_b} : b \in B \without N \}$ covering $B
    \without N$; each $\mathcal{C}_j$ is countable, since its members are disjoint and have
    positive $\LM^{\adim}$ measure. Suppose $\cball{a_1}{r_1}, \ldots, \cball{a_l}{r_l}$ are
    distinct members of one~$\mathcal{C}_j$ for which the first alternative occurs, with fields
    $g_1, \ldots, g_l$. Since $\delta V(g_i) > 0$, \ref{al:8.12(3)} gives $\textint{}{} |g_i|^2 \ud
    \|V\| > 0$, and we may multiply $g_i$ by a~positive number to achieve $\textint{}{} |g_i|^2 \ud
    \|V\| = \|V\| \cball{a_i}{r_i}$. With $g = \textsum{i=1}{l} g_i$, whose summands have disjoint
    supports in $\cylind{T}{0}{5}$, \ref{al:8.12(3)} yields
    \begin{displaymath}
        \gamma \textsum{i=1}{l} \|V\| \cball{a_i}{r_i} < \delta V(g) \le \alpha \Big(
        \textsum{i=1}{l} \|V\| \cball{a_i}{r_i} \Big)^{1/2} \,, \quad \text{hence} \quad
        \textsum{i=1}{l} \|V\| \cball{a_i}{r_i} < \gamma^{-2} \alpha^2 \,.
    \end{displaymath}
    If the second alternative occurs for distinct members $\cball{a_1}{r_1}, \ldots,
    \cball{a_l}{r_l}$ of~$\mathcal{C}_j$, then $\gamma^2 \textsum{i=1}{l} \|V\| \cball{a_i}{r_i} <
    \tau^2$. Summing over $\mathcal{C}_j$ and then over~$j$,
    \begin{equation}
        \label{eq:8.12:11}
        \|V\|(B) \le \besicovitch{\adim} \gamma^{-2} (\alpha^2 + \tau^2) \,.
    \end{equation}

    \emph{Estimate of $\HM^{\vdim}(C)$.} Choose a~Borel set $\tilde B$ with $B \subseteq \tilde B
    \subseteq \cylind{T}{0}{4}$ and $\|V\|(\tilde B) = \|V\|(B)$, see~\cite[2.2.5]{Federer1969}.
    By~\ref{al:2.8(2)}(c), applied to $\|V\| \restrict \cylind{T}{0}{5}$ and the set
    $\cylind{T}{0}{5} \without \tilde B$, we have $\udensity^{\vdim}(\|V\| \restrict \tilde B, x) =
    \udensity^{\vdim}(\|V\|, x) \ge \epsilon / \unitmeasure{\vdim}$ for $\HM^{\vdim}$ almost all $x
    \in B$, by~\eqref{eq:8.12:a}. Applying~\ref{al:2.8(2)}(b) with $0 < t < \epsilon /
    \unitmeasure{\vdim}$ and letting $t \to \epsilon / \unitmeasure{\vdim}$, and
    using~\cite[2.10.11]{Federer1969} with $\Lip T \le 1$, we obtain
    \begin{equation}
        \label{eq:8.12:12}
        \HM^{\vdim}(C) \le \HM^{\vdim}(B) \le \unitmeasure{\vdim} \epsilon^{-1} \|V\|(B) \le
        \unitmeasure{\vdim} \epsilon^{-1} \besicovitch{\adim} \gamma^{-2} (\alpha^2 + \tau^2) \,.
    \end{equation}

    \emph{Estimate of $\HM^{\vdim}(D)$.} We show
    \begin{equation}
        \label{eq:8.12:13}
        \|\delta V\|(A \cap E) \le \gamma \|V\|(E) \quad \text{for $E \subseteq \cylind{T}{0}{5}$}
        \,.
    \end{equation}
    By outer regularity of~$\|V\|$ it suffices to treat open~$E$. The closed balls $\cball{a}{r}$
    with $a \in A \cap E$, $0 < r < 1$ and $\cball{a}{r} \subseteq E$ form a~relation which is fine
    at each point of $A \cap E$; since the closed balls in~$\R^{\adim}$ form a~$\|\delta V\|
    \restrict \cylind{T}{0}{5}$ Vitali relation by~\cite[2.8.9, 2.8.18]{Federer1969}, there is
    a~countable disjointed subfamily $\{ \cball{a_i}{r_i} : i = 1, 2, \ldots \}$ covering $\|\delta
    V\|$ almost all of $A \cap E$. By~\eqref{eq:8.12:c}, $\|\delta V\|(A \cap E) \le
    \textsum{i}{} \|\delta V\| \cball{a_i}{r_i} \le \gamma \textsum{i}{} \|V\| \cball{a_i}{r_i} \le
    \gamma \|V\|(E)$.

    Suppose $y \in D$ and let $s = \inf \{ |T(x - y)| : x \in \spt \|V\| \}$. By~\ref{al:8.12(4)}
    there is $x \in \spt \|V\|$ with $|T(x)| < 1$, hence $s < 2$; since $\spt \|V\| \cap \{ x :
    |T(x)| \le 3 \}$ is compact and $y \notin T[\spt \|V\|]$, we have $s > 0$ and $|T(x_0 - y)| =
    s$ for some $x_0 \in \spt \|V\|$. Let $\zeta = (2\vdim + 1)/(2 \vdim + 2)$ and $r = s / \zeta$;
    then $\zeta (\vdim + 1) - \vdim = 1/2$, $r < 8/3$, $\|V\| \cylind{T}{y}{\zeta r} = 0$, and
    $\cylind{T}{y}{r} \subseteq \cylind{T}{0}{4}$. By~\ref{al:8.11},
    \begin{displaymath}
        \tfrac12 \|V\| \cylind{T}{y}{r} \le \textint{\cylind{T}{y}{r} \times
        \pgrass{\adim}{\vdim}}{} \|S - T\|^2 \ud V(x,S) + r \|\delta V\| \cylind{T}{y}{r} \,.
    \end{displaymath}
    Since $\cylind{T}{y}{r} \cap \spt \|V\| \subseteq A \cup \tilde B$ and $\|\delta V\|$ vanishes
    off $\spt \|V\|$ in $\cylind{T}{0}{5}$, \eqref{eq:8.12:13} and $3 \gamma \le 1/4$ give $r
    \|\delta V\| \cylind{T}{y}{r} \le \frac14 \|V\| \cylind{T}{y}{r} + 3 \|\delta V\|(\tilde B \cap
    \cylind{T}{y}{r})$, hence
    \begin{displaymath}
        \tfrac14 \|V\| \cylind{T}{y}{r} \le \textint{\cylind{T}{y}{r} \times
        \pgrass{\adim}{\vdim}}{} \|S - T\|^2 \ud V(x,S) + 3 \|\delta V\|(\tilde B \cap
        \cylind{T}{y}{r}) \,.
    \end{displaymath}
    Since $\oball{x_0}{r - s} \subseteq \cylind{T}{y}{r}$ and $r - s = r / (2 \vdim + 2) < 1$,
    \eqref{eq:8.12:a} gives $\|V\| \cylind{T}{y}{r} \ge \epsilon (r / (2 \vdim + 2))^{\vdim}$.
    Therefore
    \begin{displaymath}
        \HM^{\vdim}(D \cap \cball{y}{r}) \le \unitmeasure{\vdim} r^{\vdim} \le c \Big(
        \textint{\cylind{T}{y}{r} \times \pgrass{\adim}{\vdim}}{} \|S - T\|^2 \ud V(x,S) + 3 \|\delta
        V\|(\tilde B \cap \cylind{T}{y}{r}) \Big) \,.
    \end{displaymath}
    Apply~\ref{al:2.7(1)} in $\im T$, identified with $\R^{\vdim}$ through a~linear isometry, to the
    family of the sets $\im T \cap \cball{y}{r}$ with $y \in D$ and $r$ as above. If $\im T \cap
    \cball{y_i}{r_i}$, $i = 1, 2, \ldots$, are disjoint, then so are the cylinders
    $\cylind{T}{y_i}{r_i} \subseteq \cylind{T}{0}{4}$. Hence, by~\eqref{eq:8.12:b}
    and~\eqref{eq:8.12:11},
    \begin{equation}
        \label{eq:8.12:16}
        \begin{multlined}
        \HM^{\vdim}(D) \le \besicovitch{\vdim} c \big( \tau^2 + 3 \|\delta V\|(\tilde B) \big) \le
        \besicovitch{\vdim} c \big( \tau^2 + 3 \alpha \besicovitch{\adim}^{1/2} \gamma^{-1}
        (\alpha^2 + \tau^2)^{1/2} \big) \\ \le \besicovitch{\vdim} c (1 + 3 \besicovitch{\adim}^{1/2}
        \gamma^{-1}) (\alpha^2 + \tau^2) \,.
        \end{multlined}
    \end{equation}

    \emph{Conclusion.} By~\eqref{eq:8.12:10}, $\cylind{T}{0}{1} \cap \spt \|V\| \without \im F
    \subseteq B$, so $\|V\|(\cylind{T}{0}{1} \without \im F) \le \|V\|(B)$. If $y \in \im T \cap
    \oball{0}{1} \without D$, then $y = T(x)$ for some $x \in \spt \|V\| \cap \cylind{T}{0}{1}$;
    either $x \in B$ and $y \in C$, or $x \in A$, in which case $x = F(y)$ by~\eqref{eq:8.12:10}
    and~\ref{al:8.12(6)}, and $y \in Y$. Therefore $\im T \cap \oball{0}{1} \without Y \subseteq C
    \cup D$, and~\eqref{eq:8.12:11}, \eqref{eq:8.12:12}, \eqref{eq:8.12:16} yield~\ref{al:8.12(9)}.
\end{proof}

\begin{remark}
    \label{rem:8.12}
    Allard estimates $\HM^{\vdim}(C) \le \epsilon^{-1} \|V\|(B)$; the factor
    $\unitmeasure{\vdim}$ is needed because upper densities are normalised by
    $\unitmeasure{\vdim}$. His constant~$c$ lacks the factor~$4$ coming from the coefficient
    $\frac14$. The inequality~\eqref{eq:8.12:13}, which Allard derives from~\ref{al:4.3}\eqref{eq:4.3:3},
    is proved here directly by Vitali's covering theorem. The extension of $f_0$ is composed with the
    retraction~$\sigma$ to obtain~\ref{al:8.12(7)}, which Kirszbraun's theorem alone does not give.
    In~\ref{al:8.12(2)} we require the density to exist, as Allard's notation indicates.
\end{remark}

\begin{lemma}[{\cite[8.13]{Allard1972}}]
    \label{al:8.13}
    Suppose $G$ and~$U$ are open subsets of~$\R^{\adim}$, $0 < \delta < \infty$, $\oball{x}{\delta}
    \subseteq G$ for $x \in U$, $T \in \pgrass{\adim}{\vdim}$, $V \in \Var{\vdim}(G)$, $0 \le \alpha <
    \infty$, and
    \begin{displaymath}
        \delta V(g) \le \alpha \Big( \textint{}{} |g|^2 \ud \|V\| \Big)^{1/2} \quad \text{for $g \in
        \VF(G)$} \,.
    \end{displaymath}
    Then
    \begin{displaymath}
        \textint{U \times \pgrass{\adim}{\vdim}}{} \|S - T\|^2 \ud V(x,S) \le \sup \Big\{ \delta^2
        \alpha^2, 9 \delta^{-2} \textint{}{} \dist(x, \im T)^2 \ud \|V\| x \Big\} \,.
    \end{displaymath}
\end{lemma}

\begin{proof}
    Since $U = \tbcup \{ U \cap \oball{0}{j} : j = 1, 2, \ldots \}$, we may assume $U$ is bounded,
    by~\cite[2.4.7]{Federer1969}. Let $\Lambda$ be a~number with $\delta^{-1} < \Lambda < \infty$.
    By~\ref{new:cutoff} with $K = \Clos U$, $\rho = \delta$ and $\lambda = \delta \Lambda$ there is
    $\phi \in \dspace{\R^{\adim}}{\R}$ with $0 \le \phi \le 1$, $\phi = 1$ on~$U$, $|\grad \phi| \le
    \Lambda$, and $\spt \phi \subseteq \{ x : \dist(x, \Clos U) < \delta \} \subseteq G$. Let $g(x)
    = \phi(x)^2 T^\perp(x)$ for $x \in G$; then $g \in \VF(G)$ and $\uD g(x) = \phi(x)^2 T^\perp + 2
    \phi(x) \, \uD \phi(x) \cdot T^\perp(x)$. By~\ref{al:2.3(1)}\eqref{eq:2.3:4},
    \begin{displaymath}
        \uD g(x) \bullet S = \phi(x)^2 T^\perp \bullet S + 2 \phi(x) \grad \phi(x) \bullet
        S(T^\perp(x)) \quad \text{for $(x,S) \in G \times \pgrass{\adim}{\vdim}$} \,.
    \end{displaymath}
    By~\ref{al:8.9(3)} and~\ref{al:8.9(2)}, and since the operator norm does not exceed the
    Euclidean norm, $\|S - T\|^2 = \|T^\perp \circ S\|^2 \le |T^\perp \circ S|^2 = T^\perp \bullet S$,
    and $|S(T^\perp(x))| = |S(T^\perp(T^\perp(x)))| \le \|S - T\| \dist(x, \im T)$. Hence
    \begin{displaymath}
        (\phi(x) \|S - T\|)^2 \le \uD g(x) \bullet S + 2 \Lambda \phi(x) \|S - T\| \dist(x, \im T)
        \quad \text{for $(x,S) \in G \times \pgrass{\adim}{\vdim}$} \,.
    \end{displaymath}
    Let $X = (\textint{}{} (\phi(x) \|S - T\|)^2 \ud V(x,S))^{1/2}$ and $Q = (\textint{}{}
    \dist(x, \im T)^2 \ud \|V\| x)^{1/2}$, and assume $Q < \infty$. Integrating, using
    \ref{al:4.2:lem}, the hypothesis with $|g(x)| \le \dist(x, \im T)$, and Schwarz's inequality, we
    obtain $X^2 \le \alpha Q + 2 \Lambda X Q$. If $X > \delta \alpha$, it follows that $X^2 \le
    (\delta^{-1} + 2 \Lambda) X Q$, hence $X \le (\delta^{-1} + 2 \Lambda) Q$. In either case
    \begin{displaymath}
        \textint{U \times \pgrass{\adim}{\vdim}}{} \|S - T\|^2 \ud V(x,S) \le X^2 \le \sup \{ \delta^2
        \alpha^2, (\delta^{-1} + 2 \Lambda)^2 Q^2 \} \,,
    \end{displaymath}
    and letting $\Lambda \to \delta^{-1}+$ yields the assertion.
\end{proof}

\begin{lemma}[{\cite[8.14]{Allard1972}}]
    \label{al:8.14}
    There is a~positive number $\mathbf{c}_2$ with the following property. Suppose $T \in
    \pgrass{\adim}{\vdim}$, $i : \im T \to \R^{\adim}$ is the inclusion map, $l \in \Hom(\im T, \im
    T^\perp)$, and $L = i + l$, regarded as a~member of $\Hom(\im T, \R^{\adim})$. Then
    \begin{equation}
        \label{eq:8.14:1}
        |\extpower{\vdim} (i + \kappa)|^2 = 1 + \textsum{j=1}{\vdim} |\extpower{j} \kappa|^2 \quad
        \text{for $\kappa \in \Hom(\im T, \im T^\perp)$} \,,
    \end{equation}
    in particular $|\extpower{\vdim} L| \ge 1$, and, with $P$ denoting the orthogonal projection of
    $\R^{\adim}$ onto $\im L$,
    \begin{displaymath}
        \big| [(\theta \circ T) \bullet P] \, |\extpower{\vdim} L| - l \bullet \theta \big| \le
        \mathbf{c}_2 \|\theta\| \sup \{ \|l\|^3, \|l\|^{3 \vdim - 1} \} \quad \text{for $\theta \in
        \Hom(\im T, \im T^\perp)$} \,,
    \end{displaymath}
    where $\theta \circ T$ and $l$, $\theta$ are regarded as members of $\Hom(\R^{\adim},
    \R^{\adim})$ and $\Hom(\im T, \R^{\adim})$, respectively.
\end{lemma}

\begin{proof}
    Let $e_1, \ldots, e_{\adim}$ be an orthonormal basis of~$\R^{\adim}$ with $e_1, \ldots,
    e_{\vdim} \in \im T$. By~\cite[1.7.5]{Federer1969} the $j$-vectors $e_\lambda$ corresponding to
    $\lambda \in \Lambda(\adim,j)$ form an orthonormal basis of $\extpower{j} \R^{\adim}$, and those
    with $\im \lambda \subseteq \{1, \ldots, \vdim\}$ form an orthonormal basis of $\extpower{j}
    \im T$. Let $\kappa \in \Hom(\im T, \im T^\perp)$. Expanding $(e_1 + \kappa(e_1)) \wedge \cdots
    \wedge (e_{\vdim} + \kappa(e_{\vdim}))$ we obtain the sum over the subsets $J$ of $\{1, \ldots,
    \vdim\}$ of the $\vdim$-vectors $\pm e_{I} \wedge (\extpower{|J|} \kappa)(e_J)$, where $I = \{1,
    \ldots, \vdim\} \without J$ and $e_I$, $e_J$ are the corresponding increasing products. Writing
    $(\extpower{|J|} \kappa)(e_J)$ in the basis $e_\lambda$ with $\im \lambda \subseteq \{\vdim+1,
    \ldots, \adim\}$, we see that the summands corresponding to different~$J$ are orthogonal and
    that $|e_I \wedge (\extpower{|J|}\kappa)(e_J)| = |(\extpower{|J|} \kappa)(e_J)|$. By~\ref{al:2.4},
    \begin{displaymath}
        |\extpower{\vdim} (i + \kappa)|^2 = \textsum{J}{} |(\extpower{|J|} \kappa)(e_J)|^2 = 1 +
        \textsum{j=1}{\vdim} |\extpower{j} \kappa|^2 \,,
    \end{displaymath}
    which is~\eqref{eq:8.14:1}.

    Next suppose $u_1, \ldots, u_{\vdim}$ is an orthonormal basis of $\im L$, $\eta = L(e_1) \wedge
    \cdots \wedge L(e_{\vdim}) = |\eta| u_1 \wedge \cdots \wedge u_{\vdim}$, and $f \in
    \Hom(\R^{\adim}, \R^{\adim})$. Then $|\extpower{\vdim} ((\id{\R^{\adim}} + t f) \circ L)| =
    |\extpower{\vdim}(\id{\R^{\adim}} + t f)(\eta)|$ for $t \in \R$, by~\ref{al:2.4}, and
    \begin{multline*}
        \frac{\ud}{\ud t} |\extpower{\vdim}(\id{\R^{\adim}} + t f)(\eta)| \Big|_{t=0} = |\eta|
        \textsum{j=1}{\vdim} (u_1 \wedge \cdots \wedge f(u_j) \wedge \cdots \wedge u_{\vdim}) \bullet
        (u_1 \wedge \cdots \wedge u_{\vdim}) \\ = |\eta| \textsum{j=1}{\vdim} f(u_j) \bullet u_j = |\eta|
        \, f \bullet P \,;
    \end{multline*}
    here we split $f(u_j)$ into its components in $\im P$ and $\ker P$ and used~\cite[1.7.5]{Federer1969}
    and~\ref{al:2.3(1)}\eqref{eq:2.3:5}. Apply this with $f = \theta \circ T$; since $(\theta \circ
    T) \circ L = \theta$, we have $(\id{\R^{\adim}} + t f) \circ L = i + l + t \theta$, and
    by~\eqref{eq:8.14:1}
    \begin{displaymath}
        [(\theta \circ T) \bullet P] \, |\extpower{\vdim} L| = \frac{\ud}{\ud t} \Big( 1 +
        \textsum{j=1}{\vdim} |\extpower{j}(l + t \theta)|^2 \Big)^{1/2} \Big|_{t=0} = (1 + X)^{-1/2}
        \textsum{j=1}{\vdim} \extpower{j} l \bullet \Psi_j(l,\theta) \,,
    \end{displaymath}
    where $X = \textsum{j=1}{\vdim} |\extpower{j} l|^2$ and $\Psi_j(l,\theta) = \frac{\ud}{\ud t}
    \extpower{j}(l + t \theta) |_{t=0}$. Note $\Psi_1(l,\theta) = \theta$. Since $\extpower{j}(s l)
    = s^j \extpower{j} l$ and $\Psi_j(s l, \theta) = s^{j-1} \Psi_j(l,\theta)$ for $s \in \R$, and
    $\Psi_j(l,\theta)$ is linear in~$\theta$, compactness of the unit spheres yields numbers $a_j$
    depending only on $\adim$ and~$\vdim$ with $|\extpower{j} l| \le a_j \|l\|^j$ and
    $|\Psi_j(l,\theta)| \le a_j \|l\|^{j-1} \|\theta\|$. Hence
    \begin{displaymath}
        \big| [(\theta \circ T) \bullet P] \, |\extpower{\vdim} L| - l \bullet \theta \big| \le
        \textsum{j=2}{\vdim} a_j^2 \|l\|^{2j-1} \|\theta\| + |(1 + X)^{-1/2} - 1|
        \textsum{j=1}{\vdim} a_j^2 \|l\|^{2j-1} \|\theta\| \,.
    \end{displaymath}
    Since $|(1 + X)^{-1/2} - 1| \le \inf \{ 1, X \}$ and $X \le \textsum{j=1}{\vdim} a_j^2
    \|l\|^{2j}$, the right side does not exceed $\mathbf{c}_2 \|\theta\| \|l\|^3$ if $\|l\| \le 1$
    and $\mathbf{c}_2 \|\theta\| \sup \{ \|l\|^3, \|l\|^{2 \vdim - 1} \}$ if $\|l\| \ge 1$, for
    a~suitable $\mathbf{c}_2$ depending only on $\adim$ and~$\vdim$. Since $2 \vdim - 1 \le 3 \vdim -
    1$, the assertion follows.
\end{proof}

\begin{remark}
    \label{rem:8.14}
    Allard states the one-sided inequality; the proof gives the bound for the absolute value, and in
    fact with $\|l\|^{2\vdim - 1}$ in place of $\|l\|^{3 \vdim - 1}$. The identity~\eqref{eq:8.14:1}
    is Allard's first formula in the proof.
\end{remark}

\begin{miniremark}
    \label{mr:harmonic}
    Lemmas~\ref{new:ibp}--\ref{new:rellich} supply the facts on harmonic functions which Allard
    quotes from~\cite[5.2.5, 5.2.6]{Federer1969} and the compactness argument which he calls
    ``well-known reasoning'' in the proof of~\ref{al:8.16}. Throughout them, $Z$ is a~finite
    dimensional inner product space, $e_1, \ldots, e_{\vdim}$ is the standard orthonormal basis
    of~$\R^{\vdim}$, and, for an open subset $\Omega$ of~$\R^{\vdim}$, a~map $u : \Omega \to Z$ is
    \emph{harmonic} if and only if $u$ is of class~$2$ and $\textsum{j=1}{\vdim} \uD^2 u(x)(e_j,e_j)
    = 0$ for $x \in \Omega$. A~map of an open subset of $\im T$, where $T \in \pgrass{\adim}{\vdim}$,
    is harmonic if its composition with a~linear isometry of $\R^{\vdim}$ onto $\im T$ is harmonic;
    this does not depend on the isometry, since $\textsum{j=1}{\vdim} \beta(\iota(e_j),\iota(e_j))$
    is the same for all orthogonal maps $\iota$ of~$\R^{\vdim}$ and bilinear~$\beta$. We fix
    $\chi(t) = \exp(-1/t)$ for $0 < t < \infty$, $\chi(t) = 0$ for $-\infty < t \le 0$,
    \begin{displaymath}
        \Psi_0(s) = \chi(1/16 - s) / \big( \chi(1/16 - s) + \chi(s - 1/64) \big) \quad \text{for $s
        \in \R$} \,,
    \end{displaymath}
    and $\psi(z) = c_0 \Psi_0(|z|^2)$ for $z \in \R^{\vdim}$, with $c_0$ chosen so that $\textint{}{}
    \psi \ud \LM^{\vdim} = 1$. Then $\Psi_0$ is of class~$\infty$, $\Psi_0(s) = 1$ for $s \le 1/64$,
    $\Psi_0(s) = 0$ for $s \ge 1/16$, $\psi \ge 0$, $\psi \in \dspace{\R^{\vdim}}{\R}$ and $\spt \psi
    \subseteq \oball{0}{1/4}$. For $0 < \epsilon < \infty$ let $\psi_\epsilon(z) = \epsilon^{-\vdim}
    \psi(z/\epsilon)$. Finally let $\Psi_1(t) = \chi(1 - t) / (\chi(1 - t) + \chi(t))$ for $t \in
    \R$; then $\Psi_1$ is of class~$\infty$, $0 \le \Psi_1 \le 1$, $\Psi_1(t) = 1$ for $t \le 0$ and
    $\Psi_1(t) = 0$ for $t \ge 1$.
\end{miniremark}

\begin{lemma}[new]
    \label{new:ibp}
    Suppose $\Omega$ is an open subset of~$\R^{\vdim}$, $f : \Omega \to \R$ is Lipschitzian on each
    compact subset of~$\Omega$, $\phi \in \dspace{\Omega}{\R}$ and $v \in \R^{\vdim}$. Then
    \begin{displaymath}
        \textint{\Omega}{} \uD f(x)(v) \, \phi(x) \ud \LM^{\vdim} x = - \textint{\Omega}{} f(x) \,
        \uD \phi(x)(v) \ud \LM^{\vdim} x \,,
    \end{displaymath}
    where $\uD f(x)$ exists for $\LM^{\vdim}$ almost all $x \in \Omega$ by Rademacher's
    theorem~\cite[3.1.6]{Federer1969}.
\end{lemma}

\begin{proof}
    Let $F(x) = f(x) \phi(x)$ for $x \in \Omega$ and $F(x) = 0$ for $x \in \R^{\vdim} \without
    \Omega$. Since $f$ is Lipschitzian and bounded on the compact set $\{ x : \dist(x, \spt \phi)
    \le \delta \}$, where $2 \delta = \dist(\spt \phi, \R^{\vdim} \without \Omega)$, the function $F$
    is Lipschitzian. By~\cite[3.1.6]{Federer1969} the set $N$ of points at which $F$ is not
    differentiable is an $\LM^{\vdim}$ null set, and $\uD F(x) = \phi(x) \uD f(x) + f(x) \uD
    \phi(x)$ for $\LM^{\vdim}$ almost all $x \in \Omega$, while $\uD F(x) = 0$ for $x \in
    \R^{\vdim} \without (N \cup \spt \phi)$. By Fubini's theorem~\cite[2.6.2]{Federer1969}, for
    $\LM^{\vdim}$ almost all~$x$ the point $x + t v$ lies outside~$N$ for $\LM^1$ almost all $0 \le t
    \le 1$; for such~$x$, the Lipschitzian function $[t \mapsto F(x + t v)]$ has derivative $\uD
    F(x + t v)(v)$ at these~$t$, and~\cite[2.9.20]{Federer1969} gives $F(x + v) - F(x) =
    \textint{0}{1} \uD F(x + t v)(v) \ud \LM^1 t$. Integrating over~$x$, using Fubini's theorem and
    the invariance of~$\LM^{\vdim}$ under translations, we obtain
    \begin{displaymath}
        0 = \textint{}{} (F(x + v) - F(x)) \ud \LM^{\vdim} x = \textint{0}{1} \textint{}{} \uD F(x +
        t v)(v) \ud \LM^{\vdim} x \ud \LM^1 t = \textint{}{} \uD F(x)(v) \ud \LM^{\vdim} x \,,
    \end{displaymath}
    which is the assertion.
\end{proof}

\begin{lemma}[new]
    \label{new:meanvalue}
    Suppose $x \in \R^{\vdim}$, $0 < R < \infty$, $u : \oball{x}{R} \to Z$ is harmonic, $\Psi : \R
    \to \R$ is of class~$\infty$ and constant on a~neighbourhood of~$0$, $\varphi(z) = \Psi(|z|^2)$
    for $z \in \R^{\vdim}$, and $\varphi \in \dspace{\R^{\vdim}}{\R}$ with $\spt \varphi \subseteq
    \oball{0}{R}$. Then
    \begin{displaymath}
        \textint{}{} u(x - z) \varphi(z) \ud \LM^{\vdim} z = u(x) \textint{}{} \varphi \ud
        \LM^{\vdim} \,.
    \end{displaymath}
\end{lemma}

\begin{proof}
    Replacing $u$ by $[\oball{0}{R} \ni z \mapsto u(x - z)]$, which is harmonic, we assume $x = 0$
    and prove $\textint{}{} u \varphi \ud \LM^{\vdim} = u(0) \textint{}{} \varphi \ud \LM^{\vdim}$.
    For $0 < t \le 1$ let $\varphi_t(z) = t^{-\vdim} \varphi(z/t)$; then $\varphi_t$ has the same
    form as~$\varphi$, $\spt \varphi_t \subseteq \oball{0}{R}$ and $\textint{}{} \varphi_t \ud
    \LM^{\vdim} = \textint{}{} \varphi \ud \LM^{\vdim}$. Fix~$t$. Let $\omega(s) = \Psi(s^2) -
    t^{-\vdim} \Psi(s^2/t^2)$ for $s \in \R$, so that $w = \varphi - \varphi_t$ satisfies $w(z) =
    \omega(|z|)$. Choose $0 < s_1 < r_1 < R$ such that $\omega$ is constant on $\{ s : |s| \le s_1
    \}$ and $\spt w \subseteq \oball{0}{r_1}$. Define
    \begin{displaymath}
        \beta(s) = s^{1-\vdim} \textint{0}{s} \rho^{\vdim-1} \omega(\rho) \ud \LM^1 \rho \,, \qquad
        \Phi(s) = - \textint{s}{\infty} \beta \ud \LM^1 \quad \text{for $0 < s < \infty$} \,.
    \end{displaymath}
    By~\cite[3.2.13]{Federer1969}, $\vdim \unitmeasure{\vdim} \textint{0}{s} \rho^{\vdim-1}
    \omega(\rho) \ud \LM^1 \rho = \textint{\oball{0}{s}}{} w \ud \LM^{\vdim}$, which vanishes for $s
    \ge r_1$; hence $\beta(s) = 0$ and $\Phi(s) = 0$ for $s \ge r_1$. For $0 < s \le s_1$ we have
    $\beta(s) = \omega(0) s/\vdim$, hence $\Phi(s) = \omega(0) s^2/(2 \vdim) + \Phi(s_1) -
    \omega(0) s_1^2/(2\vdim)$. Therefore $\phi(z) = \Phi(|z|)$ for $z \ne 0$ and $\phi(0) =
    \lim_{s \to 0+} \Phi(s)$ define $\phi \in \dspace{\R^{\vdim}}{\R}$ with $\spt \phi \subseteq
    \oball{0}{R}$, and for $z \ne 0$, $s = |z|$,
    \begin{displaymath}
        \textsum{j=1}{\vdim} \uD^2 \phi(z)(e_j,e_j) = \Phi''(s) + (\vdim - 1) s^{-1} \Phi'(s) =
        \beta'(s) + (\vdim - 1) s^{-1} \beta(s) = \omega(s) = w(z) \,,
    \end{displaymath}
    because $\beta'(s) = (1 - \vdim) s^{-1} \beta(s) + \omega(s)$; by continuity this holds also for
    $z = 0$. Applying~\ref{new:ibp} twice to the components of~$u$ and of its first partial
    derivatives, which are of class~$1$, we obtain
    \begin{displaymath}
        \textint{}{} u \, w \ud \LM^{\vdim} = \textsum{j=1}{\vdim} \textint{}{} u \, \uD^2 \phi(\cdot)
        (e_j,e_j) \ud \LM^{\vdim} = \textint{}{} \Big( \textsum{j=1}{\vdim} \uD^2 u(\cdot)(e_j,e_j)
        \Big) \phi \ud \LM^{\vdim} = 0 \,.
    \end{displaymath}
    Thus $\textint{}{} u \varphi \ud \LM^{\vdim} = \textint{}{} u \varphi_t \ud \LM^{\vdim} =
    \textint{}{} u(t z) \varphi(z) \ud \LM^{\vdim} z$ for $0 < t \le 1$. Since $u$ is uniformly
    continuous on $\spt \varphi$, the last integral tends to $u(0) \textint{}{} \varphi \ud
    \LM^{\vdim}$ as $t \to 0+$.
\end{proof}

\begin{lemma}[{\cite[8.15]{Allard1972}}]
    \label{al:8.15}
    There is a~positive number $\mathbf{c}_3$, depending only on~$\vdim$, with the following
    property. If $T \in \pgrass{\adim}{\vdim}$ and $h : \im T \cap \oball{0}{1} \to \im T^\perp$ is
    harmonic, then, with $I = (\textint{\im T \cap \oball{0}{1}}{} |h|^2 \ud \HM^{\vdim})^{1/2}$,
    \begin{gather*}
        \sup \{ |h(0)|, \|\uD h(0)\| \} \le \mathbf{c}_3 I \,, \qquad |h(y) - h(0) - \uD h(0)(y)| \le
        \mathbf{c}_3 I |y|^2 \quad \text{for $y \in \im T$ with $2 |y| < 1$} \,.
    \end{gather*}
\end{lemma}

\begin{proof}
    Let $\iota : \R^{\vdim} \to \im T$ be a~linear isometry and $u = h \circ \iota$; then $u$ is
    harmonic on $\oball{0}{1}$, and $I = (\textint{\oball{0}{1}}{} |u|^2 \ud
    \LM^{\vdim})^{1/2}$ by~\cite[2.10.35]{Federer1969}, because $\HM^{\vdim}$ is invariant under
    isometries. Let $\psi$ be as in~\ref{mr:harmonic} and $m_j = \sup \{ \|\uD^j \psi(z)\| : z \in
    \R^{\vdim} \}$ for $j \in \{0, 1, 2\}$. If $|x| < 3/4$, then $\oball{x}{1/4} \subseteq
    \oball{0}{1}$, and~\ref{new:meanvalue} gives $u(x) = \textint{}{} u(z) \psi(x - z) \ud
    \LM^{\vdim} z$. Differentiating under the integral sign, as in~\cite[4.1.2]{Federer1969}, and
    using Schwarz's inequality,
    \begin{displaymath}
        \|\uD^j u(x)\| \le m_j \textint{\oball{0}{1}}{} |u| \ud \LM^{\vdim} \le m_j
        \unitmeasure{\vdim}^{1/2} I \quad \text{for $|x| < 3/4$ and $j \in \{0, 1, 2\}$} \,.
    \end{displaymath}
    If $|y| < 1/2$, applying~\cite[2.9.20]{Federer1969} twice to $[t \mapsto u(t y)]$ gives $u(y) -
    u(0) - \uD u(0)(y) = \textint{0}{1} (1 - t) \uD^2 u(t y)(y,y) \ud \LM^1 t$, whose norm does not
    exceed $\frac12 m_2 \unitmeasure{\vdim}^{1/2} I |y|^2$. Since $\uD h(0) = \uD u(0) \circ
    \iota^{-1}$, the assertion holds with $\mathbf{c}_3 = \unitmeasure{\vdim}^{1/2} \sup \{ m_0,
    m_1, m_2 \}$.
\end{proof}

\begin{remark}
    \label{rem:8.15}
    Allard refers to~\cite[5.2.5]{Federer1969}, which lies beyond the range of chapters~1--3
    and~4.1.1--4.1.27; the proof above uses only~\ref{new:ibp} and~\ref{new:meanvalue}.
\end{remark}

\begin{lemma}[new]
    \label{new:weyl}
    Suppose $\Omega$ is an open subset of~$\R^{\vdim}$, $h : \Omega \to Z$ is $\LM^{\vdim}
    \restrict \Omega$ measurable, $\textint{K}{} |h| \ud \LM^{\vdim} < \infty$ for every compact
    subset $K$ of~$\Omega$, and
    \begin{displaymath}
        \textint{\Omega}{} h(z) \bullet \textsum{j=1}{\vdim} \uD^2 \phi(z)(e_j,e_j) \ud \LM^{\vdim} z
        = 0 \quad \text{for $\phi \in \dspace{\Omega}{Z}$} \,.
    \end{displaymath}
    Then there is a~harmonic map $H : \Omega \to Z$ with $h(z) = H(z)$ for $\LM^{\vdim}$ almost all
    $z \in \Omega$.
\end{lemma}

\begin{proof}
    For $0 < \epsilon < \infty$ let $\Omega_\epsilon = \Omega \cap \{ x : \oball{x}{\epsilon}
    \subseteq \Omega \}$, and let $h_\epsilon(x) = \textint{}{} h(z) \psi_\epsilon(x - z) \ud
    \LM^{\vdim} z$ for $x \in \Omega_\epsilon$, with $\psi_\epsilon$ as in~\ref{mr:harmonic}. By~\cite[4.1.2]{Federer1969}, $h_\epsilon$ is
    of class~$\infty$ and its derivatives are obtained by differentiating $\psi_\epsilon$. For $x \in
    \Omega_\epsilon$ and $w \in Z$ the map $\phi(z) = \psi_\epsilon(x - z) w$ belongs to
    $\dspace{\Omega}{Z}$, and the hypothesis gives $w \bullet \textsum{j=1}{\vdim} \uD^2
    h_\epsilon(x)(e_j,e_j) = 0$; thus $h_\epsilon$ is harmonic on~$\Omega_\epsilon$. If $x \in
    \Omega_{\epsilon + \delta}$, then $\oball{x}{\delta} \subseteq \Omega_\epsilon$ and
    $\oball{x}{\epsilon} \subseteq \Omega_\delta$, so~\ref{new:meanvalue} and Fubini's
    theorem~\cite[2.6.2]{Federer1969} yield
    \begin{multline*}
        h_\epsilon(x) = \textint{}{} h_\epsilon(x - y) \psi_\delta(y) \ud \LM^{\vdim} y =
        \textint{}{} \textint{}{} h(z) \psi_\epsilon(x - y - z) \psi_\delta(y) \ud \LM^{\vdim} z \ud
        \LM^{\vdim} y \\ = \textint{}{} h_\delta(x - y') \psi_\epsilon(y') \ud \LM^{\vdim} y' =
        h_\delta(x) \,.
    \end{multline*}
    Hence $H(x) = h_\epsilon(x)$ whenever $x \in \Omega_{2 \epsilon}$ defines a~harmonic map $H :
    \Omega \to Z$. Applying~\cite[2.8.9, 2.8.18, 2.9.9]{Federer1969} to the products of~$h$ with the
    characteristic functions of the compact sets $\Omega_{1/m} \cap \cball{0}{m}$, we infer that for
    $\LM^{\vdim}$ almost all $x \in \Omega$ $\lim_{\epsilon \to 0+} \epsilon^{-\vdim} \textint{\cball{x}{\epsilon}}{} |h(z) -
    h(x)| \ud \LM^{\vdim} z = 0$; at such~$x$, $|h_\epsilon(x) - h(x)| \le \sup \im \psi \cdot
    \epsilon^{-\vdim} \textint{\cball{x}{\epsilon}}{} |h(z) - h(x)| \ud \LM^{\vdim} z \to 0$, since
    $\textint{}{} \psi_\epsilon \ud \LM^{\vdim} = 1$, whence $h(x) = H(x)$.
\end{proof}

\begin{remark}
    \label{rem:weyl}
    This replaces Allard's reference ``[FE 2.5.6]'', a~misprint for~\cite[5.2.6]{Federer1969},
    which lies beyond the range of chapters~1--3 and~4.1.1--4.1.27.
\end{remark}

\begin{lemma}[new]
    \label{new:rellich}
    Suppose $0 < \rho < R < \infty$, $g_1, g_2, \ldots$ are Lipschitzian maps of $\oball{0}{R}
    \subseteq \R^{\vdim}$ into~$Z$, and
    \begin{displaymath}
        \sup \Big\{ \textint{\oball{0}{R}}{} \big( |g_i|^2 + \|\uD g_i\|^2 \big) \ud \LM^{\vdim} :
        i = 1, 2, \ldots \Big\} = M < \infty \,.
    \end{displaymath}
    Then there are $i(1) < i(2) < \cdots$ and an $\LM^{\vdim} \restrict \oball{0}{\rho}$
    measurable $h : \oball{0}{\rho} \to Z$ with $\lim_{j \to \infty} \textint{\oball{0}{\rho}}{}
    |g_{i(j)} - h|^2 \ud \LM^{\vdim} = 0$.
\end{lemma}

\begin{proof}
    Suppose $0 < \epsilon < R - \rho$ and $g : \oball{0}{R} \to Z$ is Lipschitzian. By Rademacher's
    theorem and Fubini's theorem, as in the proof of~\ref{new:ibp}, for $z \in
    \oball{0}{\epsilon}$ and $\LM^{\vdim}$ almost all $x \in \oball{0}{\rho}$ we have $g(x) - g(x -
    z) = \textint{0}{1} \uD g(x - t z)(z) \ud \LM^1 t$. Since $\psi_\epsilon \ge 0$, $\textint{}{}
    \psi_\epsilon \ud \LM^{\vdim} = 1$ and $\spt \psi_\epsilon \subseteq \oball{0}{\epsilon}$,
    Schwarz's inequality with respect to the measure $\psi_\epsilon(z) \ud \LM^{\vdim} z \ud \LM^1
    t$ gives
    \begin{displaymath}
        \Big| g(x) - \textint{}{} g(x - z) \psi_\epsilon(z) \ud \LM^{\vdim} z \Big|^2 \le
        \epsilon^2 \textint{}{} \textint{0}{1} \|\uD g(x - t z)\|^2 \psi_\epsilon(z) \ud \LM^1 t \ud
        \LM^{\vdim} z
    \end{displaymath}
    for $\LM^{\vdim}$ almost all $x \in \oball{0}{\rho}$, and integrating over $\oball{0}{\rho}$,
    by Fubini's theorem,
    \begin{equation}
        \label{eq:rellich:a}
        \textint{\oball{0}{\rho}}{} |g - \psi_\epsilon \ast g|^2 \ud \LM^{\vdim} \le \epsilon^2
        \textint{\oball{0}{R}}{} \|\uD g\|^2 \ud \LM^{\vdim} \,, \quad \text{where} \quad
        (\psi_\epsilon \ast g)(x) = \textint{}{} g(x - z) \psi_\epsilon(z) \ud \LM^{\vdim} z \,.
    \end{equation}
    For fixed~$\epsilon$, the maps $\psi_\epsilon \ast g_i$ restricted to $\cball{0}{\rho}$ satisfy
    $|\psi_\epsilon \ast g_i| \le \sup \im \psi_\epsilon \cdot \unitmeasure{\vdim}^{1/2} R^{\vdim/2} M^{1/2}$ and
    $\Lip (\psi_\epsilon \ast g_i) \le \sup \im \|\uD \psi_\epsilon\| \cdot \unitmeasure{\vdim}^{1/2}
    R^{\vdim/2} M^{1/2}$, by~\cite[4.1.2]{Federer1969} and Schwarz's inequality. By~\cite[2.10.21]{Federer1969}
    and a~diagonal argument there are $i(1) < i(2) < \cdots$ such that, for every positive
    integer~$m$ with $1/m < R - \rho$, the sequence $\psi_{1/m} \ast g_{i(j)}$ converges uniformly on
    $\cball{0}{\rho}$ as $j \to \infty$. By~\eqref{eq:rellich:a}, with $\|\cdot\|_2$ denoting the
    norm of $\Lp{2}(\LM^{\vdim} \restrict \oball{0}{\rho}, Z)$,
    \begin{multline*}
        \|g_{i(j)} - g_{i(l)}\|_2 \le 2 M^{1/2} / m \\ + (\unitmeasure{\vdim} \rho^{\vdim})^{1/2} \sup
        \{ |(\psi_{1/m} \ast g_{i(j)})(x) - (\psi_{1/m} \ast g_{i(l)})(x)| : |x| \le \rho \}
        \quad \text{for $1/m < R - \rho$} \,,
    \end{multline*}
    so $g_{i(1)}, g_{i(2)}, \ldots$ is a~Cauchy sequence, which converges to some~$h$ by the
    completeness of the Lebesgue space $\Lp{2}$, see~\cite[2.4.12]{Federer1969}.
\end{proof}

\begin{miniremark}
    \label{mr:8.16}
    Theorem~\ref{al:8.16} is the decay of the height excess: if the excess~$\mu$ and the first
    variation are small, then, after the plane is tilted by an amount controlled by~$\mu$, the
    excess at the smaller scale $\theta r$ improves by the factor $\theta^{1-\vdim/p}$. The proof is
    by contradiction. A~sequence of counterexamples with excess $\mu_i \to 0$ is approximated
    by~\ref{al:8.12} by Lipschitzian graphs of maps~$f_i$; the maps $\mu_i^{-1} f_i$ converge
    in~$\Lp{2}$ by~\ref{new:rellich}, and the limit is harmonic, by the first variation formula
    and~\ref{al:8.14}, \ref{new:ibp}, \ref{new:weyl}. The estimates~\ref{al:8.15} for the harmonic
    limit then yield planes violating the assumption. This is the ``blowing up'' technique of
    Almgren~\cite[6.2, 6.4]{Almgren1968}.
\end{miniremark}

\begin{theorem}[{\cite[8.16]{Allard1972}}]
    \label{al:8.16}
    Corresponding to every $\epsilon$ with $0 < \epsilon < 1$ there are $\Delta$, $\theta$ and~$M$
    with $0 < \Delta < 1$, $0 < 14 \theta < 1$ and $1 < M < \infty$, and with the following
    property. If $V$, $a$, $r$, $T$, $U$, $A$, $\mu$ and~$\alpha$ satisfy
    \begin{enumerate}[label=(\arabic*)]
        \item \label{al:8.16(1)} $V \in \Var{\vdim}(\R^{\adim})$, $a \in \spt \|V\|$, $0 < r <
            \infty$;
        \item \label{al:8.16(2)} $T, U \in \pgrass{\adim}{\vdim}$, $\|T - U\| \le \Delta$, $A \in
            \affgrass{\adim}{\vdim}$, and $A$ is parallel to~$T$;
        \item \label{al:8.16(3)} $\mu = (r^{-\vdim-2} \textint{\cylind{U}{a}{r}}{} \dist(x,A)^2
            \ud \|V\| x)^{1/2} \le \Delta$;
        \item \label{al:8.16(4)} $0 \le \alpha < \infty$ and $\delta V(g) \le \alpha (\textint{}{}
            |g|^q \ud \|V\|)^{1/q}$ for $g \in \VF(\R^{\adim})$ with $\spt g \subseteq
            \cylind{U}{a}{r}$;
        \item \label{al:8.16(5)} for $\|V\|$ almost all $x \in \cylind{U}{a}{r}$, $\epsilon \le
            s^{-\vdim} \|V\| \cball{x}{s} \le (2 - \epsilon) \unitmeasure{\vdim}$ for $0 < s <
            r$;
        \item \label{al:8.16(6)} for $\|V\|$ almost all $x \in \cylind{U}{a}{r}$, the density
            $\density^{\vdim}(\|V\|,x)$ exists and $1 \le \density^{\vdim}(\|V\|,x) \le 1 + \Delta$;
    \end{enumerate}
    then there are $\tilde T$ and $\tilde A$ such that
    \begin{enumerate}[label=(\arabic*),start=7]
        \item \label{al:8.16(7)} $\tilde T \in \pgrass{\adim}{\vdim}$, $\tilde A \in
            \affgrass{\adim}{\vdim}$, and $\tilde A$ is parallel to~$\tilde T$;
        \item \label{al:8.16(8)} $\|\tilde T - T\| \le M \mu$;
        \item \label{al:8.16(9)} $((\theta r)^{-\vdim-2} \textint{\cylind{U}{a}{\theta r}}{}
            \dist(x, \tilde A)^2 \ud \|V\| x)^{1/2} \le \theta^{1 - \vdim/p} \sup \{ \mu, M r^{1 -
            \vdim/p} \alpha \}$.
    \end{enumerate}
\end{theorem}

\begin{proof}
    Let $\mathbf{c}_2$ and~$\mathbf{c}_3$ be as in~\ref{al:8.14} and~\ref{al:8.15}. Since $2 (1 -
    \vdim/p) < 2$, we may choose $\theta$ with $0 < 14 \theta < 1$ and
    \begin{equation}
        \label{eq:8.16:10}
        2^{\vdim+1} \mathbf{c}_3^2 \big( \vdim \unitmeasure{\vdim} / (\vdim + 4) \big) 7^{\vdim + 4}
        \theta^2 < \theta^{2 (1 - \vdim/p)} \,.
    \end{equation}
    Suppose that no $\Delta$ and~$M$ have the stated property with this~$\theta$. Let $\Delta_i =
    (i+4)^{-1}$ and choose numbers $M_i \ge 2$ with
    \begin{equation}
        \label{eq:8.16:12}
        \mathbf{c}_3 7^{(\vdim+2)/2} \le M_i \,, \qquad \theta^{-(\vdim+2)/2} \Delta_i^{-1} \kappa
        \le M_i (7 \theta)^{1 - \vdim/p} \,, \quad \text{where} \quad \kappa = \big( (2 - \epsilon)
        \unitmeasure{\vdim} 7^{\vdim} \big)^{1/q - 1/2} \,.
    \end{equation}
    For each $i$ there are $V_i$, $a_i$, $r_i$, $T_i$, $U_i$, $A_i$, $\mu_i$, $\alpha_i$
    satisfying~\ref{al:8.16(1)}--\ref{al:8.16(6)} with $\Delta = \Delta_i$, such that no $\tilde T$,
    $\tilde A$ satisfy~\ref{al:8.16(7)}--\ref{al:8.16(9)} with $M = M_i$. Fix $T \in
    \pgrass{\adim}{\vdim}$. Since $\|T_i - U_i\| < 1$, the map $U_i | \im T_i$ is univalent, hence
    maps $\im T_i$ onto $\im U_i$, and there is $v_i \in A_i$ with $U_i(v_i) = U_i(a_i)$. Choose
    a~linear isometry $\iota_i$ of~$\R^{\adim}$ with $\iota_i[\im T_i] = \im T$, and replace $V_i$,
    $a_i$, $r_i$, $T_i$, $U_i$, $A_i$, $\alpha_i$ by $(\iota_i \circ \scale{7/r_i} \circ
    \trans{-v_i})_\# V_i$, $7 r_i^{-1} \iota_i(a_i - v_i)$, $7$, $T$, $\iota_i \circ U_i \circ
    \iota_i^{-1}$, $\im T$ and $(r_i/7)^{1 - \vdim/p} \alpha_i$. By~\ref{al:3.2(1)}, \ref{al:3.2(2)}
    and~\ref{al:8.2}, the hypotheses remain valid with the same $\mu_i$, and the left side
    of~\ref{al:8.16(9)}, the number $r^{1-\vdim/p} \alpha$ and the statement~\ref{al:8.16(8)} are
    transformed consistently. Thus we may assume $r_i = 7$, $T_i = T$, $A_i = \im T$ and $U_i(a_i) =
    0$, so that $\cylind{U_i}{a_i}{s} = \cylind{U_i}{0}{s}$; and for every $\tilde T$, $\tilde A$ at
    least one of the following three statements is false:
    \begin{gather}
        \label{eq:8.16:19}
        \tilde T \in \pgrass{\adim}{\vdim} \,, \ \tilde A \in \affgrass{\adim}{\vdim} \,, \text{ and
        $\tilde A$ is parallel to $\tilde T$} \,; \\
        \label{eq:8.16:20}
        \|\tilde T - T\| \le M_i \mu_i \,; \\
        \label{eq:8.16:21}
        \Big( (7 \theta)^{-\vdim-2} \textint{\cylind{U_i}{0}{7 \theta}}{} \dist(x, \tilde A)^2 \ud
        \|V_i\| x \Big)^{1/2} \le \theta^{1 - \vdim/p} \sup \{ \mu_i, M_i 7^{1 - \vdim/p} \alpha_i \}
        \,.
    \end{gather}
    Here $\mu_i = (7^{-\vdim-2} \textint{\cylind{U_i}{0}{7}}{} \dist(x, \im T)^2 \ud \|V_i\|
    x)^{1/2}$, and $\mu_i > 0$, since otherwise $\tilde T = T$ and $\tilde A = \im T$ would satisfy
    \eqref{eq:8.16:19}--\eqref{eq:8.16:21}.

    As in the proof of~\ref{al:8.12}, \ref{al:8.16(5)} extends by closure to all points of the
    support, and \ref{al:8.16(4)} with H\"older's inequality~\cite[2.4.14]{Federer1969} and outer
    regularity gives
    \begin{gather}
        \label{eq:8.16:a}
        \epsilon s^{\vdim} \le \|V_i\| \cball{x}{s} \,, \quad \|V_i\| \oball{x}{s} \le (2 -
        \epsilon) \unitmeasure{\vdim} s^{\vdim} \quad \text{for $x \in \spt \|V_i\| \cap
        \cylind{U_i}{0}{7}$ and $0 < s < 7$} \,, \\
        \label{eq:8.16:b}
        \|\delta V_i\|(E) \le \alpha_i \|V_i\|(E)^{1/q} \quad \text{for $E \subseteq
        \cylind{U_i}{0}{7}$} \,.
    \end{gather}
    Let
    \begin{displaymath}
        \nu_i = \sup \big\{ (7 - |U_i(x)|)^{\vdim} \dist(x, \im T) : x \in \spt \|V_i\| \cap
        \cylind{U_i}{0}{7} \big\} \,.
    \end{displaymath}
    Suppose $x \in \spt \|V_i\| \cap \cylind{U_i}{0}{7}$, $s = 7 - |U_i(x)|$, $d = \dist(x, \im T) >
    0$ and $t = \inf \{ s, d \}/2$. Then $\cball{x}{t} \subseteq \cylind{U_i}{0}{7}$ and
    $\dist(z, \im T) \ge d/2$ for $z \in \cball{x}{t}$, so~\eqref{eq:8.16:a} gives $7^{\vdim+2}
    \mu_i^2 \ge (d/2)^2 \epsilon t^{\vdim}$. If $d \le s$, then $d^{\vdim+2} \le 2^{\vdim+2}
    \epsilon^{-1} 7^{\vdim+2} \mu_i^2$ and $s^{\vdim} d \le 7^{\vdim} d$; if $s < d$, then
    $s^{\vdim} d^2 \le 2^{\vdim+2} \epsilon^{-1} 7^{\vdim+2} \mu_i^2$ and $s^{\vdim} d \le
    (s^{\vdim} d^2)^{1/2} 7^{\vdim/2}$. Since $\mu_i \le \Delta_i \to 0$, we conclude
    \begin{equation}
        \label{eq:8.16:23}
        \lim_{i \to \infty} \nu_i = 0 \,.
    \end{equation}
    Moreover $|T^\perp(a_i)| \le 7^{-\vdim} \nu_i$ and $|T(a_i)| = |(T - U_i)(a_i)| \le \Delta_i
    |a_i|$, hence
    \begin{equation}
        \label{eq:8.16:22}
        |a_i| \le (1 - \Delta_i)^{-1} 7^{-\vdim} \nu_i \to 0 \quad \text{as $i \to \infty$} \,.
    \end{equation}
    Let $\nu_i' = 2^{\vdim} \nu_i$. If $x \in \spt \|V_i\| \cap \oball{0}{13/2}$, then $|U_i(x)| <
    13/2$, hence $\dist(x, \im T) \le \nu_i'$. Henceforth we consider only~$i$ so large that $\nu_i'
    < 1/2$ and $|a_i| < 1/2$. Let
    \begin{displaymath}
        W_i = V_i \restrict (\oball{0}{13/2} \times \pgrass{\adim}{\vdim}) \,.
    \end{displaymath}
    Then $\spt \|W_i\| \subseteq \cball{0}{13/2} \cap \{ x : \dist(x, \im T) \le \nu_i' \}$, and,
    since $|x| \le |T(x)| + |T^\perp(x)|$,
    \begin{equation}
        \label{eq:8.16:24}
        \spt \|W_i\| \cap \cylind{T}{0}{5} \subseteq \oball{0}{11/2} \,.
    \end{equation}

    We show that
    \begin{equation}
        \label{eq:8.16:25}
        \delta V_i(g) \le \Delta_i \mu_i \Big( \textint{}{} |g|^2 \ud \|V_i\| \Big)^{1/2} \quad
        \text{for $g \in \VF(\R^{\adim})$ with $\spt g \subseteq \oball{a_i}{7}$} \,.
    \end{equation}
    Since $U_i(a_i) = 0$, $\oball{a_i}{7} \subseteq \cylind{U_i}{0}{7}$. Since $\vdim < p$ and $p
    \ge 2$, we have $q \le 2$; by H\"older's inequality and~\eqref{eq:8.16:a} at~$a_i$,
    \begin{displaymath}
        \delta V_i(g) \le \alpha_i \|V_i\| \oball{a_i}{7}^{1/q - 1/2} \Big( \textint{}{} |g|^2 \ud
        \|V_i\| \Big)^{1/2} \le \alpha_i \kappa \Big( \textint{}{} |g|^2 \ud \|V_i\| \Big)^{1/2}
    \end{displaymath}
    for such~$g$. If~\eqref{eq:8.16:25} failed for some such~$g$, then $\textint{}{} |g|^2 \ud
    \|V_i\| > 0$ by~\ref{al:8.16(4)}, hence $\Delta_i \mu_i < \alpha_i \kappa$, and, with $\tilde T
    = T$ and $\tilde A = \im T$, the inclusion $\cylind{U_i}{0}{7\theta} \subseteq
    \cylind{U_i}{0}{7}$ and~\eqref{eq:8.16:12} would give
    \begin{displaymath}
        \Big( (7 \theta)^{-\vdim-2} \textint{\cylind{U_i}{0}{7 \theta}}{} \dist(x, \tilde A)^2 \ud
        \|V_i\| x \Big)^{1/2} \le \theta^{-(\vdim+2)/2} \mu_i < \theta^{-(\vdim+2)/2}
        \Delta_i^{-1} \kappa \alpha_i \le \theta^{1 - \vdim/p} M_i 7^{1 - \vdim/p} \alpha_i \,,
    \end{displaymath}
    so that~\eqref{eq:8.16:19}--\eqref{eq:8.16:21} would hold. This proves~\eqref{eq:8.16:25}.

    By~\ref{new:cutoff} choose $\chi \in \dspace{\R^{\adim}}{\R}$ with $0 \le \chi \le 1$, $\chi =
    1$ on a~neighbourhood of $\cball{0}{6}$ and $\spt \chi \subseteq \oball{0}{13/2}$. If $g \in
    \VF(\R^{\adim})$ and $\spt g \subseteq \cylind{T}{0}{5}$, then $\uD g(x) = \uD(\chi g)(x)$ for
    $x \in \spt \|W_i\|$, by~\eqref{eq:8.16:24}, and $\spt (\chi g) \subseteq \oball{0}{13/2}
    \subseteq \oball{a_i}{7}$; hence~\ref{al:4.2:lem} and~\eqref{eq:8.16:25} give $\delta W_i(g) =
    \delta V_i(\chi g)$ and
    \begin{equation}
        \label{eq:8.16:25a}
        \delta W_i(g) \le \Delta_i \mu_i \Big( \textint{}{} |g|^2 \ud \|W_i\| \Big)^{1/2} \quad
        \text{for $g \in \VF(\R^{\adim})$ with $\spt g \subseteq \cylind{T}{0}{5}$} \,.
    \end{equation}
    Let
    \begin{displaymath}
        \beta_i = \Big( \textint{\cylind{T}{0}{5} \times \pgrass{\adim}{\vdim}}{} \|S - T\|^2 \ud
        W_i(x,S) \Big)^{1/2} \,.
    \end{displaymath}
    Applying~\ref{al:8.13} with $G = \oball{a_i}{7}$, $U = \oball{0}{6}$, $\delta = 1/2$, the
    restriction of~$V_i$ to~$G$ and $\alpha = \Delta_i \mu_i$, which is permitted
    by~\eqref{eq:8.16:25}, and using~\eqref{eq:8.16:24} and $\oball{a_i}{7} \subseteq
    \cylind{U_i}{0}{7}$, we obtain
    \begin{displaymath}
        \beta_i^2 \le \textint{\oball{0}{6} \times \pgrass{\adim}{\vdim}}{} \|S - T\|^2 \ud V_i(x,S)
        \le \Delta_i^2 \mu_i^2 + 36 \cdot 7^{\vdim+2} \mu_i^2 \,.
    \end{displaymath}

    Let $\xi$ and $P$ correspond to $\epsilon$ and $\mu = 2^{1/2}$ as in~\ref{al:8.12}, and
    consider $i$ so large that $\nu_i' \le \xi$. The hypotheses of~\ref{al:8.12} hold with $V$
    and~$\alpha$ replaced by $W_i$ and $\Delta_i \mu_i$: \ref{al:8.12(2)} holds since $\spt \|W_i\|
    \cap \cylind{T}{0}{5} \subseteq \oball{0}{11/2} \subseteq \cylind{U_i}{0}{7}$ and $\|W_i\|$
    agrees with $\|V_i\|$ on $\oball{0}{13/2}$; \ref{al:8.12(3)} is~\eqref{eq:8.16:25a};
    \ref{al:8.12(4)} holds because $a_i \in \spt \|W_i\| \cap \cylind{T}{0}{1}$. We obtain $f_i :
    \im T \to \im T^\perp$ and $F_i : \im T \to \R^{\adim}$ with
    \begin{gather}
        \label{eq:8.16:27}
        T(F_i(y)) = y \,, \quad T^\perp(F_i(y)) = f_i(y) \quad \text{for $y \in \im T$} \,, \\
        \label{eq:8.16:28}
        |f_i(y)| \le \nu_i' \,, \quad |f_i(y) - f_i(z)| \le |y - z| \quad \text{for $y, z \in \im
        T$} \,,
    \end{gather}
    and, with $C = \cylind{T}{0}{1}$, $D = \im T \cap \oball{0}{1}$, $X_i = C \cap \im F_i =
    F_i[D]$ and $Y_i = D \cap \{ y : \density^{\vdim}(\|W_i\|,F_i(y)) \text{ exists and is at least
    } 1 \}$,
    \begin{equation}
        \label{eq:8.16:30}
        \|W_i\|(C \without X_i) + \HM^{\vdim}(D \without Y_i) \le P \big( (\Delta_i \mu_i)^2 +
        \beta_i^2 \big) \,.
    \end{equation}
    With $N = (P + 1)(2 + 36 \cdot 7^{\vdim+2})$ we therefore have
    \begin{equation}
        \label{eq:8.16:33}
        \sup \{ \beta_i^2, \|W_i\|(C \without X_i) + \HM^{\vdim}(D \without Y_i) \} \le N \mu_i^2
        \,.
    \end{equation}
    The set $Y_i$ is a~Borel set, since upper and lower densities are Borel functions, see the
    proof of~\ref{al:2.8(5)}, and $F_i$ is continuous.

    By Rademacher's theorem~\cite[3.1.6]{Federer1969}, $f_i$ and~$F_i$ are differentiable at
    $\HM^{\vdim}$ almost all points of $\im T$; at such~$y$, $\uD F_i(y) = j + \uD f_i(y)$, where
    $j : \im T \to \R^{\adim}$ is the inclusion map, and by~\eqref{eq:8.16:28}
    \begin{equation}
        \label{eq:8.16:31}
        \|\uD f_i(y)\| \le 1 \,, \qquad \|\uD F_i(y)\| \le 2 \,, \qquad 1 \le |\extpower{\vdim}
        \uD F_i(y)| \le 2^{\vdim} \,,
    \end{equation}
    the last by~\ref{al:8.14}\eqref{eq:8.14:1} and~\cite[1.7.6]{Federer1969}. Let $S_i(y) \in
    \pgrass{\adim}{\vdim}$ be the orthogonal projection onto $\im \uD F_i(y)$. We show
    \begin{equation}
        \label{eq:8.16:32}
        \textint{X_i \times \pgrass{\adim}{\vdim}}{} \gamma \ud W_i = \textint{Y_i}{} \gamma(F_i(y),
        S_i(y)) \, \density^{\vdim}(\|W_i\|,F_i(y)) \, |\extpower{\vdim} \uD F_i(y)| \ud
        \HM^{\vdim} y
    \end{equation}
    for every bounded Borel function $\gamma$ on $\R^{\adim} \times \pgrass{\adim}{\vdim}$. Let $O =
    \oball{0}{13/2}$. On~$O$, $W_i$ agrees with~$V_i$, $\|\delta W_i\|$ is a~Radon measure
    by~\eqref{eq:8.16:b}, and $\density^{\vdim}(\|W_i\|,x) \ge 1$ for $\|W_i\|$ almost all $x \in
    O$; hence the restriction of~$W_i$ to $O \times \pgrass{\adim}{\vdim}$ is rectifiable
    by~\ref{al:5.5(1)}. By~\ref{al:3.5(1)}(b) and~\ref{rem:2.6(2)}, $\|W_i\| \restrict O = \HM^{\vdim}
    \restrict \density^{\vdim}(\|W_i\|,\cdot)$ on~$O$ and $\textint{}{} \gamma \ud W_i =
    \textint{}{} \gamma(x, \Tan^{\vdim}(\|W_i\|,x)) \density^{\vdim}(\|W_i\|,x) \ud \HM^{\vdim} x$
    whenever $\gamma$ is a~bounded Borel function vanishing outside $K \times \pgrass{\adim}{\vdim}$
    for some compact $K \subseteq O$. The set $X_i \subseteq \cball{0}{2}$ is such a~$K$ up to its
    closure, and $\|W_i\|(X_i \without F_i[Y_i]) = 0$ by~\ref{al:8.16(6)}. By~\cite[3.2.17]{Federer1969},
    applied through a~linear isometry of $\R^{\vdim}$ onto $\im T$ with $K = D$, $\psi = F_i$ and
    $\lambda = 2$ (note $T \circ F_i = \id{\im T}$), we have $\Tan^{\vdim}(\HM^{\vdim} \restrict
    F_i[D], F_i(y)) = \im \uD F_i(y)$ whenever $y \in D$ and $F_i$ is differentiable at~$y$.
    By~\ref{new:density_tangent}(c) with $E = F_i[D]$, and since $F_i$ and $T | X_i$ map
    $\HM^{\vdim}$ null sets onto $\HM^{\vdim}$ null sets~\cite[2.10.11]{Federer1969}, we infer
    $\Tan^{\vdim}(\|W_i\|, F_i(y)) = \im S_i(y)$ for $\HM^{\vdim}$ almost all $y \in Y_i$. Now the
    area formula~\cite[3.2.3, 3.2.5]{Federer1969} for the univalent Lipschitzian map $F_i | Y_i$,
    whose Jacobian is $|\extpower{\vdim} \uD F_i|$ by~\ref{al:2.4}, yields~\eqref{eq:8.16:32}.
    Moreover, the set of $x \in F_i[Y_i]$ with $\density^{\vdim}(\|W_i\|,x) > 1 + \Delta_i$ has
    $\|W_i\|$ measure zero by~\ref{al:8.16(6)}, hence $\HM^{\vdim}$ measure zero, so
    \begin{equation}
        \label{eq:8.16:32a}
        1 \le \density^{\vdim}(\|W_i\|, F_i(y)) \le 1 + \Delta_i \quad \text{for $\HM^{\vdim}$ almost
        all $y \in Y_i$} \,.
    \end{equation}

    Since $\uD f_i(y) = T^\perp \circ S_i(y) \circ \uD F_i(y)$, \ref{al:8.9(3)} gives $\|\uD
    f_i(y)\| \le 2 \|S_i(y) - T\|$, and \ref{al:8.9(5)} with $\eta_1 = \uD f_i(y)$, $\eta_2 = 0$
    gives $\|S_i(y) - T\| \le \|\uD f_i(y)\|$. Using \eqref{eq:8.16:31}, \eqref{eq:8.16:32}
    with $\gamma(x,S) = \inf \{ 1, \dist(x, \im T) \}^2$ and with $\gamma(x,S) = \|S - T\|^2$,
    \eqref{eq:8.16:32a}, \eqref{eq:8.16:28}, \eqref{eq:8.16:33}, and $X_i \subseteq O \subseteq
    \cylind{U_i}{0}{7}$, we estimate
    \begin{gather}
        \label{eq:8.16:34}
        \begin{aligned}
            \textint{Y_i}{} |f_i|^2 \ud \HM^{\vdim} & \le \textint{X_i}{} \dist(x, \im T)^2 \ud
            \|W_i\| x \le 7^{\vdim+2} \mu_i^2 \,, & \textint{D \without Y_i}{} |f_i|^2 \ud
            \HM^{\vdim} & \le \nu_i'^2 N \mu_i^2 \,, \\
            \textint{Y_i}{} \|\uD f_i\|^2 \ud \HM^{\vdim} & \le 4 \textint{X_i \times
            \pgrass{\adim}{\vdim}}{} \|S - T\|^2 \ud W_i(x,S) \le 4 N \mu_i^2 \,, &
            \textint{D \without Y_i}{} \|\uD f_i\|^2 \ud \HM^{\vdim} & \le N \mu_i^2 \,.
        \end{aligned}
    \end{gather}

    Let $\iota : \R^{\vdim} \to \im T$ be a~linear isometry. By~\eqref{eq:8.16:34}, the maps
    $\mu_i^{-1} f_i \circ \iota | \oball{0}{1}$ satisfy the hypotheses of~\ref{new:rellich} with $R
    = 1$ and $Z = \im T^\perp$. Applying~\ref{new:rellich} with $\rho = 1 - 1/m$ for $m = 2, 3,
    \ldots$ and a~diagonal argument, and passing to a~subsequence, which does not affect the
    preceding statements, we obtain an $\HM^{\vdim} \restrict D$ measurable map $h : D \to \im
    T^\perp$ with
    \begin{equation}
        \label{eq:8.16:35}
        \lim_{i \to \infty} \textint{\im T \cap \oball{0}{\rho}}{} |\mu_i^{-1} f_i - h|^2 \ud
        \HM^{\vdim} = 0 \quad \text{for $0 < \rho < 1$} \,.
    \end{equation}
    By~\eqref{eq:8.16:34} and~\eqref{eq:8.16:23}, $\textint{\im T \cap \oball{0}{\rho}}{} |h|^2
    \ud \HM^{\vdim} \le \limsup_{i \to \infty} \mu_i^{-2} \textint{D}{} |f_i|^2 \ud \HM^{\vdim} \le
    7^{\vdim+2}$ for $0 < \rho < 1$, hence
    \begin{equation}
        \label{eq:8.16:36}
        \textint{D}{} |h|^2 \ud \HM^{\vdim} \le 7^{\vdim+2} \,.
    \end{equation}

    Suppose $\phi : \im T \to \im T^\perp$ is of class~$\infty$ and $\spt \phi$ is a~compact subset
    of~$D$. We show
    \begin{equation}
        \label{eq:8.16:37}
        \lim_{i \to \infty} \mu_i^{-1} \textint{D}{} \uD f_i(y) \bullet \uD \phi(y) \ud \HM^{\vdim}
        y = 0 \,.
    \end{equation}
    Let $b_\phi = \sup \{ |\phi(y)| + |\uD \phi(y)| : y \in \im T \}$ and $g_\phi(x) = \chi(x)
    \phi(T(x))$ for $x \in \R^{\adim}$. Then $g_\phi \in \VF(\R^{\adim})$, $\spt g_\phi \subseteq C$,
    and, since $\spt \|W_i\| \cap C \subseteq \oball{0}{2}$ and $\chi = 1$ near $\cball{0}{6}$,
    \ref{al:4.2:lem} gives $\delta W_i(g_\phi) = \textint{}{} (\uD \phi(T(x)) \circ T) \bullet S \ud
    W_i(x,S)$, where the integrand vanishes unless $x \in C$. Let
    \begin{align*}
        a_{1,i} & = \textint{D \without Y_i}{} \uD f_i(y) \bullet \uD \phi(y) \ud \HM^{\vdim} y \,,
        \\
        a_{2,i} & = \textint{Y_i}{} \big( \uD f_i(y) \bullet \uD \phi(y) - [(\uD \phi(y) \circ T)
        \bullet S_i(y)] \, |\extpower{\vdim} \uD F_i(y)| \big) \ud \HM^{\vdim} y \,, \\
        a_{3,i} & = \textint{Y_i}{} [(\uD \phi(y) \circ T) \bullet S_i(y)] \, \big(1 -
        \density^{\vdim}(\|W_i\|,F_i(y)) \big) \, |\extpower{\vdim} \uD F_i(y)| \ud \HM^{\vdim} y
        \,, \\
        a_{4,i} & = \textint{Y_i}{} [(\uD \phi(y) \circ T) \bullet S_i(y)] \,
        \density^{\vdim}(\|W_i\|,F_i(y)) \, |\extpower{\vdim} \uD F_i(y)| \ud \HM^{\vdim} y -
        \delta W_i(g_\phi) \,, \qquad a_{5,i} = \delta W_i(g_\phi) \,,
    \end{align*}
    so that $\textint{D}{} \uD f_i \bullet \uD \phi \ud \HM^{\vdim} = \textsum{j=1}{5} a_{j,i}$. By
    \eqref{eq:8.16:31} and~\eqref{eq:8.16:33}, $|a_{1,i}| \le \vdim^{1/2} b_\phi N \mu_i^2$, since
    $|\uD f_i(y)| \le \vdim^{1/2} \|\uD f_i(y)\|$. By~\ref{al:8.14} with $l = \uD f_i(y)$, $L = \uD
    F_i(y)$ and $\theta$ replaced by $\uD \phi(y)$, and since $\|\uD f_i(y)\| \le 1$,
    by~\eqref{eq:8.16:34}, $|a_{2,i}| \le \mathbf{c}_2 b_\phi \textint{Y_i}{} \|\uD f_i\|^2 \ud
    \HM^{\vdim} \le 4 \mathbf{c}_2 b_\phi N \mu_i^2$. The map $\tau = \uD \phi(y) \circ T$ satisfies
    $\tau = T^\perp \circ \tau \circ T$; \ref{al:8.9(4)} and $|S_i(y) - T|^2 = 2 |T^\perp \circ
    S_i(y)|^2 \le 2 \vdim \|S_i(y) - T\|^2$ give $|\tau \bullet S_i(y)| \le \vdim^{1/2} \|S_i(y) -
    T\| \, |\tau| \le \vdim^{1/2} b_\phi \|\uD f_i(y)\|$. Hence, by~\eqref{eq:8.16:32a},
    \eqref{eq:8.16:31}, Schwarz's inequality and~\eqref{eq:8.16:34},
    \begin{displaymath}
        |a_{3,i}| \le \Delta_i 2^{\vdim} \vdim^{1/2} b_\phi \textint{Y_i}{} \|\uD f_i\| \ud
        \HM^{\vdim} \le \Delta_i 2^{\vdim} \vdim^{1/2} b_\phi (4 N \mu_i^2)^{1/2}
        \unitmeasure{\vdim}^{1/2} \,.
    \end{displaymath}
    By~\eqref{eq:8.16:32} with $\gamma(x,S) = (\uD \phi(T(x)) \circ T) \bullet S$, and $T(F_i(y)) =
    y$, we have $a_{4,i} = - \textint{(C \without X_i) \times \pgrass{\adim}{\vdim}}{} (\uD
    \phi(T(x)) \circ T) \bullet S \ud W_i(x,S)$, hence $|a_{4,i}| \le \vdim^{1/2} b_\phi N \mu_i^2$.
    By~\eqref{eq:8.16:25a} applied to $\pm g_\phi$ and~\eqref{eq:8.16:a}, $|a_{5,i}| \le \Delta_i
    \mu_i b_\phi \|W_i\|(C)^{1/2} \le \Delta_i \mu_i b_\phi ((2 - \epsilon) \unitmeasure{\vdim}
    7^{\vdim})^{1/2}$, because $\spt \|W_i\| \cap C \subseteq \oball{a_i}{7}$. Since $\mu_i \to 0$
    and $\Delta_i \to 0$, \eqref{eq:8.16:37} follows.

    By~\ref{new:ibp}, applied through~$\iota$ to the components of $f_i$ with respect to an
    orthonormal basis of $\im T^\perp$, $\textint{D}{} \uD f_i \bullet \uD \phi \ud \HM^{\vdim} = -
    \textint{D}{} f_i \bullet \textsum{j=1}{\vdim} \uD^2 \phi(\cdot)(\iota(e_j),\iota(e_j)) \ud
    \HM^{\vdim}$. Hence~\eqref{eq:8.16:35} and~\eqref{eq:8.16:37} give $\textint{D}{} h \bullet
    \textsum{j=1}{\vdim} \uD^2 \phi(\cdot)(\iota(e_j),\iota(e_j)) \ud \HM^{\vdim} = 0$
    for all such~$\phi$. By~\ref{new:weyl}, applied through~$\iota$, there is a~harmonic map $H : D
    \to \im T^\perp$ with $h = H$ almost everywhere on~$D$; we replace $h$ by~$H$. By~\ref{al:8.15}
    and~\eqref{eq:8.16:36},
    \begin{equation}
        \label{eq:8.16:H}
        \begin{gathered}
        \sup \{ |h(0)|, \|\uD h(0)\| \} \le \mathbf{c}_3 7^{(\vdim+2)/2} \,, \\ |h(y) - h(0) - \uD
        h(0)(y)| \le \mathbf{c}_3 7^{(\vdim+2)/2} |y|^2 \quad \text{for $y \in D$ with $2 |y| < 1$}
        \,.
        \end{gathered}
    \end{equation}
    Let
    \begin{displaymath}
        L_i(y) = y + \mu_i \uD h(0)(y) \quad \text{for $y \in \im T$} \,, \qquad K_i(x) = L_i(T(x)) +
        \mu_i h(0) \quad \text{for $x \in \R^{\adim}$} \,,
    \end{displaymath}
    let $\tilde T_i \in \pgrass{\adim}{\vdim}$ be the orthogonal projection onto $\im L_i$, and let
    $\tilde A_i = \im K_i = \trans{\mu_i h(0)}[\im L_i]$. Then $\tilde T_i$ and $\tilde A_i$
    satisfy~\eqref{eq:8.16:19}, and by~\ref{al:8.9(5)}, \eqref{eq:8.16:H} and~\eqref{eq:8.16:12},
    \begin{equation}
        \label{eq:8.16:39}
        \|\tilde T_i - T\| \le \mu_i \|\uD h(0)\| \le \mathbf{c}_3 7^{(\vdim+2)/2} \mu_i \le M_i
        \mu_i \,,
    \end{equation}
    which is~\eqref{eq:8.16:20}. If $x \in \R^{\adim}$ and $|T(x)| < 1$, then $x - K_i(x) =
    T^\perp(x) - \mu_i h(0) - \mu_i \uD h(0)(T(x))$, hence
    \begin{equation}
        \label{eq:8.16:40}
        \dist(x, \tilde A_i) \le \dist(x, \im T) + 2 \mathbf{c}_3 7^{(\vdim+2)/2} \mu_i \,;
    \end{equation}
    if $y \in D$ and $2 |y| < 1$, then $F_i(y) - K_i(F_i(y)) = f_i(y) - \mu_i h(y) + \mu_i [h(y) -
    h(0) - \uD h(0)(y)]$, hence
    \begin{equation}
        \label{eq:8.16:41}
        \dist(F_i(y), \tilde A_i) \le |f_i(y) - \mu_i h(y)| + \mathbf{c}_3 7^{(\vdim+2)/2} \mu_i
        |y|^2 \,.
    \end{equation}

    By~\eqref{eq:8.16:10} and $14 \theta < 1$ we may choose $\sigma$ with $7 < \sigma$, $\sigma
    \theta < 1/2$ and $2^{\vdim+1} \mathbf{c}_3^2 (\vdim \unitmeasure{\vdim} / (\vdim + 4))
    \sigma^{\vdim+4} \theta^2 < \theta^{2(1 - \vdim/p)}$. If $x \in \spt \|V_i\| \cap
    \cylind{U_i}{0}{7 \theta}$, then $7 - |U_i(x)| > 6$, so $\dist(x, \im T) \le \nu_i$, and $|T(x)|
    \le |U_i(x)| + \Delta_i |x|$, $|x| \le |T(x)| + \nu_i$, whence $|x| \le (1 - \Delta_i)^{-1} (7
    \theta + \nu_i)$. Therefore, for large~$i$,
    \begin{displaymath}
        \spt \|V_i\| \cap \cylind{U_i}{0}{7 \theta} \subseteq \oball{0}{13/2} \cap
        \cylind{T}{0}{\sigma \theta} \,.
    \end{displaymath}
    Using \eqref{eq:8.16:40}, \eqref{eq:8.16:33} on $C \without X_i$, and~\eqref{eq:8.16:32} with
    $\gamma(x,S) = \inf \{ 1, \dist(x, \tilde A_i) \}^2$ for $x \in \cylind{T}{0}{\sigma \theta}$
    and $\gamma(x,S) = 0$ otherwise, together with~\eqref{eq:8.16:41}, \eqref{eq:8.16:32a},
    \eqref{eq:8.16:31} and~\cite[3.2.13]{Federer1969}, we obtain for large~$i$
    \begin{multline*}
        \textint{\cylind{U_i}{0}{7 \theta}}{} \dist(x, \tilde A_i)^2 \ud \|V_i\| x \le
        \textint{\cylind{T}{0}{\sigma \theta} \cap X_i}{} \dist(x, \tilde A_i)^2 \ud \|W_i\| x +
        \textint{C \without X_i}{} \dist(x, \tilde A_i)^2 \ud \|W_i\| x \\
        \le 2 (1 + \Delta_i) 2^{\vdim} \Big( \textint{\im T \cap \oball{0}{1/2}}{} |f_i - \mu_i
        h|^2 \ud \HM^{\vdim} + \mathbf{c}_3^2 7^{\vdim+2} \mu_i^2 \frac{\vdim
        \unitmeasure{\vdim}}{\vdim + 4} (\sigma \theta)^{\vdim+4} \Big) + 2 (\nu_i'^2 + 4
        \mathbf{c}_3^2 7^{\vdim+2} \mu_i^2) N \mu_i^2 \,.
    \end{multline*}
    By~\eqref{eq:8.16:35}, \eqref{eq:8.16:23} and $\Delta_i \to 0$,
    \begin{displaymath}
        \limsup_{i \to \infty} \mu_i^{-2} (7 \theta)^{-\vdim-2} \textint{\cylind{U_i}{0}{7
        \theta}}{} \dist(x, \tilde A_i)^2 \ud \|V_i\| x \le 2^{\vdim+1} \mathbf{c}_3^2 \frac{\vdim
        \unitmeasure{\vdim}}{\vdim + 4} \sigma^{\vdim+4} \theta^2 < \theta^{2(1 - \vdim/p)} \,.
    \end{displaymath}
    Hence~\eqref{eq:8.16:21} holds with $\tilde A = \tilde A_i$ for large~$i$. Together
    with~\eqref{eq:8.16:19} and~\eqref{eq:8.16:39}, this contradicts the choice of $V_i$, $a_i$,
    $U_i$, $\mu_i$, $\alpha_i$.
\end{proof}

\begin{remark}
    \label{rem:8.16}
    Allard assumes $\spt \|V_i\| \cap \cylind{T}{0}{5} \subseteq \oball{a_i}{7}$ and states that
    the proof needs only a~``trivial modification'' otherwise; the modification is the passage to
    $W_i$ above. The convergence~\eqref{eq:8.16:35}, which Allard obtains by ``well-known
    reasoning'' as a~$\liminf$ over all of~$D$, is established here on the balls of radius less
    than~$1$, which suffices. Allard's reference~\cite[2.5.6]{Federer1969} for the regularity of
    weakly harmonic functions is a~misprint for~\cite[5.2.6]{Federer1969}, and is replaced
    by~\ref{new:weyl}. Allard's bound $\sup \im |f_i| \le \nu_i$ is replaced by $2^{\vdim} \nu_i$, and
    his constants $9 \cdot 7^{\vdim+2}$ in the estimate of~$\beta_i$ and $n^{1/2}$ in the estimate
    of $a_{3,i}$ by $36 \cdot 7^{\vdim+2}$ and~$\vdim^{1/2}$; these changes do not affect the
    argument.
\end{remark}

\begin{corollary}[{\cite[8.17]{Allard1972}}]
    \label{al:8.17}
    Corresponding to every $\epsilon$ with $0 < \epsilon < 1$ there are $\gamma$ and~$N$ with $0 <
    \gamma < 1$ and $1 < N < \infty$, and with the following property. If $V$, $a$, $r$, $T$, $U$,
    $A$, $\mu$ and~$\alpha$ satisfy \ref{al:8.16(1)}--\ref{al:8.16(6)} of~\ref{al:8.16} with $\Delta$
    replaced by~$\gamma$, and $r^{1 - \vdim/p} \alpha \le \gamma$, then
    \begin{enumerate}[label=(\arabic*)]
        \item \label{al:8.17(1)} $\spt \|V\| \cap \{ x : U(x) = U(a) \} = \{ a \}$;
        \item \label{al:8.17(2)} there is $P \in \pgrass{\adim}{\vdim}$ with $\Tan(\spt \|V\|, a) =
            \im P$ and $\|T - P\| \le N \sup \{ \mu, r^{1 - \vdim/p} \alpha \}$;
        \item \label{al:8.17(3)} whenever $0 < s < r$, there are $T_s \in \pgrass{\adim}{\vdim}$ and
            $A_s \in \affgrass{\adim}{\vdim}$ such that $A_s$ is parallel to $T_s$ and neither of the
            numbers
            \begin{displaymath}
                \|T_s - P\| \,, \qquad \Big( s^{-\vdim-2} \textint{\cylind{U}{a}{s}}{} \dist(x,
                A_s)^2 \ud \|V\| x \Big)^{1/2}
            \end{displaymath}
            exceeds $N \sup \{ \mu, r^{1 - \vdim/p} \alpha \} (s/r)^{1 - \vdim/p}$.
    \end{enumerate}
\end{corollary}

\begin{proof}
    Let $\Delta$, $\theta$ and~$M$ correspond to~$\epsilon$ as in~\ref{al:8.16}, and let $\lambda
    = \theta^{1 - \vdim/p}$; then $0 < \lambda < 1$. Choose $\gamma$ and~$N$ with $0 < \gamma < 1 <
    N < \infty$ and
    \begin{displaymath}
        M \gamma \le \Delta \,, \qquad \gamma + M^2 \gamma (1 - \lambda)^{-1} \le \Delta \,, \qquad
        N \ge M^2 \lambda^{-1} (1 - \lambda)^{-1} \,, \qquad N \ge M \theta^{-(\vdim+2)/2}
        \lambda^{-1} \,.
    \end{displaymath}
    Suppose $V$, $a$, $r$, $T$, $U$, $A$, $\mu$, $\alpha$ are as in the hypotheses, and let $Q =
    \sup \{ \mu, M r^{1 - \vdim/p} \alpha \}$; then $Q \le M \sup \{ \mu, r^{1-\vdim/p} \alpha \}
    \le M \gamma \le \Delta$. We construct inductively $T_j \in \pgrass{\adim}{\vdim}$ and $A_j \in
    \affgrass{\adim}{\vdim}$ parallel to~$T_j$, for $j = 0, 1, 2, \ldots$, with $T_0 = T$, $A_0 =
    A$,
    \begin{equation}
        \label{eq:8.17:a}
        \mu_j = \Big( (r \theta^j)^{-\vdim-2} \textint{\cylind{U}{a}{r \theta^j}}{} \dist(x, A_j)^2
        \ud \|V\| x \Big)^{1/2} \le Q \lambda^j \,, \qquad \|T_j - T_{j-1}\| \le M Q \lambda^{j-1}
        \quad \text{for $j \ge 1$} \,.
    \end{equation}
    The case $j = 0$ holds since $\mu_0 = \mu$. Suppose $T_0, \ldots, T_l$ and $A_0, \ldots, A_l$
    satisfy~\eqref{eq:8.17:a}. Then $\mu_l \le Q \le \Delta$ and
    \begin{displaymath}
        \|T_l - U\| \le \|T - U\| + \textsum{j=1}{l} \|T_j - T_{j-1}\| \le \gamma + M Q (1 -
        \lambda)^{-1} \le \gamma + M^2 \gamma (1 - \lambda)^{-1} \le \Delta \,.
    \end{displaymath}
    Since $\cylind{U}{a}{r \theta^l} \subseteq \cylind{U}{a}{r}$ and $\gamma \le \Delta$, the
    hypotheses of~\ref{al:8.16} hold with $r$, $T$, $A$, $\mu$ replaced by $r \theta^l$, $T_l$,
    $A_l$, $\mu_l$. Hence there are $T_{l+1} \in \pgrass{\adim}{\vdim}$ and $A_{l+1} \in
    \affgrass{\adim}{\vdim}$ parallel to $T_{l+1}$ with $\|T_{l+1} - T_l\| \le M \mu_l \le M Q
    \lambda^l$ and
    \begin{displaymath}
        \mu_{l+1} \le \lambda \sup \{ \mu_l, M (r \theta^l)^{1 - \vdim/p} \alpha \} = \lambda \sup
        \{ \mu_l, \lambda^l M r^{1 - \vdim/p} \alpha \} \le Q \lambda^{l+1} \,.
    \end{displaymath}
    By~\eqref{eq:8.17:a} and~\ref{new:grass_compact}, $T_\infty = \lim_{j \to \infty} T_j$ exists in
    $\pgrass{\adim}{\vdim}$, and
    \begin{equation}
        \label{eq:8.17:b}
        \|T_\infty - T_j\| \le M Q \lambda^j (1 - \lambda)^{-1} \quad \text{for $j = 0, 1, 2,
        \ldots$} \,, \qquad \|T_\infty - U\| \le \Delta < 1 \,.
    \end{equation}

    As in the proof of~\ref{al:8.12}, \ref{al:8.16(5)} gives $\|V\| \cball{x}{s} \ge \epsilon
    s^{\vdim}$ for $x \in \spt \|V\| \cap \cylind{U}{a}{r}$ and $0 < s < r$. Let
    \begin{displaymath}
        \omega_j = \sup \{ (r \theta^j)^{-1} \dist(x, A_j) : x \in \spt \|V\| \cap \cylind{U}{a}{r
        \theta^j / 2} \} \,.
    \end{displaymath}
    If $x \in \spt \|V\| \cap \cylind{U}{a}{r \theta^j/2}$, $d = \dist(x, A_j) > 0$ and $t = \inf \{
    d, r \theta^j \}/2$, then $\cball{x}{t} \subseteq \cylind{U}{a}{r \theta^j}$ and $\dist(z,
    A_j) \ge d/2$ for $z \in \cball{x}{t}$, hence $(r \theta^j)^{\vdim+2} \mu_j^2 \ge (d/2)^2
    \epsilon t^{\vdim}$; considering the cases $t = d/2$ and $t = r \theta^j/2$, we obtain $(r
    \theta^j)^{-1} d \le \sup \{ (2^{\vdim+2} \epsilon^{-1} \mu_j^2)^{1/(\vdim+2)}, (2^{\vdim+2}
    \epsilon^{-1})^{1/2} \mu_j \}$. Since $\mu_j \to 0$, we infer $\lim_{j \to \infty} \omega_j =
    0$.

    Suppose $x \in \spt \|V\|$, $x \ne a$ and $U(x) = U(a)$. Then $x \in \cylind{U}{a}{r
    \theta^j/2}$ for all~$j$, so $|T_j^\perp(x - a)| \le \dist(x, A_j) + \dist(a, A_j) \le 2
    \omega_j r \theta^j$, because $A_j$ is a~translate of $\im T_j$; letting $j \to \infty$ we obtain
    $T_\infty^\perp(x - a) = 0$. But then $|U(x - a)| \ge |T_\infty(x - a)| - \|T_\infty - U\| |x -
    a| = (1 - \|T_\infty - U\|) |x - a| > 0$ by~\eqref{eq:8.17:b}, a~contradiction.
    This proves~\ref{al:8.17(1)}.

    Suppose $x_m \in \spt \|V\| \without \{ a \}$ and $x_m \to a$. Let $j_m$ be the largest
    integer~$j$ with $|x_m - a| < r \theta^j/2$; then $j_m \to \infty$ and $|x_m - a| \ge r
    \theta^{j_m + 1}/2$. As above, $|T_{j_m}^\perp(x_m - a)| \le 2 \omega_{j_m} r \theta^{j_m} \le
    4 \theta^{-1} \omega_{j_m} |x_m - a|$, hence by~\eqref{eq:8.17:b}
    \begin{displaymath}
        |T_\infty^\perp(x_m - a)| / |x_m - a| \le \|T_\infty - T_{j_m}\| + 4 \theta^{-1}
        \omega_{j_m} \to 0 \quad \text{as $m \to \infty$} \,.
    \end{displaymath}
    By the definition of the tangent cone in~\ref{al:2.1}, $\Tan(\spt \|V\|, a) \subseteq \im
    T_\infty$. To prove the opposite inclusion we note that, by~\ref{al:8.6} applied with $G =
    \cylind{U}{a}{r}$, $\density^{\vdim}(\|V\|,a)$ exists and is at least~$1$; and that
    $\density^{\vdim-1}(\|\delta V\|, a) = 0$, because $\|\delta V\| \cball{a}{s} \le \alpha \|V\|
    \cball{a}{s}^{1/q} \le \alpha (2 \unitmeasure{\vdim} s^{\vdim})^{1/q}$ for small~$s$
    by~\ref{al:8.16(4)}, \ref{al:8.16(5)} and outer regularity, and $s^{1 - \vdim} s^{\vdim/q} =
    s^{1 - \vdim/p} \to 0$. By~\ref{al:3.4(1)}(a) there is $C \in \VarTan(V,a)$; by~\ref{al:4.12},
    applied to $V \restrict (\cylind{U}{a}{r} \times \pgrass{\adim}{\vdim})$ as in~\ref{al:6.5},
    $C$ is stationary, and by~\ref{al:3.4(1)}(b), $\|C\| \cball{0}{1} = \density^{\vdim}(\|V\|,a)
    \unitmeasure{\vdim} > 0$. If $y \in \spt \|C\| \without \{0\}$ and $C = \lim_{i \to \infty}
    (\scale{r_i} \circ \trans{-a})_\# V$ with $r_i \to \infty$, then~\ref{al:2.6(2)(a)}(c) yields
    $x_i \in \spt \|V\|$ with $r_i(x_i - a) \to y$, so $x_i \to a$ and $y \in \Tan(\spt \|V\|, a)$.
    Hence $\spt \|C\| \subseteq \im T_\infty$, and~\ref{al:4.6(3)}, applied with $M = \im
    T_\infty$, gives $C = c \var{}(\im T_\infty)$ for some $0 < c < \infty$. Therefore $\im T_\infty
    = \spt \|C\| \subseteq \Tan(\spt \|V\|, a)$, and $P = T_\infty$ satisfies
    \begin{displaymath}
        \|T - P\| \le \textsum{j=1}{\infty} \|T_j - T_{j-1}\| \le M Q (1 - \lambda)^{-1} \le M^2 (1
        - \lambda)^{-1} \sup \{ \mu, r^{1 - \vdim/p} \alpha \} \le N \sup \{ \mu, r^{1 - \vdim/p}
        \alpha \} \,,
    \end{displaymath}
    which is~\ref{al:8.17(2)}.

    Suppose $0 < s < r$ and choose~$j$ with $r \theta^{j+1} \le s < r \theta^j$; then $\lambda^j \le
    \lambda^{-1} (s/r)^{1 - \vdim/p}$. Let $T_s = T_j$ and $A_s = A_j$. By~\eqref{eq:8.17:b},
    $\|T_s - P\| \le M Q \lambda^j (1 - \lambda)^{-1} \le M^2 \lambda^{-1} (1 - \lambda)^{-1} \sup
    \{ \mu, r^{1 - \vdim/p} \alpha \} (s/r)^{1 - \vdim/p}$, and by~\eqref{eq:8.17:a}
    \begin{displaymath}
        \Big( s^{-\vdim-2} \textint{\cylind{U}{a}{s}}{} \dist(x, A_s)^2 \ud \|V\| x \Big)^{1/2} \le
        \theta^{-(\vdim+2)/2} \mu_j \le \theta^{-(\vdim+2)/2} \lambda^{-1} M \sup \{ \mu, r^{1 -
        \vdim/p} \alpha \} (s/r)^{1 - \vdim/p} \,.
    \end{displaymath}
    The choice of~$N$ yields~\ref{al:8.17(3)}.
\end{proof}

\begin{remark}
    \label{rem:8.17}
    In Allard's choice of~$N$ the condition $\theta^{-(\vdim+2)/2 - (1 - \vdim/p)} \le N$ lacks the
    factor~$M$ needed in the last estimate, since $\sup \{ \mu, M r^{1-\vdim/p} \alpha \}$ may exceed
    $\sup \{ \mu, r^{1-\vdim/p} \alpha \}$ by the factor~$M$. In the proof, Allard writes $\|V\|
    \cball{x}{s} \ge \epsilon \unitmeasure{\vdim} s^{\vdim}$ where~\ref{al:8.16(5)} gives $\epsilon
    s^{\vdim}$. He states the inclusion $\Tan(\spt \|V\|, a) \subseteq \im T_\infty$ and the
    equality~\ref{al:8.17(1)} as evident.
\end{remark}

\begin{lemma}[{\cite[8.18]{Allard1972}}]
    \label{al:8.18}
    Suppose $U$ is an open subset of~$\R^{\vdim}$, $B$ is a~closed subset of~$\R^{\vdim}$, and
    $\Tan(B,b) = \R^{\vdim}$ for every $b \in B \cap U$. Then $B \cap U$ is open.
\end{lemma}

\begin{proof}
    Suppose $B \cap U$ is not open, and let $b_0 \in B \cap U$ be a~point which is not interior to
    $B \cap U$. Choose $\rho > 0$ with $\oball{b_0}{\rho} \subseteq U$ and $u \in \oball{b_0}{\rho/2}
    \without B$. Then $\dist(u, B) \le |u - b_0| < \rho/2 < \dist(u, \R^{\vdim} \without U)$. Since
    $B$ is closed, there is $b \in B$ with $|u - b| = \dist(u,B)$, and $b \in U$. Let $v = u - b \ne
    0$. Since $v \in \Tan(B,b)$, there are $b_m \in B \without \{ b \}$ with $b_m \to b$ and
    $|b_m - b|^{-1} (b_m - b) \to |v|^{-1} v$. With $w_m = b_m - b$,
    \begin{displaymath}
        |u - b_m|^2 - |u - b|^2 = |w_m| \big( |w_m| - 2 v \bullet |w_m|^{-1} w_m \big) \,,
    \end{displaymath}
    which is negative for large~$m$, because $v \bullet |w_m|^{-1} w_m \to |v| > 0$ and $|w_m| \to 0$.
    This contradicts $|u - b| = \dist(u, B)$.
\end{proof}

\begin{theorem}[{\cite[8.19]{Allard1972}}]
    \label{al:8.19}
    Suppose $0 < \epsilon < 1$. There are $\eta$ and~$C$ with $0 < \eta < 1$ and $1 < C < \infty$,
    and with the following property. If $V$, $r$, $T$, $\mu$ and~$\alpha$ are such that
    \begin{enumerate}[label=(\arabic*)]
        \item \label{al:8.19(1)} $V \in \Var{\vdim}(\R^{\adim})$, $0 \in \spt \|V\|$, $0 < r <
            \infty$, $T \in \pgrass{\adim}{\vdim}$;
        \item \label{al:8.19(2)} $\mu = (r^{-\vdim-2} \textint{\cylind{T}{0}{2r}}{} \dist(x, \im T)^2
            \ud \|V\| x)^{1/2} \le \eta$;
        \item \label{al:8.19(3)} $0 \le r^{1 - \vdim/p} \alpha \le \eta$, and $\delta V(g) \le \alpha
            (\textint{}{} |g|^q \ud \|V\|)^{1/q}$ for $g \in \VF(\R^{\adim})$ with $\spt g \subseteq
            \cylind{T}{0}{2r}$;
        \item \label{al:8.19(4)} for $\|V\|$ almost all $x \in \cylind{T}{0}{2r}$, $\epsilon \le
            s^{-\vdim} \|V\| \cball{x}{s} \le (2 - \epsilon) \unitmeasure{\vdim}$ for $0 < s < r$;
        \item \label{al:8.19(5)} for $\|V\|$ almost all $x \in \cylind{T}{0}{2r}$, the density
            $\density^{\vdim}(\|V\|,x)$ exists and $1 \le \density^{\vdim}(\|V\|,x) \le 1 + \eta$,
    \end{enumerate}
    then there are $f : \im T \cap \oball{0}{r} \to \im T^\perp$ and $F : \im T \cap \oball{0}{r}
    \to \R^{\adim}$ such that
    \begin{enumerate}[label=(\arabic*),start=6]
        \item \label{al:8.19(6)} $T(F(y)) = y$ and $T^\perp(F(y)) = f(y)$ for $y \in \im T \cap
            \oball{0}{r}$;
        \item \label{al:8.19(7)} $\spt \|V\| \cap \cylind{T}{0}{r} = \im F$;
        \item \label{al:8.19(8)} $f$ is differentiable at every point of $\im T \cap \oball{0}{r}$;
        \item \label{al:8.19(9)} $\|\uD f(y) - \uD f(z)\| \le C \sup \{ \mu, r^{1 - \vdim/p} \alpha
            \} (|y - z|/r)^{1 - \vdim/p}$ for $y, z \in \im T \cap \oball{0}{r}$;
        \item \label{al:8.19(10)} $\|\uD f(y)\| \le C \sup \{ \mu, r^{1 - \vdim/p} \alpha \}$ for $y
            \in \im T \cap \oball{0}{r}$.
    \end{enumerate}
\end{theorem}

\begin{proof}
    Let $\gamma$ and~$N$ correspond to~$\epsilon$ as in~\ref{al:8.17}. Choose $\eta$ with $0 < \eta
    < 1$, $\eta \le \gamma$, $2^{(\vdim+2)/2} N \eta \le \gamma$ and $N \eta \le 1/2$, and let $C = 8
    N + 3 \cdot 2^{(\vdim+2)/2} N^2$. Suppose $V$, $r$, $T$, $\mu$, $\alpha$ satisfy the hypotheses,
    and let $Q = \sup \{ \mu, r^{1 - \vdim/p} \alpha \}$.

    As in the proof of~\ref{al:8.12}, \ref{al:8.19(4)} holds at all points of $\spt \|V\| \cap
    \cylind{T}{0}{2r}$. If $x \in \spt \|V\| \cap \cylind{T}{0}{r}$, $d = \dist(x, \im T) > 0$ and
    $t = \inf \{ d, r \}/2$, then $\cball{x}{t} \subseteq \cylind{T}{0}{2r}$ and, as in the proof
    of~\ref{al:8.17}, $r^{\vdim+2} \mu^2 \ge (d/2)^2 \epsilon t^{\vdim}$; hence $d \le \sigma r$,
    where $\sigma = \sup \{ (2^{\vdim+2} \epsilon^{-1})^{1/(\vdim+2)}, (2^{\vdim+2}
    \epsilon^{-1})^{1/2} \}$.

    Suppose $a \in \spt \|V\| \cap \cylind{T}{0}{r}$. Then $\cylind{T}{a}{r} \subseteq
    \cylind{T}{0}{2r}$, and the hypotheses of~\ref{al:8.17} hold with $U = T$, $A = \im T$,
    $\Delta$ replaced by~$\gamma$, and $\mu$ replaced by $\mu_a = (r^{-\vdim-2}
    \textint{\cylind{T}{a}{r}}{} \dist(x, \im T)^2 \ud \|V\| x)^{1/2} \le \mu$. By~\ref{al:8.17}
    there is $P_a \in \pgrass{\adim}{\vdim}$ with
    \begin{equation}
        \label{eq:8.19:a}
        \spt \|V\| \cap \{ x : T(x) = T(a) \} = \{ a \} \,, \qquad \Tan(\spt \|V\|, a) = \im P_a \,,
        \qquad \|T - P_a\| \le N Q \le \tfrac12 \,.
    \end{equation}
    By~\ref{al:8.9(3)}, $|T(v)| \ge |P_a(v)| - \|P_a - T\| |v| \ge \frac12 |v|$ for $v \in \im P_a$,
    so $T$ maps $\im P_a$ onto $\im T$.

    Let $K = \spt \|V\| \cap \{ x : |T(x)| \le r, \ \dist(x, \im T) \le \sigma r \}$; then $K$ is
    compact, $B = T[K]$ is closed, and $B \cap \oball{0}{r} = T[\spt \|V\| \cap \cylind{T}{0}{r}]$.
    Suppose $a \in \spt \|V\| \cap \cylind{T}{0}{r}$ and $v \in \im P_a \without \{0\}$. There are
    $x_m \in \spt \|V\| \without \{a\}$ with $x_m \to a$ and $|x_m - a|^{-1} (x_m - a) \to |v|^{-1}
    v$; for large~$m$, $x_m \in K$, $T(x_m) \ne T(a)$ by~\eqref{eq:8.19:a}, and $|T(x_m) -
    T(a)|^{-1} (T(x_m) - T(a)) \to |T(v)|^{-1} T(v)$. Hence $T[\im P_a] \subseteq \Tan(B, T(a))$, and
    $\Tan(B, T(a)) = \im T$. Applying~\ref{al:8.18} through a~linear isometry of $\R^{\vdim}$ onto
    $\im T$, we see that $B \cap \oball{0}{r}$ is open in $\im T$; it is also relatively closed in
    the connected set $\im T \cap \oball{0}{r}$, and contains $T(0) = 0$. Therefore $T$ maps $\spt
    \|V\| \cap \cylind{T}{0}{r}$ univalently onto $\im T \cap \oball{0}{r}$. Let $F$ be the inverse
    map and $f = T^\perp \circ F$; then~\ref{al:8.19(6)} and~\ref{al:8.19(7)} hold. For $0 < \rho <
    r$, $T$ maps the compact set $K \cap \{ x : |T(x)| \le \rho \}$ univalently onto $\im T \cap
    \cball{0}{\rho}$, hence $F$ is continuous.

    Suppose $y \in \im T \cap \oball{0}{r}$, $a = F(y)$, and let $L_a = (T | \im P_a)^{-1}$,
    regarded as a~map of $\im T$ into~$\R^{\adim}$, so $\|L_a\| \le 2$. If $z \to y$, $z \ne y$, then
    $w = F(z) - a \to 0$, $w \ne 0$, and every limit point of $|w|^{-1} w$ lies in $\im P_a$, by the
    definition of $\Tan(\spt \|V\|, a)$; thus $|P_a^\perp(w)| / |w| \to 0$. Since $T(w) = z - y$, we
    have $P_a(w) = L_a(z - y - T(P_a^\perp(w)))$, hence $w - L_a(z - y) = P_a^\perp(w) - L_a(T(P_a^\perp
    (w)))$ and $|w - L_a(z - y)| \le 3 |P_a^\perp(w)|$. In particular $|w| \le 2 |z - y| + 3
    |P_a^\perp(w)|$, so $\limsup |w|/|z - y| \le 2$, and $|F(z) - F(y) - L_a(z - y)| / |z - y| \to
    0$. Thus $F$ and~$f$ are differentiable at~$y$, $\uD F(y) = L_a$, $\im \uD F(y) = \im P_a$, and
    $\uD f(y) = T^\perp \circ L_a$; this is~\ref{al:8.19(8)}, and $P_a$ is the orthogonal projection
    onto $\{ v + \uD f(y)(v) : v \in \im T \}$. By~\ref{al:8.9(5)} with $S = T$, $\eta_1 = \uD
    f(y)$, $\eta_2 = 0$, and~\eqref{eq:8.19:a},
    \begin{equation}
        \label{eq:8.19:b}
        \|\uD f(y)\|^2 \le (1 - \|P_a - T\|^2)^{-1} \|P_a - T\|^2 \le \tfrac43 N^2 Q^2 \,, \qquad
        \|\uD f(y)\|^2 \le \tfrac13 \,,
    \end{equation}
    which implies~\ref{al:8.19(10)}.

    Suppose $y, z \in \im T \cap \oball{0}{r}$, $y \ne z$, $a = F(y)$, $b = F(z)$. If $|y - z| \ge
    r/2$, then~\eqref{eq:8.19:b} gives $\|\uD f(y) - \uD f(z)\| \le 4 N Q \le 8 N Q (|y - z|/r)^{1 -
    \vdim/p}$. Suppose $|y - z| < r/2$ and let $s = 2 |y - z|$. By~\ref{al:8.17(3)} at~$a$ there are
    $T_s$ and~$A_s$ with $A_s$ parallel to~$T_s$, $\|T_s - P_a\| \le N Q (s/r)^{1 - \vdim/p}$ and
    \begin{displaymath}
        \Big( s^{-\vdim-2} \textint{\cylind{T}{a}{s}}{} \dist(x, A_s)^2 \ud \|V\| x \Big)^{1/2} \le N
        Q (s/r)^{1 - \vdim/p} \,.
    \end{displaymath}
    Then $\|T_s - T\| \le \|T_s - P_a\| + \|P_a - T\| \le 2 N \eta \le \gamma$, and, since
    $\cylind{T}{b}{|y - z|} \subseteq \cylind{T}{a}{s}$,
    \begin{displaymath}
        \Big( |y - z|^{-\vdim-2} \textint{\cylind{T}{b}{|y - z|}}{} \dist(x, A_s)^2 \ud \|V\| x
        \Big)^{1/2} \le 2^{(\vdim+2)/2} N Q (s/r)^{1 - \vdim/p} \le 2^{(\vdim+2)/2} N \eta \le
        \gamma \,.
    \end{displaymath}
    Moreover $b \in \spt \|V\|$, $\cylind{T}{b}{|y - z|} \subseteq \cylind{T}{0}{2r}$ and $|y -
    z|^{1 - \vdim/p} \alpha \le \eta \le \gamma$. Hence~\ref{al:8.17} applies with $a$, $r$, $T$,
    $U$, $A$ replaced by $b$, $|y - z|$, $T_s$, $T$, $A_s$, and~\ref{al:8.17(2)} yields
    \begin{displaymath}
        \|T_s - P_b\| \le N \sup \{ 2^{(\vdim+2)/2} N Q (s/r)^{1 - \vdim/p}, |y - z|^{1 - \vdim/p}
        \alpha \} \le 2^{(\vdim+2)/2} N^2 Q (s/r)^{1 - \vdim/p} \,,
    \end{displaymath}
    because $|y - z|^{1-\vdim/p} \alpha = r^{1 - \vdim/p} \alpha (|y-z|/r)^{1 - \vdim/p}$. Thus
    $\|P_a - P_b\| \le (N + 2^{(\vdim+2)/2} N^2) Q (s/r)^{1 - \vdim/p}$. By~\ref{al:8.9(5)} with $S
    = T$, $\eta_1 = \uD f(y)$, $\eta_2 = \uD f(z)$, and by~\eqref{eq:8.19:a}, \eqref{eq:8.19:b},
    \begin{displaymath}
        \|\uD f(y) - \uD f(z)\|^2 \le (1 - \|P_a - T\|^2)^{-1} (1 + \|\uD f(z)\|^2) \|P_a - P_b\|^2
        \le \tfrac{16}{9} \|P_a - P_b\|^2 \,,
    \end{displaymath}
    hence $\|\uD f(y) - \uD f(z)\| \le \frac43 \cdot 2^{1 - \vdim/p} (N + 2^{(\vdim+2)/2} N^2) Q
    (|y - z|/r)^{1 - \vdim/p} \le C Q (|y - z|/r)^{1 - \vdim/p}$.
\end{proof}

\begin{remark}
    \label{rem:8.19}
    Allard's hypothesis~\ref{al:8.19(3)} is printed as $0 \le r^{p - \vdim} \alpha \le \eta$; the
    proof uses $r^{1 - \vdim/p} \alpha \le \eta$, which is the scale invariant quantity
    of~\ref{al:8.2}. His proof begins ``Let $\gamma$ and~$N$ be as in~8.18'', which should read
    \ref{al:8.17}. The conclusion~\ref{al:8.19(10)}, which is needed in the proof of~\ref{al:8.RT},
    is contained in Allard's proof. Allard does not verify that the tangent cone of the support at
    $a$ projects onto a~neighbourhood, nor the differentiability of~$F$; both are supplied above.
\end{remark}

\begin{miniremark}
    \label{mr:8.RT}
    The regularity theorem is stated by Allard at the beginning of his section~8 with the remark
    that it follows from 8.8 and~8.19. After normalisation, \ref{al:8.8} yields a~plane $\im T$
    close to the support in a~large ball and the bounds on densities needed by~\ref{al:8.19}
    at every point of the support and every small scale. Then~\ref{al:8.19}, applied at each point
    of the support with the fixed plane $\im T$, represents the support near that point as the graph
    of a~continuously differentiable map over $\im T$; a~connectedness argument assembles these
    graphs into one graph over a~disc, and a~cut-off function extends it to $\im T$.
\end{miniremark}

\begin{theorem}[{\cite[section~8, Regularity Theorem]{Allard1972}}]
    \label{al:8.RT}
    Corresponding to every $\epsilon$ with $0 < \epsilon < 1$ there is $\eta > 0$ with the
    following property. Suppose $0 < R < \infty$, $0 < d < \infty$, $V \in \Var{\vdim}(\R^{\adim})$, $a
    \in \spt \|V\|$,
    \begin{gather*}
        \udensity^{\vdim}(\|V\|,x) \ge d \quad \text{for $\|V\|$ almost all $x \in \oball{a}{R}$}
        \,, \qquad \|V\| \oball{a}{R} \le (1 + \eta) d \unitmeasure{\vdim} R^{\vdim} \,, \\
        \delta V(g) \le \eta d^{1/p} R^{\vdim/p - 1} \Big( \textint{}{} |g|^q \ud \|V\| \Big)^{1/q}
        \quad \text{for $g \in \VF(\R^{\adim})$ with $\spt g \subseteq \oball{a}{R}$} \,.
    \end{gather*}
    Then there are $T \in \pgrass{\adim}{\vdim}$ and a~continuously differentiable map $F : \im T
    \to \R^{\adim}$ such that $T \circ F = \id{\im T}$,
    \begin{gather*}
        \|\uD F(y) - \uD F(z)\| \le \epsilon (|y - z|/R)^{1 - \vdim/p} \quad \text{for $y, z \in \im
        T$} \,, \\
        \oball{a}{(1 - \epsilon) R} \cap \spt \|V\| = \oball{a}{(1 - \epsilon) R} \cap \im F \,.
    \end{gather*}
\end{theorem}

\begin{proof}
    If the assertion holds for some~$\epsilon$, it holds with the same~$\eta$ for every larger
    $\epsilon < 1$; hence we assume $\epsilon \le 1/2$. Let $\epsilon_0 = \inf \{ 1/2,
    \unitmeasure{\vdim}/2 \}$, let $\eta_0$ and~$C_0$ correspond to $\epsilon_0$ as $\eta$ and~$C$
    in~\ref{al:8.19}, and let
    \begin{displaymath}
        c_5 = 16 (8 \cdot 3^{\vdim} \unitmeasure{\vdim})^{1/2} \,, \qquad c_6 = 1 + 36 \sup \im
        |\Psi_1''| + 18 \sup \im |\Psi_1'| \,,
    \end{displaymath}
    with $\Psi_1$ as in~\ref{mr:harmonic}. Choose $\epsilon_1 > 0$ with
    \begin{displaymath}
        \epsilon_1 \le \inf \{ \epsilon/32, \eta_0, \epsilon^3 / (64 c_6) \} \,, \qquad c_5
        \epsilon_1 / \epsilon \le \eta_0 \,, \qquad C_0 c_5 \epsilon_1 / \epsilon \le \epsilon^2
        / (64 c_6) \,,
    \end{displaymath}
    and let $\eta_1$ correspond to $\epsilon_1$ as $\eta$ in~\ref{al:8.8}. Choose $\eta > 0$ with
    $\eta \le \inf \{ 1, \eta_0, \eta_1 \}$ and $C_0 \eta \le \epsilon^2 / (64 c_6)$. Let
    $\lambda_0 = \sup \{ c_5 \epsilon_1 / \epsilon, \eta \}$ and $\sigma_0 = C_0 \lambda_0$; then
    $\lambda_0 \le \eta_0$ and $\sigma_0 \le \epsilon^2 / (64 c_6) \le \epsilon^2 / 64$.

    Suppose $R$, $d$, $V$, $a$ satisfy the hypotheses. Replacing $V$ by $d^{-1} (\scale{1/R} \circ
    \trans{-a})_\# V$ and using~\ref{al:3.2(2)} and~\ref{al:8.2}, we may assume $a = 0$, $R = 1$
    and $d = 1$; the conclusion for the original~$V$ follows with $T$ unchanged and $F$ replaced by
    $[y \mapsto a + R F(R^{-1}(y - T(a)))]$. Then
    \begin{displaymath}
        \delta V(g) \le \eta \Big( \textint{}{} |g|^q \ud \|V\| \Big)^{1/q} \quad \text{for $g \in
        \VF(\R^{\adim})$ with $\spt g \subseteq \oball{0}{1}$} \,.
    \end{displaymath}
    By~\ref{al:8.6}, $\density^{\vdim}(\|V\|,x)$ exists for $x \in \oball{0}{1}$ and is at least~$1$
    for $x \in \spt \|V\| \cap \oball{0}{1}$. By~\ref{al:8.8} with $\epsilon$ replaced
    by~$\epsilon_1$, and since $\eta \le \eta_1$,
    \begin{equation}
        \label{eq:8.RT:B}
        (1 - \epsilon_1) \unitmeasure{\vdim} \le s^{-\vdim} \|V\| \cball{x}{s} \le (1 + \epsilon_1)
        \unitmeasure{\vdim} \,, \qquad 1 \le \density^{\vdim}(\|V\|,x) \le 1 + \epsilon_1
    \end{equation}
    whenever $x \in \spt \|V\| \cap \oball{0}{1 - \epsilon_1}$ and $0 < s < (1 - \epsilon_1)(1 -
    |x|)$; and, applying~\ref{al:8.8} at $b = 0$ with $r = r_0 = 1 - 2 \epsilon_1$, there is $T \in
    \pgrass{\adim}{\vdim}$ with $\im T = X(0,r_0)$ and
    \begin{equation}
        \label{eq:8.RT:A}
        \dist(x, \im T) \le \epsilon_1 \quad \text{for $x \in \spt \|V\| \cap \oball{0}{r_0}$}
        \,,
    \end{equation}
    since $X(0,r_0)$ is an affine plane containing~$0$, hence a~linear subspace. Let $W = V
    \restrict (\oball{0}{r_0} \times \pgrass{\adim}{\vdim})$. Then $\spt \|W\|$ is a~compact subset
    of $\Clos(\spt \|V\| \cap \oball{0}{r_0})$; since $\dist(\cdot, \im T)$ is continuous,
    \eqref{eq:8.RT:A} holds for $x \in \spt \|W\|$.

    Let $r = \epsilon/16$ and suppose $b \in \spt \|V\| \cap \oball{0}{1 - \epsilon/2}$; then $b \in
    \spt \|W\|$. If $x \in \spt \|W\| \cap \cylind{T}{b}{2r}$, then $|T(x)| < |b| + \epsilon/8$ and
    $|x| \le |T(x)| + \epsilon_1 < 1 - \epsilon/4$; moreover $|T^\perp(x - b)| \le 2 \epsilon_1 \le
    r$, hence $|x - b| < 3r$. By~\ref{new:cutoff} choose $\chi \in \dspace{\R^{\adim}}{\R}$ with
    $0 \le \chi \le 1$, $\chi = 1$ on a~neighbourhood of $\cball{0}{1 - \epsilon/4}$ and $\spt \chi \subseteq \oball{0}{1
    - \epsilon/8} \subseteq \oball{0}{r_0}$. If $g \in \VF(\R^{\adim})$ and $\spt g \subseteq
    \cylind{T}{b}{2r}$, then $\spt \|V\| \cap \oball{0}{r_0} \cap \spt g \subseteq \spt \|W\| \cap
    \cylind{T}{b}{2r}$, hence $\uD g(x) = \uD(\chi g)(x)$ for $x \in \spt \|V\| \cap
    \oball{0}{r_0}$, and $\chi g = 0$ outside $\oball{0}{r_0}$; consequently $\delta W(g) = \delta
    V(\chi g) \le \eta (\textint{}{} |\chi g|^q \ud \|V\|)^{1/q} \le \eta (\textint{}{} |g|^q \ud
    \|W\|)^{1/q}$. Let $V_b = \trans{-b \#} W$.
    We verify the hypotheses of~\ref{al:8.19} for $V_b$, $r$, $T$ and $\alpha = \eta$, with
    $\epsilon$ replaced by~$\epsilon_0$. By~\ref{al:8.2}, \ref{al:8.19(3)} holds since $r^{1 -
    \vdim/p} \eta \le \eta \le \eta_0$. Using~\eqref{eq:8.RT:B} with $s = 3r < (1 - \epsilon_1)(1 -
    |b|)$, we obtain
    \begin{displaymath}
        \mu_b = \Big( r^{-\vdim-2} \textint{\cylind{T}{b}{2r}}{} \dist(x, \trans{b}[\im T])^2 \ud
        \|W\| x \Big)^{1/2} \le \big( r^{-\vdim-2} (2 \epsilon_1)^2 \cdot 2 \unitmeasure{\vdim} (3
        r)^{\vdim} \big)^{1/2} = c_5 \epsilon_1 / \epsilon \le \eta_0 \,,
    \end{displaymath}
    which gives~\ref{al:8.19(2)}. If $x \in \spt \|W\| \cap \cylind{T}{b}{2r}$ and $0 < s < r$, then
    $\cball{x}{s} \subseteq \oball{0}{r_0}$, so that $\|W\| \cball{x}{s} = \|V\| \cball{x}{s}$, and
    $s < (1 - \epsilon_1)(1 - |x|)$,
    so~\eqref{eq:8.RT:B} yields~\ref{al:8.19(4)} and~\ref{al:8.19(5)}, because $(1 - \epsilon_1)
    \unitmeasure{\vdim} \ge \epsilon_0$, $1 + \epsilon_1 \le 2 - \epsilon_0$ and $\epsilon_1 \le
    \eta_0$. By~\ref{al:8.19} there are $f_b : \im T \cap \oball{0}{r} \to \im T^\perp$ and $F_b$
    with $F_b(y) = y + f_b(y)$,
    \begin{gather}
        \label{eq:8.RT:C}
        \trans{b}[\spt \|V_b\|] \cap \cylind{T}{b}{r} = \spt \|W\| \cap \cylind{T}{b}{r} = \trans{b}[\im
        F_b] \,, \qquad F_b(0) = 0 \,, \\
        \label{eq:8.RT:D}
        \|\uD f_b(y) - \uD f_b(z)\| \le \sigma_0 (|y - z|/r)^{1 - \vdim/p} \,, \qquad \|\uD f_b(y)\|
        \le \sigma_0 \quad \text{for $y, z \in \im T \cap \oball{0}{r}$} \,,
    \end{gather}
    where $F_b(0) = 0$ because $0 \in \spt \|V_b\|$ and $T(0) = 0$.

    If $x, x' \in \spt \|W\| \cap \oball{0}{1 - \epsilon/2}$ and $T(x) = T(x')$, then $x' - x \in
    \im F_x$ and $T(x' - x) = 0$, so $x' = x + F_x(0) = x$ by~\eqref{eq:8.RT:C}. Let $\Omega = \im T
    \cap \oball{0}{1 - 3\epsilon/4}$ and $\Sigma = \spt \|W\| \cap \cylind{T}{0}{1 - 3 \epsilon/4}$.
    If $x \in \Sigma$, then $|x| \le |T(x)| + \epsilon_1 < 1 - \epsilon/2$; thus $T | \Sigma$ is
    univalent and $T[\Sigma] \subseteq \Omega$. By~\eqref{eq:8.RT:C}, $T[\Sigma]$ contains $\Omega \cap
    \oball{T(x)}{r}$ for $x \in \Sigma$, so $T[\Sigma]$ is open; it is relatively closed in~$\Omega$
    because $\spt \|W\|$ is compact; and $0 \in T[\Sigma]$. Since $\Omega$ is connected, $T[\Sigma] =
    \Omega$. Let $F_\Omega = (T | \Sigma)^{-1}$ and $f = T^\perp \circ F_\Omega$. If $y \in \Omega$, $b
    = F_\Omega(y)$ and $z \in \Omega \cap \oball{y}{r}$, then $b + F_b(z - y) \in \spt \|W\|$ and
    $T(b + F_b(z - y)) = z$, hence $F_\Omega(z) = b + F_b(z - y)$ and $f(z) = f(y) + f_b(z - y)$. By
    \eqref{eq:8.RT:D}, $f$ is continuously differentiable, $\|\uD f\| \le \sigma_0$, $|f| \le
    \epsilon_1$ by~\eqref{eq:8.RT:A}, and $\|\uD f(y) - \uD f(z)\| \le \sigma_0 (16/\epsilon) |y -
    z|^{1 - \vdim/p} \le (\epsilon/4) |y - z|^{1-\vdim/p}$ for $y, z \in \Omega$ with $|y - z| <
    r$. If $y, z \in \Omega$ and $|y - z| \ge r$, then $|y - z|^{1 - \vdim/p} \ge \epsilon/16$ and
    $\|\uD f(y) - \uD f(z)\| \le 2 \sigma_0 \le \epsilon^2/32$. Hence
    \begin{equation}
        \label{eq:8.RT:E}
        \|\uD f(y) - \uD f(z)\| \le (\epsilon/2) |y - z|^{1 - \vdim/p} \quad \text{for $y, z \in
        \Omega$} \,.
    \end{equation}
    Moreover $\spt \|V\| \cap \oball{0}{1 - \epsilon} \subseteq \Sigma$.

    Let $\zeta(y) = \Psi_1(16 (3 \epsilon)^{-1} (|y| - 1 + \epsilon))$ for $y \in \im T$. Then
    $\zeta$ is of class~$\infty$, $\zeta(y) = 1$ for $|y| \le 1 - \epsilon$, $\spt \zeta \subseteq
    \cball{0}{1 - 13 \epsilon/16}$, and, since $|y| \ge 1/2$ on $\spt \uD \zeta$, $\|\uD \zeta\| \le
    c_6/\epsilon$ and $\|\uD^2 \zeta\| \le c_6/\epsilon^2$. Let $\tilde f(y) = \zeta(y) f(y)$ for $y
    \in \Omega$, $\tilde f(y) = 0$ for $y \in \im T \without \Omega$, and $F(y) = y + \tilde f(y)$.
    Then $F$ is continuously differentiable, $T \circ F = \id{\im T}$, and $\uD \tilde f = \zeta \uD
    f + \uD \zeta \cdot f$ on~$\Omega$. For $y, z \in \Omega$, using $|y - z| \le 2 |y - z|^{1 -
    \vdim/p}$ and~\eqref{eq:8.RT:E},
    \begin{multline*}
        \|\uD \tilde f(y) - \uD \tilde f(z)\| \le \|\uD f(y) - \uD f(z)\| + |\zeta(y) - \zeta(z)| \,
        \|\uD f(z)\| \\ + \|\uD \zeta(y) - \uD \zeta(z)\| \, |f(y)| + \|\uD \zeta(z)\| \, |f(y) -
        f(z)| \\
        \le \big( \epsilon/2 + 2 (2 c_6 \sigma_0 / \epsilon + c_6 \epsilon_1 / \epsilon^2) \big) |y
        - z|^{1 - \vdim/p} \le \epsilon |y - z|^{1 - \vdim/p} \,.
    \end{multline*}
    If $y \in \spt \zeta$ and $z \in \im T \without \Omega$, then $|y - z| \ge \epsilon/16$ and
    $\|\uD \tilde f(y) - \uD \tilde f(z)\| = \|\uD \tilde f(y)\| \le \sigma_0 + c_6 \epsilon_1 /
    \epsilon \le \epsilon^2/32 \le \epsilon |y - z|^{1 - \vdim/p}$; in the remaining cases both
    derivatives vanish. Since $\uD F(y) - \uD F(z) = \uD \tilde f(y) - \uD \tilde f(z)$, the
    H\"older estimate holds. Finally, if $x \in \spt \|V\| \cap \oball{0}{1 - \epsilon}$, then $x \in
    \Sigma$, $|T(x)| < 1 - \epsilon$ and $x = F_\Omega(T(x)) = F(T(x))$; conversely, if $x = F(y)
    \in \oball{0}{1 - \epsilon}$, then $|y| \le |x| < 1 - \epsilon$, hence $x = F_\Omega(y) \in
    \Sigma \subseteq \spt \|V\|$.
\end{proof}

\begin{remark}
    \label{rem:8.RT}
    Allard assumes $\density^{\vdim}(\|V\|,x) \ge d$ almost everywhere; the upper density suffices
    by~\ref{al:8.6}. His proof consists of the sentence ``it will follow from 8.8 and 8.19''; the
    argument above supplies the deduction, including the univalence of the projection and the
    extension of~$F$ to $\im T$, which Allard does not discuss. Conclusion~\ref{al:8.19(10)} is used
    to control $\uD f$ in the extension.
\end{remark}

\begin{remark}[{\cite[8.1(1)]{Allard1972}}]
    \label{al:8.1(1)}
    Suppose $V \in \Var{\vdim}(\R^{\adim})$, $G$ is an open subset of~$\R^{\adim}$ and $0 \le \alpha <
    \infty$. Then
    \begin{equation}
        \label{eq:8.1(1):a}
        \delta V(g) \le \alpha \Big( \textint{}{} |g|^q \ud \|V\| \Big)^{1/q} \quad \text{for $g \in
        \VF(\R^{\adim})$ with $\spt g \subseteq G$}
    \end{equation}
    if and only if $\|\delta V\| \restrict G$ is a~Radon measure, $\|\delta V\|_{\mathrm{sing}}(G) =
    0$ and $(\textint{G}{} |h(V;x)|^p \ud \|V\| x)^{1/p} \le \alpha$, where $\|\delta
    V\|_{\mathrm{sing}}$ and $h(V;\cdot)$ are defined by~\ref{al:4.3} applied to the restriction
    of~$V$ to~$G$. If~\eqref{eq:8.1(1):a} holds, $0 < d < \infty$, and $\udensity^{\vdim}(\|V\|,x)
    \ge d$ for $\|V\|$ almost all $x \in G$, then, for every $\eta > 0$, the points $b$ of $G \cap
    \spt \|V\|$ for which there are $0 < R < \infty$ and $0 < d' < \infty$ such that the hypotheses
    of~\ref{al:8.RT} hold with $a$, $R$, $d$ replaced by $b$, $R$, $d'$ form a~dense subset of $G
    \cap \spt \|V\|$.
\end{remark}

\begin{proof}
    Suppose~\eqref{eq:8.1(1):a}. Testing with $|g| \le 1$, $\spt g \subseteq O$, gives $\|\delta
    V\|(O) \le \alpha \|V\|(O)^{1/q}$ for open $O \subseteq G$, and outer regularity extends this
    to all subsets of~$G$; thus $\|\delta V\| \restrict G$ is a~Radon measure and vanishes on
    $\|V\|$ null subsets of~$G$. If $Z = G \cap \{ x : \|\delta V\| / \|V\|(x) = \infty \}$ and $K
    \subseteq G$ is compact, then~\ref{al:4.3}\eqref{eq:4.3:3} gives $\textint{Z \cap K}{} \infty \ud
    \|V\| \le \|\delta V\|(K) < \infty$, so $\|V\|(Z) = 0$, hence $\|\delta V\|_{\mathrm{sing}}(G) =
    \|\delta V\|(Z) = 0$. By~\ref{al:4.3}, $\delta V(g) = - \textint{}{} g \bullet h(V;\cdot) \ud
    \|V\|$ for $g \in \VF(G)$, and $\textint{K}{} |h(V;\cdot)| \ud \|V\| = \|\delta V\|(K) <
    \infty$ for compact $K \subseteq G$.

    We show that $|\textint{}{} u \bullet h(V;\cdot) \ud \|V\|| \le \alpha (\textint{}{} |u|^q \ud
    \|V\|)^{1/q}$ whenever $u$ is a~bounded Borel map of $G$ into~$\R^{\adim}$ vanishing outside
    a~compact subset $K$ of~$G$. Let $c = 1 + \sup \im |u|$, and let $K'$ be a~compact
    neighbourhood of~$K$ in~$G$. Rounding the coordinates of~$u$ to integral multiples of~$2^{-j}$
    yields maps $u_j$ with finite image, vanishing outside~$K$, with $|u_j - u| \le \adim^{1/2}
    2^{-j}$. For a~Borel set $A \subseteq K$ and $\delta > 0$, \cite[2.2.5]{Federer1969} gives
    a~compact set $C \subseteq A$ and an open set $O$ with $A \subseteq O \subseteq \Int K'$ and
    $\|V\|(O \without C) < \delta$, and~\ref{new:cutoff} gives $\phi \in \dspace{\R^{\adim}}{\R}$
    with $0 \le \phi \le 1$, $\phi = 1$ on~$C$ and $\spt \phi \subseteq O$, so that $\textint{}{}
    |\phi - \CF{A}| \ud \|V\| < \delta$. Hence there are $g_j \in \VF(\R^{\adim})$ with $\spt g_j
    \subseteq K'$ and $\textint{}{} |g_j - u| \ud \|V\| \to 0$; passing to a~subsequence, $g_j \to
    u$ at $\|V\|$ almost all points. Let $\varrho(w) = \Psi_1(c^{-2} |w|^2 - 1) w$ for $w \in
    \R^{\adim}$, with $\Psi_1$ as in~\ref{mr:harmonic}; then $\varrho$ is of class~$\infty$,
    $\varrho(w) = w$ for $|w| \le c$ and $|\varrho| \le 2^{1/2} c$. The maps $\varrho \circ g_j$
    belong to $\VF(\R^{\adim})$, vanish outside~$K'$, are bounded by $2^{1/2} c$, and converge to
    $\varrho \circ u = u$ almost everywhere. Applying~\eqref{eq:8.1(1):a} to $\pm \varrho \circ g_j$
    and Lebesgue's bounded convergence theorem~\cite[2.4.9]{Federer1969}, with $|h(V;\cdot)|$
    summable over~$K'$, we obtain the claim.

    Let $K$ be a~compact subset of~$G$, $0 < M < \infty$, $E = K \cap \{ x : |h(V;x)| \le M \}$, and
    $u = |h(V;\cdot)|^{p-2} h(V;\cdot)$ on~$E$, $u = 0$ elsewhere. Then $\textint{E}{} |h(V;\cdot)|^p
    \ud \|V\| = \textint{}{} u \bullet h(V;\cdot) \ud \|V\| \le \alpha (\textint{E}{}
    |h(V;\cdot)|^p \ud \|V\|)^{1/q}$, hence $(\textint{E}{} |h(V;\cdot)|^p \ud \|V\|)^{1/p} \le
    \alpha$, and the increasing convergence theorem~\cite[2.4.7]{Federer1969} gives the bound
    on $G$. Conversely, under the three conditions, $\delta V(g) = - \textint{}{} g \bullet
    h(V;\cdot) \ud \|V\|$ for $g \in \VF(G)$ by~\ref{al:4.3}, and H\"older's
    inequality~\cite[2.4.14]{Federer1969} yields~\eqref{eq:8.1(1):a}.

    For the last assertion, let $O$ be an open ball with $\Clos O \subseteq G$ and $O \cap \spt \|V\|
    \ne \varnothing$, and let $\eta > 0$. By~\ref{al:8.6}, $\density^{\vdim}(\|V\|,\cdot)$ is real
    valued on~$G$ and at least~$d$ on $G \cap \spt \|V\|$. Let $d'$ be the supremum of the numbers
    $t$ with $\density^{\vdim}(\|V\|,x) \ge t$ for $\|V\|$ almost all $x \in O$; then $d \le d' <
    \infty$, because $\|V\|(O) > 0$. The set $O \cap \{ x : \density^{\vdim}(\|V\|,x) < (1 + \eta/2)
    d' \}$ has positive $\|V\|$ measure, hence contains a~point $b \in \spt \|V\|$. Then
    $\density^{\vdim}(\|V\|,x) \ge d'$ for $\|V\|$ almost all $x \in O$, $\lim_{R \to 0+} R^{-\vdim}
    \|V\| \oball{b}{R} = \unitmeasure{\vdim} \density^{\vdim}(\|V\|,b) < (1 + \eta) d'
    \unitmeasure{\vdim}$, and $\alpha \le \eta d'^{1/p} R^{\vdim/p - 1}$ for small~$R$, since
    $\vdim < p$. Hence the hypotheses of~\ref{al:8.RT} hold with $a$, $d$ replaced by $b$, $d'$ and
    every sufficiently small~$R$ with $\oball{b}{R} \subseteq O$.
\end{proof}

\begin{example}[{\cite[8.1(2)]{Allard1972}}]
    \label{al:8.1(2)}
    Let $F$ be the subset of $\{ t : 0 \le t \le 1 \}$ obtained by removing, for $m = 1, 2,
    \ldots$, the open middle intervals of length $4^{-m}$ from each of the $2^{m-1}$ closed intervals
    remaining after $m - 1$ steps. Then $F$ is closed, $\LM^1(F) = 1 - \textsum{m=1}{\infty} 2^{m-1}
    4^{-m} = 1/2$, and $F$ contains no interval. The set $\R \without F$ is the union of
    disjointed open intervals $I_1, I_2, \ldots$; let $b_j(t) = \chi(t - s_j) \chi(t_j - t)$ if $I_j
    = \{ t : s_j < t < t_j \}$, $b_j(t) = \chi(t_j - t)$ if $I_j = \{ t : t < t_j \}$, and $b_j(t) =
    \chi(t - s_j)$ if $I_j = \{ t : s_j < t \}$, with $\chi$ as in~\ref{mr:harmonic}. Each $b_j$ is
    of class~$\infty$ with bounded derivatives, $b_j > 0$ on~$I_j$ and $b_j = 0$ off~$I_j$. Choosing
    $\epsilon_j > 0$ with $\epsilon_j \sup \im |b_j^{(m)}| \le 2^{-j}$ for $m \le j$, the series $f
    = \textsum{j=1}{\infty} \epsilon_j b_j$ converges with all derivatives uniformly on~$\R$, so $f$
    is of class~$\infty$, $f \ge 0$ and $F = \R \cap \{ t : f(t) = 0 \}$. Let $E_1 = \R^2 \cap \{ x :
    x_2 = 0 \}$, $E_2 = \R^2 \cap \{ x : x_2 = f(x_1) \}$, and $V = \var{}(E_1) + \var{}(E_2) \in
    \IVar{1}(\R^2)$. Then $h(V;\cdot)$ is locally bounded, so that~\ref{al:8.1(1)}\eqref{eq:8.1(1):a}
    holds for every bounded open~$G$ with a~suitable~$\alpha$;
    \begin{gather*}
        \density^1(\|V\|,x) = 0 \quad \text{for $x \in \R^2 \without \spt \|V\|$} \,, \qquad
        \density^1(\|V\|,x) = 2 \quad \text{for $x \in F \times \{0\}$} \,, \\
        \density^1(\|V\|,x) = 1 \quad \text{for $x \in \spt \|V\| \without (F \times \{0\})$} \,;
    \end{gather*}
    and for no $a \in F \times \{0\}$ is $\spt \|V\|$ near~$a$ the image of a~map~$F'$ as
    in~\ref{al:8.RT}, although $\HM^1(F \times \{0\}) = 1/2$.
\end{example}

\begin{proof}
    Let $\gamma(t) = (t, f(t))$ and $\tau = |\gamma'|^{-1} \gamma'$. For $g \in \VF(\R^2)$, by
    \ref{al:3.5:lem}, the area formula~\cite[3.2.3]{Federer1969}, \ref{al:4.2:lem}
    and~\ref{al:2.3(1)}\eqref{eq:2.3:4}, $\delta \var{}(E_2)(g) = \textint{}{} (g \circ \gamma)'
    \bullet \tau \ud \LM^1 = - \textint{}{} g(\gamma(t)) \bullet \tau'(t) \ud \LM^1 t$, by
    integration by parts~\ref{new:ibp}; similarly $\delta \var{}(E_1) = 0$. Hence $\delta V(g) = -
    \textint{E_2}{} g \bullet \kappa \ud \HM^1$ with $\kappa(\gamma(t)) = |\gamma'(t)|^{-1}
    \tau'(t)$, a~continuous function; since $\|V\| \ge \HM^1 \restrict E_2$, we have $\|\delta V\|
    \le |\kappa| \, \|V\|$, and H\"older's inequality gives~\eqref{eq:8.1(1):a} on bounded~$G$ with
    $\alpha = \sup \{ |\kappa(x)| : x \in G \cap E_2 \} \|V\|(G)^{1/p}$. The densities follow
    from the area formula~\cite[3.2.3]{Federer1969} applied to the maps $[t \mapsto (t,0)]$
    and~$\gamma$: since $\gamma'$ is continuous, $\lim_{s \to 0+} (2s)^{-1} \HM^1(E_i \cap
    \cball{x}{s}) = 1$ for $x \in E_i$ and $i \in \{1,2\}$, while $E_1 \cap E_2 = F \times \{0\}$
    and every point of $E_i \without E_{3-i}$ has a~neighbourhood disjoint from $E_{3-i}$. Suppose $a = (s,0)$ with $s
    \in F$, $\rho > 0$, and $\oball{a}{\rho} \cap \spt \|V\| = \oball{a}{\rho} \cap \im F'$ with
    $T \circ F' = \id{\im T}$. Since every point of $F$ is a~limit of removed intervals whose
    lengths tend to zero, there is $j$ with $\Clos I_j \subseteq \{ t : |t - s| < \rho/2 \}$ and
    $\sup \{ f(t) : t \in I_j \} < \rho/2$. The set $\Gamma = \{ (t,0) : t \in \Clos I_j \} \cup \{
    \gamma(t) : t \in \Clos I_j \}$ is a~simple closed curve contained in $\oball{a}{\rho} \cap
    \spt \|V\|$, and $T | \Gamma$ is univalent and continuous. This is impossible, since the image
    of~$\Gamma$ is a~compact connected subset of the line $\im T$, that is a~segment, which is
    disconnected by the removal of an interior point, whereas $\Gamma$ is not disconnected by the
    removal of a~point.
\end{proof}

\begin{remark}
    \label{rem:8.1(2)}
    Allard refers to~\cite[3.1.14]{Federer1969} for the existence of~$f$; the direct construction
    above avoids it. The example shows that the density hypothesis of~\ref{al:8.RT} cannot be
    replaced by a~bound on the first variation alone.
\end{remark}

\begin{corollary}[{\cite[8.1(3)]{Allard1972}}]
    \label{al:8.1(3)}
    Suppose $E$ is an $\HM^{\vdim}$ measurable subset of~$\R^{\adim}$ such that $E \cap K$ is
    $(\HM^{\vdim},\vdim)$ rectifiable for every compact $K \subseteq \R^{\adim}$, $0 < \alpha <
    \infty$, and $G$ is an open subset of~$\R^{\adim}$ with compact closure such that
    \begin{displaymath}
        \frac{\ud}{\ud t} \HM^{\vdim}(G \cap h_t[E]) \Big|_{t=0} \le \alpha \Big( \textint{G \cap
        E}{} |\dot h_0|^q \ud \HM^{\vdim} \Big)^{1/q}
    \end{displaymath}
    whenever $h$, $h_t$ and $\dot h_0$ are as in~\ref{al:4.1} with $M = \R^{\adim}$, $\epsilon =
    \infty$, and each $h_t$ is univalent. Then there is a~relatively closed subset $K$ of~$G$ with
    $\HM^{\vdim}(K) = 0$ such that to each $a \in G \cap \spt(\HM^{\vdim} \restrict E) \without K$
    correspond $T \in \pgrass{\adim}{\vdim}$, an open subset $W$ of $\im T$, positive numbers $R$
    and~$C$, and a~continuously differentiable map $F : W \to \R^{\adim}$ with $T \circ F =
    \id{W}$,
    \begin{displaymath}
        \oball{a}{R} \cap \spt(\HM^{\vdim} \restrict E) = \im F \,, \qquad \|\uD F(y) - \uD F(z)\|
        \le C |y - z|^{1 - \vdim/p} \quad \text{for $y, z \in W$} \,.
    \end{displaymath}
\end{corollary}

\begin{proof}
    Let $V = \var{}(E \cap G) \in \Var{\vdim}(\R^{\adim})$; by~\ref{al:3.5:lem}, $\|V\| = \HM^{\vdim}
    \restrict (E \cap G)$, and $\spt \|V\| \cap G = \spt(\HM^{\vdim} \restrict E) \cap G$. Suppose
    $g \in \VF(\R^{\adim})$ with $\spt g \subseteq G$, and let $h_t(x) = x + t g(x)$. For small
    $|t|$, $h_t$ is univalent, $h_t[G] = G$ and $h_t[E] \cap G = h_t[E \cap G]$, because $h_t$ is the
    identity outside $\spt g$. Hence~\ref{al:4.1}\eqref{eq:4.1:2}, \ref{al:3.5:lem}
    and~\ref{al:4.2:lem} give
    \begin{displaymath}
        \delta V(g) = \frac{\ud}{\ud t} \HM^{\vdim}(G \cap h_t[E]) \Big|_{t=0} \le \alpha \Big(
        \textint{}{} |g|^q \ud \|V\| \Big)^{1/q} \,,
    \end{displaymath}
    after the hypothesis is applied to $h$ restricted to a~small interval of times, which does not
    affect the derivative at~$0$. By~\ref{al:2.8(4)(a)}(a), $\density^{\vdim}(\|V\|,x) =
    \CF{E \cap G}(x)$ for $\HM^{\vdim}$ almost all~$x$, in particular $\udensity^{\vdim}(\|V\|,x) \ge
    1$ for $\|V\|$ almost all $x \in G$. By~\ref{al:8.6}, $\density^{\vdim}(\|V\|,\cdot)$ is upper
    semicontinuous on~$G$ and at least~$1$ on $G \cap \spt \|V\|$. Let $\eta$ correspond to
    $\epsilon = 1/2$ as in~\ref{al:8.RT}, and $K = G \cap \spt \|V\| \cap \{ x :
    \density^{\vdim}(\|V\|,x) \ge 1 + \eta \}$. Then $K$ is relatively closed in~$G$, and
    $\HM^{\vdim}(K) = 0$ by the preceding density formula. If $a \in G \cap \spt \|V\| \without K$,
    then, for small $R > 0$ with $\oball{a}{R} \subseteq G$, $\|V\| \oball{a}{R} \le (1 + \eta)
    \unitmeasure{\vdim} R^{\vdim}$ and $\alpha \le \eta R^{\vdim/p - 1}$; hence~\ref{al:8.RT} with
    $d = 1$ yields $T$ and $F_0 : \im T \to \R^{\adim}$ with $\oball{a}{R/2} \cap \spt \|V\| =
    \oball{a}{R/2} \cap \im F_0$ and $\|\uD F_0(y) - \uD F_0(z)\| \le \frac12 R^{\vdim/p - 1} |y -
    z|^{1 - \vdim/p}$. The open set $W = F_0^{-1}[\oball{a}{R/2}]$, the map $F = F_0 | W$, $R/2$ in
    place of~$R$, and $C = \frac12 R^{\vdim/p - 1}$ have the required properties.
\end{proof}

\begin{remark}
    \label{rem:8.1(3)}
    Allard's proof reads ``One needs only to recall 4.1 and 2.8(4)(a)''; the argument above uses in
    addition~\ref{al:8.6} for the closedness of~$K$. The hypothesis is applied to the deformations
    $h_t(x) = x + t g(x)$, which satisfy the conditions of~\ref{al:4.1} for small~$|t|$.
\end{remark}

\begin{remark}[{\cite[8.1(4)]{Allard1972}}]
    \label{al:8.1(4)}
    Allard poses the following question. Suppose $V \in \Var{\vdim}(\R^{\adim})$, $G$ is an open
    subset of~$\R^{\adim}$, $\|\delta V\|(G) = 0$, and $\density^{\vdim}(\|V\|,x) \ge 1$ for $\|V\|$
    almost all $x \in G$. Does the conclusion of~\ref{al:8.RT} hold at $\|V\|$ almost all points
    of~$G$?
\end{remark}

\begin{corollary}[{\cite[8.1(5)]{Allard1972}}]
    \label{al:8.1(5)}
    Suppose $\vdim \ge 2$ and $i : \sphere{\adim-1} \to \R^{\adim}$ is the inclusion map. There is
    a~number $\mu_0 > 1$ with the following property. Suppose $B \in
    \Var{\vdim-1}(\sphere{\adim-1})$, $B \ne 0$, $B$ is stationary, $\density^{\vdim-1}(\|i_\#
    B\|,x) \ge 1$ for $\|B\|$ almost all $x \in \sphere{\adim-1}$, and $\|B\|(\sphere{\adim-1}) <
    \mu_0 \vdim \unitmeasure{\vdim}$. Then there are $\lambda$ with $1 \le \lambda < \mu_0$ and $T \in
    \pgrass{\adim}{\vdim}$ with $B = \lambda \var{}(\im T \cap \sphere{\adim-1})$.
\end{corollary}

\begin{proof}
    Let $\eta$ correspond to $\epsilon = 1/2$ as in~\ref{al:8.RT}, and $\mu_0 = 1 + \eta$. Let $C
    \in \Var{\vdim}(\oball{0}{1})$ be defined as in~\ref{al:5.2(3)} with $R = 1$ and $B(r) = B$ for
    $0 < r < 1$; then $C$ is stationary, and
    \begin{displaymath}
        C(\alpha') = \textint{0}{1} r^{\vdim-1} \textint{}{} \alpha'(r x, T' \oplus \R x) \ud (i_\#
        B)(x,T') \ud \LM^1 r
    \end{displaymath}
    for every nonnegative Borel function $\alpha'$ on $\Gk{\oball{0}{1}}$, by~\ref{rem:2.6(2)}. Let
    $\lambda = \|B\|(\sphere{\adim-1}) / (\vdim \unitmeasure{\vdim})$; then $\lambda < \mu_0$,
    $\|C\| \cball{0}{s} = s^{\vdim} \|B\|(\sphere{\adim-1}) / \vdim$ for $0 < s < 1$,
    $\density^{\vdim}(\|C\|,0) = \lambda$, $0 \in \spt \|C\|$ because $B \ne 0$, and $\|C\|(\{ t x : x \in A \}) = t^{\vdim} \|C\|(A)$ whenever $A \subseteq
    \oball{0}{1}$ and $\{ t x : x \in A \} \subseteq \oball{0}{1}$. For $\LM^1$ almost all~$r$, the
    varifold $C'(r)$ of~\ref{al:5.2(2)}(c) is given by the inner integral above with the factor
    $r^{\vdim-1}$, and, since $T'$ is the projection onto a~subspace of $\Tan(\sphere{\adim-1},x)$
    for $i_\# B$ almost all $(x,T')$, the varifold $B(r)$ of~\ref{al:5.2(2)}(d) equals~$B$.
    Since $\|C\| \oball{0}{1} = \density^{\vdim}(\|C\|,0) \unitmeasure{\vdim}$ and $0 \in \spt
    \|C\|$, \ref{al:5.2(2)}(f) applies and gives $\density^{\vdim}(\|C\|, r x) =
    \density^{\vdim-1}(\|i_\# B\|, x) \ge 1$ for $\|C\|$ almost all $r x$. Let $j$ be the inclusion
    of $\oball{0}{1}$ into~$\R^{\adim}$. Since $\delta (j_\# C)(g) = 0$ whenever $\spt g \subseteq
    \oball{0}{1}$, \ref{al:8.6} applied to $j_\# C$ with $G = \oball{0}{1}$, $\alpha = 0$ and $d =
    1$ gives $\lambda = \density^{\vdim}(\|C\|,0) \ge 1$.

    Apply~\ref{al:8.RT} to $j_\# C$ with $a = 0$, $R = 1$, $d = 1$: the hypotheses hold since $\delta (j_\# C)(g) = 0$ for $\spt g \subseteq \oball{0}{1}$ and
    $\|C\| \oball{0}{1} = \lambda \unitmeasure{\vdim} < (1 + \eta) \unitmeasure{\vdim}$. We obtain $T
    \in \pgrass{\adim}{\vdim}$ and $F$ of class~$1$ with $T \circ F = \id{\im T}$ and
    $\oball{0}{1/2} \cap \spt \|C\| = \oball{0}{1/2} \cap \im F$; thus $F(0) = 0$. Let $W =
    F^{-1}[\oball{0}{1/2}]$. If $y \in W$ and $0 < t \le 1$, then $t F(y) \in \spt \|C\| \cap
    \oball{0}{1/2}$, by the homogeneity of~$\|C\|$, and $T(t F(y)) = t y$; hence $F(t y) = t F(y)$.
    Differentiating at $t = 0$ gives $F(y) = \uD F(0)(y)$ for $y \in W$. Thus $\oball{0}{1/2} \cap
    \spt \|C\| \subseteq \im P$, where $P$ is the orthogonal projection onto $\im \uD F(0)$, and $\im
    P \cap \oball{0}{\delta} \subseteq \spt \|C\|$ for some $\delta > 0$, since $W$ is
    a~neighbourhood of~$0$. By homogeneity, $\spt \|C\| = \im P \cap \oball{0}{1}$. Since $C$ is
    stationary, \ref{al:4.6(3)} with $M = \im P \cap \oball{0}{1}$ gives $C = c \var{}(M)$ with $c
    = \density^{\vdim}(\|C\|,0) = \lambda$.

    Consequently $C(\alpha') = 0$ for $\alpha'(y,S) = \varphi(|y|) \|S - P\|$ with $\varphi \in
    \ccspace{\{ r : 0 < r < 1 \}}$ positive, hence $T' \oplus \R x = P$ for $i_\# B$ almost all
    $(x,T')$, and then $x \in \im P$ and $T'$ is the projection onto $\im P \cap \{ x \}^\perp =
    \Tan(\im P \cap \sphere{\adim-1}, x)$. Moreover, taking $\alpha'(y,S) = \varphi(|y|)
    \CF{A}(|y|^{-1} y)$ for Borel $A \subseteq \sphere{\adim-1}$, and computing $\lambda
    \HM^{\vdim}(\im P \cap \{ r x : x \in A, 0 < r < 1 \})$ by~\cite[3.2.13]{Federer1969} in $\im P$,
    we obtain $\|B\| = \lambda \HM^{\vdim-1} \restrict (\im P \cap \sphere{\adim-1})$. Hence $B =
    \lambda \var{}(\im P \cap \sphere{\adim-1})$, and the assertion holds with $T$ replaced by~$P$.
\end{proof}

\begin{remark}
    \label{rem:8.1(5)}
    Allard's statement does not exclude $\vdim = 1$ and $B = 0$; the conclusion fails for $B = 0$,
    and for $\vdim = 1$ the construction of~\ref{al:5.2(3)} requires the hypothesis $\textint{}{} x
    \ud \|i_\# B\| x = 0$ instead of stationarity. With these provisos, the proof follows Allard's
    sketch: the cone over~$B$ satisfies the hypotheses of~\ref{al:8.RT}, its support near~$0$ is
    a~graph which is a~cone, hence a~plane, and~\ref{al:4.6(3)} identifies the cone.
\end{remark}
